\documentclass[12pt,oneside,openany]{book}
\usepackage[T1]{fontenc}
\usepackage[utf8]{inputenc}
\usepackage{amsmath,amssymb,amsthm,mathtools,mathrsfs}
\DeclareFontFamily{U}{rsfs}{\skewchar\font127 }
\DeclareFontShape{U}{rsfs}{m}{n}{<-6.5>rsfs5<6.5-7.5>rsfs7<7.5->rsfs10}{}
\usepackage{booktabs,array}
\usepackage[numbers,sort&compress]{natbib}
\usepackage{environ}
\PassOptionsToPackage{hypertexnames=false}{hyperref}

\makeatletter

\usepackage{lmodern}
\usepackage[sc]{mathpazo}
\usepackage{bm}
\linespread{1.1}
\usepackage{microtype}
\microtypesetup{stretch=15,shrink=15}

\usepackage{geometry}
\geometry{letterpaper,left=1in,right=1in,top=0.75in,bottom=0.75in,includeheadfoot,
  headheight=15pt,headsep=16pt,footskip=26pt}
\setlength{\parindent}{1.5em}
\setlength{\parskip}{0pt}
\setlength{\emergencystretch}{3em}
\tolerance=1000
\clubpenalty=3000 \widowpenalty=3000 \displaywidowpenalty=1500
\raggedbottom
\allowdisplaybreaks[2]

\usepackage{enumitem}
\setlist{topsep=4pt plus 2pt minus 1pt,itemsep=2pt plus 1pt,parsep=0pt}
\setlist[itemize]{leftmargin=1.5em}
\setlist[enumerate]{leftmargin=2em}

\usepackage[newparttoc]{titlesec}
\usepackage{titletoc,fancyhdr}

\newcommand{\manuscript@spacedsc}[1]{{\scshape\textls[90]{\MakeLowercase{#1}}}}
\newcommand{\manuscript@headpar}{\rightskip=0pt plus 0.5\hsize\relax
  \parfillskip=0pt plus 0.5\hsize\relax\hyphenpenalty=10000\relax\exhyphenpenalty=10000\relax}

\titleformat{\part}[display]
  {\normalfont\centering}
  {\large\manuscript@spacedsc{\partname\ \thepart}}
  {22pt}
  {\fontsize{22}{28}\selectfont}
\assignpagestyle{\part}{empty}

\titleformat{\chapter}[display]
  {\normalfont\manuscript@headpar}
  {\manuscript@spacedsc{\chaptertitlename\ \thechapter}}
  {12pt}
  {\fontsize{20}{25}\selectfont}
\titleformat{name=\chapter,numberless}[display]
  {\normalfont\manuscript@headpar}
  {}
  {0pt}
  {\fontsize{20}{25}\selectfont\manuscript@starmark}
\titlespacing*{\chapter}{0pt}{12pt}{30pt}
\newcommand{\manuscript@starmark}[1]{#1\markboth{#1}{#1}}

\titleformat{\section}
  {\normalfont\fontsize{13.5}{17}\selectfont\bfseries\manuscript@headpar}{\thesection}{0.75em}{}
\titlespacing*{\section}{0pt}{22pt plus 4pt minus 2pt}{9pt plus 2pt}
\titleformat{\subsection}
  {\normalfont\normalsize\bfseries\manuscript@headpar}{\thesubsection}{0.75em}{}
\titlespacing*{\subsection}{0pt}{15pt plus 3pt minus 2pt}{6pt plus 1pt}
\titleformat{\paragraph}[runin]
  {\normalfont\normalsize\bfseries}{}{0pt}{}
\titlespacing*{\paragraph}{0pt}{9pt plus 2pt minus 1pt}{0.6em}

\setcounter{tocdepth}{1}
\setcounter{secnumdepth}{1}

\newcommand{\shorthead}[2]{\expandafter\def\csname manuscript@sh@\detokenize{#1}\endcsname{#2}}
\newcommand{\manuscript@lookup}[1]{%
  \ifcsname manuscript@sh@\detokenize{#1}\endcsname
    \expandafter\let\expandafter\manuscript@headtext\csname manuscript@sh@\detokenize{#1}\endcsname
  \else
    \def\manuscript@headtext{#1}%
  \fi}
\newcommand{\manuscript@headfont}{\normalfont\small\itshape}
\DeclareRobustCommand{\manuscript@hd}[2]{%
  \ifx\relax#1\relax\else{\upshape#1}\hspace{0.8em}\fi#2}
\newsavebox{\manuscript@chapbox}
\newsavebox{\manuscript@secbox}
\newsavebox{\manuscript@cachedchapbox}
\newsavebox{\manuscript@cachedsecbox}
\newif\ifmanuscript@headtoolong
\newif\ifmanuscript@cachedheadtoolong
\let\manuscript@cachedheadkey\relax
\newcommand{\manuscript@computeheads}{%
  \manuscript@headtoolongfalse
  \sbox\manuscript@chapbox{\manuscript@headfont\nouppercase{\leftmark}}%
  \sbox\manuscript@secbox{}%
  \protected@edef\manuscript@l{\leftmark}%
  \protected@edef\manuscript@r{\rightmark}%
  \ifx\manuscript@r\@empty\else\ifx\manuscript@r\manuscript@l\else
    \sbox\manuscript@secbox{\manuscript@headfont\nouppercase{\rightmark}}%
  \fi\fi
  \dimen@=\wd\manuscript@chapbox\advance\dimen@ 4em\relax
  \ifdim\dimen@>\headwidth
    \sbox\manuscript@chapbox{\manuscript@headfont\@chapapp\ \thechapter}%
    \manuscript@headtoolongtrue
  \fi
  \dimen@=\wd\manuscript@chapbox\advance\dimen@\wd\manuscript@secbox\advance\dimen@ 8em\relax
  \ifdim\dimen@>\headwidth
    \sbox\manuscript@secbox{\manuscript@headfont
      \expandafter\let\csname manuscript@hd \endcsname\manuscript@hdnum\nouppercase{\rightmark}}%
  \fi}
\newcommand{\manuscript@heads}[1]{%
  \protected@edef\manuscript@headkey{%
    {\leftmark}{\rightmark}{\the\headwidth}%
    {\fontname\font}{\the\dimexpr1em\relax}{\thechapter}{\@chapapp}}%
  \ifx\manuscript@headkey\manuscript@cachedheadkey\else
    \manuscript@computeheads
    \global\setbox\manuscript@cachedchapbox=\copy\manuscript@chapbox
    \global\setbox\manuscript@cachedsecbox=\copy\manuscript@secbox
    \global\let\manuscript@cachedheadkey\manuscript@headkey
    \ifmanuscript@headtoolong
      \global\manuscript@cachedheadtoolongtrue
    \else
      \global\manuscript@cachedheadtoolongfalse
    \fi
  \fi
  \setbox\manuscript@chapbox=\copy\manuscript@cachedchapbox
  \setbox\manuscript@secbox=\copy\manuscript@cachedsecbox
  \ifmanuscript@cachedheadtoolong
    \ifx\relax#1\relax\else\PackageWarningNoLine{manuscript-style}{Running head too long on
      page \thepage; give the chapter a \string\shorthead}\fi
  \fi}
\newcommand{\manuscript@hdnum}[2]{\ifx\relax#1\relax\else\S\,{\upshape#1}\fi}

\pagestyle{fancy}
\fancyhf{}
\fancyhead[L]{\manuscript@heads{warn}\usebox\manuscript@chapbox}
\fancyhead[R]{\manuscript@heads{}%
  \ifdim\wd\manuscript@secbox>\z@\usebox\manuscript@secbox\hspace{2em}\fi
  \normalfont\small\thepage}
\renewcommand{\headrulewidth}{0pt}
\renewcommand{\footrulewidth}{0pt}
\fancypagestyle{plain}{\fancyhf{}\fancyfoot[C]{\normalfont\small\thepage}%
  \renewcommand{\headrulewidth}{0pt}\renewcommand{\footrulewidth}{0pt}}
\renewcommand{\chaptermark}[1]{\manuscript@lookup{#1}%
  \markboth{\manuscript@hd{\ifnum\c@secnumdepth>\m@ne\if@mainmatter\thechapter\fi\fi}{\manuscript@headtext}}{}}
\renewcommand{\sectionmark}[1]{\manuscript@lookup{#1}%
  \markright{\manuscript@hd{\ifnum\c@secnumdepth>\z@\thesection\fi}{\manuscript@headtext}}}

\contentsmargin{2.2em}
\titlecontents{part}[0pt]
  {\addvspace{16pt plus 2pt}\normalfont}
  {\manuscript@spacedsc{\partname\ \thecontentslabel}\hspace{1em}}
  {}
  {\hfill\contentspage}
  [\addvspace{4pt}]
\titlecontents{chapter}[2.2em]
  {\addvspace{7pt plus 1pt}\normalfont}
  {\contentslabel{2.2em}}
  {\hspace*{-2.2em}}
  {\hfill\contentspage}
\titlecontents{section}[4.8em]
  {\normalfont\small}
  {\contentslabel{2.6em}}
  {\hspace*{-2.6em}}
  {\titlerule*[0.9em]{.}\contentspage}

\@ifpackageloaded{natbib}{%
  \renewcommand{\bibsection}{\chapter*{References}%
    \addcontentsline{toc}{chapter}{References}}%
}{}

\newlength{\manuscriptenvskip}
\setlength{\manuscriptenvskip}{10pt plus 3pt minus 2pt}
\newtheoremstyle{plain}{\manuscriptenvskip}{\manuscriptenvskip}
  {\itshape}{}{\bfseries}{.}{0.5em}{}
\newtheoremstyle{definition}{\manuscriptenvskip}{\manuscriptenvskip}
  {\normalfont}{}{\bfseries}{.}{0.5em}{}
\newtheoremstyle{remark}{\manuscriptenvskip}{\manuscriptenvskip}
  {\normalfont}{}{\itshape}{.}{0.5em}{}

\usepackage[font=small,labelsep=period,skip=8pt]{caption}

\usepackage{etoolbox}
\patchcmd{\@footnotetext}{\footnotesize}{\small}{}{}

\usepackage[hidelinks,bookmarksnumbered=true]{hyperref}
\usepackage{bookmark}
\hypersetup{pdfauthor={},pdfpagemode=UseNone,pdfdisplaydoctitle=true,
  pdfstartview={FitH \hypercalcbp{\paperheight}}}

\makeatother

\usepackage{xurl}%
\usepackage[nameinlink,capitalize,noabbrev]{cleveref}
\usepackage{aliascnt}
\hypersetup{pdftitle={Existence of the q-Colorability Threshold for Sparse Random Graphs},pdfauthor={P. M. Aronow and Patrick Lopatto}}
\renewcommand{\chaptername}{Section}
\numberwithin{equation}{chapter}

\theoremstyle{plain}
\newtheorem{theorem}{Theorem}[chapter]
\newaliascnt{proposition}{theorem}
\newtheorem{proposition}[proposition]{Proposition}
\aliascntresetthe{proposition}
\newaliascnt{lemma}{theorem}
\newtheorem{lemma}[lemma]{Lemma}
\aliascntresetthe{lemma}
\newaliascnt{corollary}{theorem}
\newtheorem{corollary}[corollary]{Corollary}
\aliascntresetthe{corollary}
\theoremstyle{definition}
\newaliascnt{definition}{theorem}
\newtheorem{definition}[definition]{Definition}
\aliascntresetthe{definition}
\newaliascnt{example}{theorem}
\newtheorem{example}[example]{Example}
\aliascntresetthe{example}
\theoremstyle{remark}
\newaliascnt{remark}{theorem}
\newtheorem{remark}[remark]{Remark}
\aliascntresetthe{remark}
\makeatletter
\@addtoreset{proposition}{chapter}\@addtoreset{lemma}{chapter}\@addtoreset{corollary}{chapter}
\@addtoreset{definition}{chapter}\@addtoreset{example}{chapter}\@addtoreset{remark}{chapter}
\makeatother
\crefname{theorem}{Theorem}{Theorems}
\crefname{proposition}{Proposition}{Propositions}
\crefname{lemma}{Lemma}{Lemmas}
\crefname{corollary}{Corollary}{Corollaries}
\crefname{definition}{Definition}{Definitions}
\crefname{example}{Example}{Examples}
\crefname{remark}{Remark}{Remarks}
\crefname{chapter}{Section}{Sections}
\crefname{appendix}{Appendix}{Appendices}
\crefname{part}{Part}{Parts}
\crefname{section}{Section}{Sections}
\crefname{subsection}{Section}{Sections}
\crefname{table}{Table}{Tables}
\crefname{equation}{}{}
\crefformat{equation}{(#2#1#3)}
\Crefformat{equation}{(#2#1#3)}
\crefrangeformat{equation}{(#3#1#4)--(#5#2#6)}
\Crefrangeformat{equation}{(#3#1#4)--(#5#2#6)}
\crefmultiformat{equation}{(#2#1#3)}{ and~(#2#1#3)}{, (#2#1#3)}{, and~(#2#1#3)}
\Crefmultiformat{equation}{(#2#1#3)}{ and~(#2#1#3)}{, (#2#1#3)}{, and~(#2#1#3)}
\newcommand{\creflastconjunction}{, and~}
\newcommand{\creflastgroupconjunction}{, and~}

\makeatletter
\newcommand{\manuscript@noteblock}[2]{%
  \ifx\relax#1\relax\else{\centering\manuscript@spacedsc{#1}\par}\vspace{6pt}\fi
  {\centering
   \begin{minipage}{0.84\textwidth}\small\normalfont\setlength{\parindent}{1.5em}%
     \noindent\ignorespaces#2\end{minipage}\par}}
\NewEnviron{authorsnote}{%
  \clearpage\thispagestyle{empty}%
  \vspace*{\stretch{2}}%
  \manuscript@noteblock{Authors' note}{\BODY}%
  \vspace*{\stretch{3}}%
  \clearpage}
\NewEnviron{abstractpage}{%
  \clearpage\thispagestyle{empty}%
  \vspace*{\stretch{2}}%
  \manuscript@noteblock{Abstract}{\BODY}%
  \vspace*{\stretch{3}}%
  \clearpage}
\newcommand{\manuscripttitlepage}[2]{%
  \hypersetup{pageanchor=false}%
  \begin{titlepage}
  \centering\normalfont
  \vspace*{1.7in}
  {\large\strut\par}
  \vspace{1.05in}
  {\fontsize{28}{35}\selectfont #1\par}
  \vspace{1.1in}
  {\Large\def\\{\par\vspace{10pt}}#2\par}
  \vfill
  \end{titlepage}%
  \hypersetup{pageanchor=true}}
\makeatother

\usepackage{tikz}
\usetikzlibrary{arrows.meta,positioning,shapes.geometric,calc,fit,backgrounds,decorations.pathreplacing}
\newcommand{\E}{\mathbb E}
\newcommand{\Pp}{\mathbb P}
\newcommand{\R}{\mathbb R}
\newcommand{\dd}{\,\mathrm d}
\DeclareMathOperator{\Var}{Var}
\newcommand{\eps}{\varepsilon}
\newcommand{\ind}{\mathbf 1}
\newcommand{\cE}{\mathcal E}
\newcommand{\cX}{\mathcal X}
\newcommand{\abs}[1]{\left|#1\right|}
\DeclareMathOperator{\Col}{Col}
\DeclareMathOperator{\Bin}{Bin}
\DeclareMathOperator{\Pois}{Pois}
\newcommand{\cF}{\mathcal F}

\allowdisplaybreaks[2]
\setcounter{tocdepth}{0}

\clubpenalty=10000
\patchcmd{\thebibliography}{\clubpenalty4000}{\clubpenalty10000}{}{}
\begin{document}
\frontmatter
\manuscripttitlepage{Existence of the\\ $q$-Colorability Threshold\\ for Sparse Random Graphs}{P.\ M.\ Aronow\\ Patrick Lopatto}
\begin{authorsnote}
With the exception of these notes, the entire document is machine-written. We
have carefully supervised its preparation to ensure that the proofs are readable and
that previous literature is appropriately cited. We are releasing this manuscript
on Hexagon to disseminate the results as quickly as possible because they may be
of interest to the mathematical community. A more carefully written version of
this manuscript is in preparation and will be submitted for publication. We
welcome suggestions concerning mathematical corrections, omitted citations, and
attribution of ideas.

P.\ L.\ was partially supported by NSF grant DMS-2450004.
\end{authorsnote}
\begin{abstractpage}
For a fixed integer $q \ge 3$, we call a graph $q$-colorable if its vertices admit $q$ colors with adjacent vertices colored differently. We prove that there exists $d_q > 0$ such that the random graph on $n$ vertices with independent edge probability $d/n$ is $q$-colorable with probability tending to one for fixed $0 \le d < d_q$ and to zero for fixed $d > d_q$, as $n \to \infty$. The same holds for the uniform random graph with $\lfloor cn \rfloor$ edges, for fixed $c \ge 0$, with threshold $d_q/2$ in $c$. For fixed $\epsilon \in (0, 1)$, the first loss of colorability under independent uniform additions of pairs of distinct vertices lies with probability at least $1 - \epsilon$ in a deterministic window of $O_{q, \epsilon}(n(\log n)^{-3/2})$ steps, where $\log$ is natural. The proof uses interpolation to establish convergence and bounds the expected number of pairs needed to destroy colorability of the last colorable graph by identifying random vertices and controlling the growth of classes equal in every coloring.

\end{abstractpage}
\tableofcontents
\mainmatter

\chapter{Introduction}\label{sec:intro}

\section{Background}\label{sec:intro-background}

Let $n \ge 1$ and $q \ge 2$ be integers. We let $\log$ denote the natural logarithm, and let $\Pp$, $\E$, and $\Var$ denote probability, expectation, and variance, respectively. For every integer $k \ge 1$, we set $[k] = \{1, \dots, k\}$. We call a map $\sigma\colon[n] \to [q]$ a \emph{$q$-coloring} of a graph on $[n]$ if $\sigma(u) \ne \sigma(v)$ for all edges $\{u, v\}$. We call a graph \emph{$q$-colorable} if it has a $q$-coloring. We define its \emph{chromatic number} to be the least number of colors in such a coloring. For every integer $0 \le m \le \binom n2$, we let $G(n, m)$ denote a uniformly random simple graph on $[n]$ with $m$ edges. For $p \in [0, 1]$, we let $G(n, p)$ denote the random simple graph on $[n]$ in which each pair is an edge independently with probability $p$. This paper concerns $G(n, \lfloor cn\rfloor)$ and $G(n, d/n)$ for fixed constants $c, d \ge 0$ and sufficiently large $n$. The expected average degrees of these models tend to $2c$ and $d$, respectively.

The study of the chromatic number of sparse random graphs goes back to Erdős and Rényi, who asked in \cite[Section~10, p.~59]{erdos1960evolution} for the chromatic number of a random graph on $n$ vertices with about $cn$ edges, for $c > 1/2$. For every $q \ge 3$, Achlioptas and Friedgut proved that $q$-colorability has a sharp threshold \cite[Theorem~1.1]{achlioptas1999sharp}. More precisely, there exists a sequence $(\hat d_n)$ (which depends on $q$) such that for every fixed $\omega > 0$,
\begin{equation}\label{eq:intro-af}
\begin{gathered}
\lim_{n \to \infty} \Pp\bigl(G(n, \min\{1, \max\{0, (\hat d_n - \omega)/n\}\}) \text{ is $q$-colorable}\bigr) = 1, \\
\lim_{n \to \infty} \Pp\bigl(G(n, \min\{1, (\hat d_n + \omega)/n\}) \text{ is $q$-colorable}\bigr) = 0.
\end{gathered}
\end{equation}
Their proof rests on Friedgut's criterion for sharp thresholds \cite{friedgut1999sharp}, which relates a coarse threshold to the approximation of the property by the appearance of bounded subgraphs, and it leaves the dependence of $\hat d_n$ on $n$ undetermined. Achlioptas and Friedgut expressed the belief that $\hat d_n$ converges, and they recorded its convergence as an open problem \cite[Remark~1.2]{achlioptas1999sharp}. Achlioptas and Molloy had already noted that this convergence was widely conjectured \cite[p.~204]{achlioptas1997analysis}, and later works describe it as a well-known open problem \cite[Section~1.1]{cojaoghlan2016chromatica}, \cite[Section~1]{bapst2016condensation}. The goal of this paper is to prove that $\hat d_n$ converges for every $q \ge 3$.

We first recall the known bounds on the location of the transition. For $q = 3$, the random graph $G(n, d/n)$ is $3$-colorable with probability tending to one for $d \le 4.03$ \cite[Corollary~2]{achlioptas2003almost}, and with probability tending to zero for $d > 4.697$ \cite[Section~1.3]{ayre2022lower} (see \cite{chvatal1991almost,achlioptas1997analysis,kaporis2000note,dubois2002noncolourability} for earlier bounds). For general $q$, we set $\gamma_q = \log q / \log(q/(q - 1))$. We use $o_q(1)$ for a quantity tending to zero as $q \to \infty$, after the limit in $n$ at fixed $q$. A first-moment computation in $G(n, \lfloor cn \rfloor)$ gives non-colorability with probability tending to one when $c > \gamma_q$. Comparing the fixed-edge and binomial models gives the same conclusion for $G(n, d/n)$ when $d > 2\gamma_q$. We have $2\gamma_q = 2q \log q - \log q + o_q(1)$. Achlioptas and Naor showed that $G(n, d/n)$ is $q$-colorable with probability tending to one for $d < 2(q - 1) \log(q - 1)$, and that its chromatic number is concentrated on two explicit consecutive integers \cite[Theorem~1]{achlioptas2005two}. Coja-Oghlan and Vilenchik \cite[Theorem~1.1]{cojaoghlan2016chromatica} and Coja-Oghlan \cite[Theorem~1]{cojaoghlan2013upper} then proved that
\begin{equation}\label{eq:intro-largeq}
\liminf_{n \to \infty} \hat d_n \ge 2q \log q - \log q - 2 \log 2 - o_q(1), \qquad \limsup_{n \to \infty} \hat d_n \le 2q \log q - \log q - 1 + o_q(1),
\end{equation}
and Coja-Oghlan and Vilenchik also showed that the chromatic number of $G(n, d/n)$ is concentrated on a single value for all $d$ in an explicit set of asymptotic density one \cite[Theorem~1.2]{cojaoghlan2016chromatica}. These results locate the transition. However, they do not show that $\hat d_n$ converges.

Statistical physics gives precise, non-rigorous predictions. Statistical-physics computations place the $3$-colorability threshold near average degree $4.69$ \cite{mulet2002coloring}, refined numerically to $4.687(2)$ \cite[Table~II]{zdeborova2007phase}, and give the asymptotic form $2q \log q - \log q - 1 + o_q(1)$ for large $q$ \cite{krzakala2004threshold}, which agrees with the upper bound in \eqref{eq:intro-largeq}. The same theory predicts further transitions in the set of colorings below the colorability threshold \cite{krzakala2007gibbs,zdeborova2007phase}. Rigorous results on the number of colorings below that threshold were obtained for large $q$ \cite[Theorem~2.1]{bapst2016condensation} and later for all $q \ge 3$ \cite[arXiv:1611.00814v4, Theorem~1.2]{cojaoghlan2018information}.

For other random constraint satisfaction problems, the existence of a limiting threshold is known in several cases. The threshold of random $2$-SAT is at one clause per variable \cite{chvatal1992mick,goerdt1996threshold}, with a window of order $n^{2/3}$ clauses \cite[Theorem~1.1]{bollobas2001scaling}. The threshold of random $k$-XORSAT was determined for $k = 3$ in \cite{dubois2002xorsat} and for all $k \ge 3$ in \cite{pittel2016satisfiability}. Ding, Sly, and Sun proved the existence of the threshold for random $k$-SAT, and identified its value, for all sufficiently large $k$ \cite[Theorem~1]{ding2022proof}. For random $d$-regular graphs and all sufficiently large $q$, Coja-Oghlan, Efthymiou, and Hetterich located the $q$-colorability threshold in an interval of length less than one \cite[Theorem~1.1 and Corollary~1.2]{cojaoghlan2016chromaticb}. Since the degree is an integer, this leaves at most one degree at which the limiting behavior is not determined. This argument does not apply to $G(n, d/n)$, in which $d$ varies continuously.

Recently, two preprints claimed the existence of a limiting threshold for random $k$-SAT for every fixed $k \ge 3$, without identifying its value. These preprints use independent uniformly signed clauses, with distinct variables within each clause and with replacement between clauses. Carenini \cite[Theorem~1.1 and Corollary~1.2]{carenini2026polynomial} claims an upper bound of order $n^{1/2 + 1/k}$ clauses for the window, improving her earlier $O(n/\log n)$ bound \cite[Corollary~7.10]{carenini2026scaling}. She obtains the polynomial bound from a variance bound for the hitting time of unsatisfiability and deduces the existence of the threshold from a criterion of Abbe and Montanari \cite{abbe2014concentration}, in the form stated in \cite[Proposition~14]{perkins2025searching}. The preprint \cite[Theorem~1.1]{openai2026limiting} claims the same conclusion, combining a variance bound for a capped hitting time with an interpolation inequality and comparisons between systems of different sizes. The latter preprint explicitly credits Carenini with priority for the threshold-existence result \cite[abstract and Section~1]{openai2026limiting}. These preprints are unrefereed. We give self-contained proofs of the adaptations used here, while drawing on their arguments as described below. To the best of our knowledge, the convergence of $\hat d_n$ had not previously been established for a single value of $q \ge 3$.

\section{Main results}\label{sec:intro-results}

From now on, we fix an integer $q \ge 3$. Recall that
\begin{equation}\label{eq:intro-gamma}
\gamma_q = \frac{\log q}{\log(q/(q - 1))}.
\end{equation}
Our main result is the existence of the $q$-colorability threshold.

\begin{theorem}\label{thm:intro-main}
Let $q \ge 3$ be an integer. There exists $\alpha_q \in [1/2, \gamma_q]$ such that for every fixed real $c \ge 0$,
\[
\lim_{n \to \infty} \Pp\bigl(G(n, \lfloor cn \rfloor) \text{ is $q$-colorable}\bigr) =
\begin{cases}
1 & \text{if } 0 \le c < \alpha_q, \\
0 & \text{if } c > \alpha_q,
\end{cases}
\]
and for every fixed real $d \ge 0$,
\[
\lim_{n \to \infty} \Pp\bigl(G(n, d/n) \text{ is $q$-colorable}\bigr) =
\begin{cases}
1 & \text{if } 0 \le d < 2\alpha_q, \\
0 & \text{if } d > 2\alpha_q,
\end{cases}
\]
where $G(n, d/n)$ is defined for $n \ge d$.
\end{theorem}

We define $d_q = 2\alpha_q$ to be the threshold in terms of the asymptotic average degree. To relate our result to the sharp threshold sequence, we show that $\hat d_n \to d_q$. Fix $0 < \omega < d_q$. If $\hat d_n \ge d_q + \omega$ along a subsequence, the lower side of \eqref{eq:intro-af}, with shift $\omega/2$, and monotonicity imply that $G(n, (d_q + \omega/2)/n)$ is colorable with probability tending to one along that subsequence, contrary to \cref{thm:intro-main}. If $\hat d_n \le d_q - \omega$ along a subsequence, the upper side, with the same shift, implies that $G(n, (d_q - \omega/2)/n)$ is colorable with probability tending to zero, again a contradiction. This proves the claim.

The bounds on the limit inferior and limit superior in \eqref{eq:intro-largeq} now transfer directly to $d_q$. A bound at a fixed density transfers as well. If it failed for $d_q$, a density strictly between the two values would contradict monotonicity and \cref{thm:intro-main}. In particular, $4.03 \le d_3 \le 4.697$, and $d_q$ satisfies the two bounds in \eqref{eq:intro-largeq}.

The bounds $1/2 \le \alpha_q \le \gamma_q$ in \cref{thm:intro-main} are classical and are included for completeness. The upper bound is the first-moment bound mentioned above. The lower bound holds since, for $c < 1/2$, with probability tending to one every component of $G(n, \lfloor cn \rfloor)$ is a tree or has exactly one cycle \cite{erdos1960evolution}, \cite[Theorem~5.5, p.~112]{janson2000random}, and such graphs are $3$-colorable. \Cref{thm:intro-main} is an asymptotic statement. As we explain in \cref{rem:proof-finite}, our proof gives no further information on the value of $\alpha_q$ and no rate of convergence.

Our second result is a bound on the width of the transition in a random graph process. For a finite set $S$, we call a finite multiset of unordered pairs of distinct elements of $S$ a \emph{multigraph on $S$}, and we call its elements, counted with multiplicity, its \emph{occurrences}. We define a coloring of a multigraph on $S$ to be a map from $S$ to $[q]$ giving different colors to the endpoints of every occurrence. For $n \ge 2$, let $e_1, e_2, \dots$ be independent random pairs of distinct elements of $[n]$, each uniformly distributed on the $\binom{n}{2}$ pairs. We call the sequence $(e_i)_{i \ge 1}$ the \emph{edge process}. For each integer $m \ge 0$, we let $G_m$ denote its \emph{time-$m$ graph}, the multigraph on $[n]$ with occurrences $e_1, \dots, e_m$. Note that repeated pairs do not affect colorability. For $n \ge 2q$, we define $H$ to be the first integer $m \ge 0$ for which $G_m$ is not $q$-colorable. Thus
\[
H = \min\{m \ge 0 : G_m \text{ is not $q$-colorable}\}.
\]
By \cref{lem:pre-lifetime}, $H$ is finite almost surely.

\begin{theorem}\label{thm:intro-window}
Let $q \ge 3$ be an integer, let $\alpha_q$ be as in \cref{thm:intro-main}, and let $\epsilon \in (0, 1)$. There exist $C > 0$ and an integer $n_0 \ge 2q$, depending only on $q$ and $\epsilon$, and integers $m_n$ with $m_n/n \to \alpha_q$ as $n \to \infty$, such that
\begin{equation}\label{eq:intro-window}
\Pp\bigl(\abs{H - m_n} \le C n (\log n)^{-3/2}\bigr) \ge 1 - \epsilon \qquad \text{for all } n \ge n_0.
\end{equation}
\end{theorem}

Since $\Pp(H > m) = \Pp(G_m \text{ is $q$-colorable})$, \cref{thm:intro-window} implies that for all $n \ge n_0$ and integers $m \ge 0$,
\[
\begin{aligned}
\Pp(G_m \text{ is $q$-colorable}) &\ge 1 - \epsilon \qquad \text{if } m < m_n - C n (\log n)^{-3/2}, \\
\Pp(G_m \text{ is $q$-colorable}) &\le \epsilon \qquad \text{if } m \ge m_n + C n (\log n)^{-3/2}.
\end{aligned}
\]
Thus the probability that the process is colorable drops from $1 - \epsilon$ to $\epsilon$ within $O(n (\log n)^{-3/2})$ steps, and the window in the edge density $m/n$ has width $O((\log n)^{-3/2})$. We will compare this bound with known bounds on threshold widths in \cref{sec:intro-related}. We have not tried to optimize the exponent $3/2$. The convergence criterion in \cref{thm:red-convergence} requires the relative window widths to have a finite sum along exponentially growing sizes. The choice $(\log n)^{-3/2}$ meets this condition.

\Cref{thm:intro-main} is false for $q = 2$. We call a graph \emph{bipartite} if it admits a $2$-coloring. Pittel and Yeum showed that for every $d \in (0, 1)$,
\[
\lim_{n \to \infty} \Pp\bigl(G(n, d/n) \text{ is bipartite}\bigr) = e^{d/2} \Bigl(\frac{1 - d}{1 + d}\Bigr)^{1/4} \in (0, 1)
\]
\cite[Section~1, equation~(1.5)]{pittel2010frequently}. Hence there is no constant $d_2$ such that the probability of $2$-colorability tends to one below $d_2$ and to zero above it (see also the discussion in \cite[p.~64]{achlioptas1999sharp}). Our proof uses $q \ge 3$ in two places, which we describe in \cref{rem:proof-q3}. The essential one is the observation in \cref{lem:low-frozen} that, if a vertex receives the same color as another vertex in every coloring, then its neighbors use at least $q - 1 \ge 2$ colors in every coloring of the graph with this vertex removed.

Our proof has two parts, which we describe in more detail in \cref{sec:overview}. We begin by showing that the existence of $\alpha_q$ follows from the summable window bound just described, at every sufficiently small fixed probability level. To do so, we compare the probabilities of colorability at different numbers of vertices through an interpolation inequality of Bayati, Gamarnik, and Tetali \cite[Theorem~5 and Proposition~2]{bayati2013combinatorial}. We then prove the window bound of \cref{thm:intro-window}. We define the \emph{expected lifetime} of a colorable graph to be the expected number of independent uniformly random pair additions needed to make it non-colorable, including the addition that destroys colorability. A variance argument reduces the window bound to controlling this lifetime for $G_{H - 1}$.

We call $e_H$ the \emph{final edge}. For a colorable graph, we call two distinct vertices a \emph{frozen pair} if they receive the same color in every coloring, and we call each equivalence class under equality in every coloring a \emph{frozen class}. The final edge joins a frozen pair of $G_{H - 1}$. We count degrees with multiplicity. We treat separately the final edges with an endpoint of low degree (at most $\lfloor(\log n)^{1/4}\rfloor$), by identifying random vertices, and the final edges between vertices of high degree, by deleting occurrences and controlling the growth of the logarithm of the size of a frozen class.

\section{Related work}\label{sec:intro-related}

We now discuss the sources of the methods that we use, and the known bounds on the width of the transition.

The comparison of different system sizes in \cref{app:conv} uses the interpolation method. System-splitting interpolation was used by Guerra and Toninelli to prove the existence of the limiting free energy of the Sherrington--Kirkpatrick model, by interpolating between a large system and two independent subsystems \cite[Section~3]{guerra2002thermodynamic}. Franz and Leone \cite{franz2003replicaa} and Panchenko and Talagrand \cite{panchenko2004bounds} developed free-energy bounds for diluted models with Poisson numbers of interactions, including random $k$-SAT, under their respective interaction and arity hypotheses. Bayati, Gamarnik, and Tetali gave a combinatorial version for sparse random graphs \cite{bayati2013combinatorial}. They proved the convergence of the normalized optimal values of several problems and of the normalized logarithm of the probability of $q$-colorability \cite[Theorem~1]{bayati2013combinatorial}. However, such a limit does not by itself locate a transition from probability one to probability zero, since it is consistent with probabilities that tend to zero more slowly than exponentially. For coloring, they also compared the probability that a graph is colorable with the probability that two smaller disjoint graphs are colorable \cite[Theorem~5 and Proposition~2]{bayati2013combinatorial}, and their proof of this comparison uses the classes of vertices that receive the same color in every coloring. Our product inequality, \cref{lem:conv-interp}, is its form for a Poisson number of edges.

To pass from comparisons of system sizes to convergence, one uses a principle of near-superadditivity. De Bruijn and Erdős proved such a principle for sequences that are subadditive up to an error satisfying an integrability condition, for comparable sizes \cite[Theorem~23]{debruijn1952some}, and Abbe and Montanari gave a version with summable errors \cite[arXiv:1006.3786v1, Section~4.3, Lemma~8 and Remark~3]{abbe2014concentration}. In \cite[Section~5]{openai2026limiting}, the product inequality \cite[Lemma~5.2]{openai2026limiting} is used to compare the normalized center of the transition for the combined system with the minimum of the centers for the two parts, with polynomially small errors, and a balanced tree of blocks yields convergence \cite[Lemma~5.4]{openai2026limiting}. \Cref{app:conv} adapts this argument to quantiles of the hitting time and to errors that are only summable along $n = \lfloor e^k \rfloor$. We show in \cref{rem:conv-necessity} that a window of order $n/\log n$ would not suffice for this argument.

The reduction of the window to the expected lifetime of the last colorable graph, in \cref{sec:red-deletion}, has several antecedents. Aldous \cite[Theorem~1(b)]{aldous1991threshold} considered the cover time of a finite set by independent copies of a random subset, the first time at which their union is the whole set. He bounded the variance of the cover time by an absolute constant times its mean times the expected time that a fresh sequence of such subsets needs to cover the set left uncovered just before the cover time. Coloring fits this framework, if each pair $\{u, v\}$ covers the maps $\sigma \colon [n] \to [q]$ with $\sigma(u) = \sigma(v)$. In this case, $H$ is the cover time, and the set left uncovered at time $H - 1$ is the set of colorings of $G_{H - 1}$. Carenini applied Aldous's bound to the last satisfiable formula in random $k$-SAT \cite[Theorem~4.1]{carenini2026polynomial}. In \cite[Lemmas~3.1 and~3.2]{openai2026limiting}, the variance of a capped hitting time is bounded through the Efron--Stein inequality \cite{efron1981jackknife,steele1986efron}, in the form \cite[(3.6)]{boucheron2005moment} in which one variable is omitted, by deleting one clause at a time without renumbering the other clauses, and a later preprint combines expected lifetimes with such deletions \cite[Section~3.3, Lemma~3.4]{openai2026variance}. \Cref{lem:red-deletion} follows the deletion argument of \cite{openai2026limiting} and evaluates the sum of the squared deletion delays exactly, in terms of the expected lifetime of the last colorable graph.

The estimates specific to $k$-SAT in \cite{carenini2026polynomial,openai2026limiting} bound the sum over time of the fraction of variables that take the same value in every satisfying assignment \cite[Proposition~3.2]{carenini2026polynomial}, \cite[Proposition~2.2]{openai2026limiting}. These proofs delete a variable and compare the clauses that contained it, which become clauses of length $k - 1$, with clauses of length $k$. This step does not carry over to colorings. Since permuting the colors maps colorings to colorings, no vertex of a colorable graph receives a fixed named color in every coloring. What can be forced is the equality of two colors. Recall that the final edge $e_H$ joins two vertices $u$ and $v$ such that $\sigma(u) = \sigma(v)$ for every coloring $\sigma$ of $G_{H - 1}$.

Culberson and Gent studied frozen pairs and their development as edges are added to a graph, and used frozen-pair counts to compute conditional one-edge survival probabilities in a process without replacement \cite[Sections~2.1 and~3.3]{culberson2001frozen}. The hazard identity in \cref{lem:pre-hazard} is the version for independent pairs sampled with replacement. For a satisfiable Boolean formula, we call the literals true in every satisfying assignment its \emph{backbone} \cite{monasson1996entropy,monasson1999determining}. Following \cite{bollobas2001scaling}, we call the literals whose addition makes some satisfiable subformula unsatisfiable its \emph{spine}. These notions record an absolute Boolean value, whereas a frozen pair records equality without specifying a named color. In \cite[Section~3.3]{bollobas2001scaling}, the probability that the next clause makes a $2$-SAT formula unsatisfiable is expressed through the spine, and \cref{lem:pre-hazard} is the analogous identity for frozen pairs. Other notions of freezing concern a variable that has one value throughout a specified restricted family of solutions \cite{achlioptas2009random,molloy2018freezing,semerjian2008freezing}. Frozen pairs here quantify over all colorings. Our estimates for the frozen pairs of the last colorable graph, in \cref{sec:low,sec:high}, replace the estimates specific to $k$-SAT. We now describe the ingredients of these estimates.

To treat final edges with an endpoint of low degree, we identify vertices. Let $Z$ be a nonempty set of vertices of a graph. Its colorings that are constant on $Z$ correspond to the colorings of the graph in which $Z$ is identified to a single vertex. For two nonadjacent vertices, this correspondence underlies the addition--identification relation for chromatic polynomials \cite[Section~3, Theorem~1]{read1968introduction}. We combine it with the splitting of a Poisson process of edges at a vertex \cite[Theorems~5.2, 5.6, and~5.8, Corollary~5.9]{last2018lectures}, and we prove in \cref{prop:low-identify} that identifying at most $(\log n)^{1/4}$ random vertices advances the hitting time of non-colorability, in expectation, by $O(n^{3/8})$.

For final edges between vertices of high degree, we instead delete edge occurrences. The deletion identities are the analogue, for a fixed number of independent points, of the Mecke equation for Poisson processes \cite{mecke1967stationare}, \cite[Theorem~4.1]{last2018lectures}. We call an occurrence at one endpoint of a frozen pair \emph{pivotal} for that pair if deleting it removes the forced equality; otherwise we call it \emph{non-pivotal}. This is the usual notion of sensitivity to one coordinate \cite{russo1982approximate}, applied to the event that adding the given pair destroys colorability. The main estimate of \cref{sec:high} bounds the total, along the process, of the density of a relation that describes how the frozen class of a vertex grows when an edge at this vertex is added. We prove it by charging this density to the growth of the logarithm of the size of the frozen class. Bounds on hitting and occupation times through a function whose expected increments can be controlled are classical \cite[Sections~2--3]{hajek1982hitting}. In the $k$-SAT preprints, Carenini bounds the occupation of states through a concave function of their expected lifetime \cite[Lemma~3.1]{carenini2026polynomial}, and a later preprint charges the visits of the process above a level of the probability of becoming unsatisfiable to the expected lifetime at the first visit \cite[Section~4.1, Lemma~4.1]{openai2026variance}. A related bounded-potential occupation argument appears in \cite[Lemmas~2.2 and~4.1]{openai2026linear}: expected increments of a bounded monotone potential control accumulated powers of short-clause killing probabilities. The inequality that we need between the density of the relation and the expected growth of the frozen class, \cref{lem:high-onestep}, is specific to colorings.

Finally, we compare \cref{thm:intro-window} with known bounds on the width of transitions. At fixed probability levels, for every nontrivial increasing graph property that is invariant under relabeling, Friedgut and Kalai proved that the transition in $G(n, p)$ takes place within an interval of values of $p$ of length $O(1/\log n)$ \cite[Theorem~1.1, p.~2993]{friedgut1996every}, and Bourgain and Kalai improved this to $O((\log n)^{-2 + \delta})$ for every $\delta > 0$ \cite{bourgain1997influences}, \cite[Theorem~1]{bourgain1999threshold}. A recent unrefereed preprint claims the bound $O((\log n)^{-2})$ for every nontrivial monotone graph property invariant under relabeling \cite[Theorem~1.1]{openai2026graphthreshold}. These additive bounds do not give a shrinking relative window when $p$ is of order $1/n$; the derivative formulation of the Bourgain--Kalai improvement also restricts how close $p$ may be to zero or one. The small-$p$ refinement in \cite[Corollary~3.3, p.~2996]{friedgut1996every} gives an upper bound $O(1/n)$ when $p = \Theta(1/n)$, corresponding to $O(n)$ expected edges; it still does not yield a shrinking relative window. The theorem of Achlioptas and Friedgut shows that the window of $q$-colorability has width $o(n)$ in the number of edges, with no rate, and \cref{thm:intro-window} improves this to $O(n (\log n)^{-3/2})$. In the other direction, Wilson proved that this window has width at least of order $\sqrt{n}$ edges for every $q \ge 3$ \cite[Corollary~6]{wilson2002critical}.

We do not know the true order of the window. For comparison, the window of random $2$-SAT has order $n^{2/3}$ clauses \cite[Theorem~1.1]{bollobas2001scaling}. For random $k$-SAT, Carenini claims a window upper bound of order $n^{1/2 + 1/k}$ clauses \cite[Theorem~1.1]{carenini2026polynomial}, and two further preprints claim that the hitting time of unsatisfiability has variance of order $n$, for $k \ge 4$ \cite[Theorem~1.1]{openai2026variance} and for $k = 3$ \cite[Theorem~1.1]{openai2026linear}. As we explain in \cref{rem:proof-sqrt}, when the probability that $H$ lies strictly above the lower quantile and at or below the upper quantile stays bounded away from zero, the variance bound of \cref{sec:red} cannot certify a window smaller than order $\sqrt n$.

\section{Organization}\label{sec:intro-organization}

\Cref{sec:overview} gives an overview of the proof. \Cref{sec:pre} collects definitions and elementary facts on the edge process, frozen pairs, and expected lifetimes. \Cref{sec:red} states the convergence criterion and reduces the window bound to two estimates for the final edge, one for endpoints of low degree and one for endpoints of high degree. \Cref{sec:low} proves the first estimate. This is the main place where the assumption $q \ge 3$ is used. \Cref{sec:high} proves the second estimate, which holds for $q = 2$ as well. \Cref{sec:proof} combines these results and proves \cref{thm:intro-main,thm:intro-window}. \Cref{app:conv} proves the convergence criterion.

\chapter{Proof overview}\label{sec:overview}

\paragraph{Notation.}\label{sec:ov-notation}
We fix an integer $q \ge 3$ and write $\log$ for the natural logarithm. For integers $k \ge 1$, we set $[k] = \{1, \dots, k\}$. We write $\Pp$, $\E$, and $\Var$ for probability, expectation, and variance, and $\ind_A$ for the indicator of an event $A$; we also use $\ind\{ \cdot \}$ for the indicator of the condition in braces. For an event $A$ and a nonnegative or integrable expression $X$ defined on $A$, we define $\E[X; A]$ to be the integral of $X$ over $A$. When we write this as $\E[X\ind_A]$, we first extend $X$ by zero outside $A$. We write $\Pois(\mu)$ for the Poisson law with mean $\mu \ge 0$ and $\Bin(m, p)$ for the binomial law with integer $m \ge 0$ and $p \in [0, 1]$. Inside probability expressions, these symbols denote random variables with the specified laws. Constants may depend on $q$ and on fixed parameters named in a statement, but not on $n$; constants denoted by $C$ or $C'$ are local to their statements. We say \emph{for all large $n$} when a statement holds for every $n \ge n_0$, where $n_0$ depends only on those fixed parameters.

Our proof has four steps. Step 1 deduces a limiting threshold from a short interval containing the loss of colorability with high probability. Step 2 bounds the length of this interval through the expected number of further pair additions needed to destroy colorability of the last colorable graph. Step 3 bounds the advance of non-colorability under random vertex identification to treat final edges at vertices of low degree. Step 4 controls the growth of classes of vertices equal in every coloring to treat final edges at vertices of high degree. Steps 1 and 2 adapt the $k$-SAT arguments of \cite{carenini2026polynomial,openai2026limiting}; our new estimates in Steps 3 and 4 bound the expected lifetime of the last colorable graph at the final edge. In the $k$-SAT proofs, the corresponding estimate uses the variables that take the same value in every satisfying assignment and compares the clauses of length $k - 1$ left by deleting a variable with clauses of length $k$ \cite{carenini2026polynomial,openai2026limiting,openai2026variance}; this comparison has no direct analogue for colorings, since permuting the colors shows that no vertex has the same color in every coloring. The dependencies between these steps are recorded in \cref{fig:ov-dependencies}. Remarks labeled heuristic describe the motivation for the estimates.

\begin{figure}[t]
\centering
\begin{tikzpicture}[x = 1cm, y = 1cm,
  box/.style = {draw, rounded corners = 2pt, align = center, font = \footnotesize, inner sep = 3pt, text width = 3.6cm, minimum height = 1.15cm},
  wide/.style = {box, text width = 4.6cm, minimum height = 1cm},
  arr/.style = {-{Stealth[length = 2mm]}, semithick}]
\node[box] (frozen) at (-6.1, 0) {\Cref{lem:low-frozen}\\ restriction on\\ neighbor colors};
\node[box] (ident) at (-2.05, 0) {\Cref{prop:low-identify}\\ identification of\\ random vertices};
\node[box] (delete) at (2.05, 0) {\Cref{lem:high-deletion},\\ \cref{prop:high-descent}\\ deleting occurrences};
\node[box] (pot) at (6.1, 0) {\Cref{lem:high-onestep},\\ \cref{thm:high-potential}\\ growth of classes\\ of equal colors};
\node[box] (low) at (-4.1, -2.5) {\Cref{thm:low-main}\\ final edges at low degree};
\node[box] (high) at (4.1, -2.5) {\Cref{thm:high-main}\\ final edges at high degree};
\node[box] (window) at (-4.1, -4.6) {\Cref{prop:red-window}\\ window and lifetime};
\node[box] (split) at (0, -4.6) {\Cref{prop:red-split}\\ splitting by degrees};
\node[wide] (wb) at (-2.6, -6.7) {\Cref{thm:proof-window}\\ bound on the window};
\node[wide] (conv) at (3.4, -6.7) {\Cref{thm:red-convergence} (\cref{app:conv})\\ convergence of the threshold};
\node[wide] (main) at (0.4, -8.7) {\Cref{thm:intro-main,thm:intro-window}};
\draw[arr] (frozen) -- (low);
\draw[arr] (ident) -- (low);
\draw[arr] (delete) -- (high);
\draw[arr] (pot) -- (high);
\draw[arr] (low) -- (split);
\draw[arr] (high) -- (split);
\draw[arr] (window) -- (split);
\draw[arr] (split) -- (wb);
\draw[arr] (wb) -- (main);
\draw[arr] (conv) -- (main);
\end{tikzpicture}

\caption{The main steps of the proof. An arrow means that the first result is used to prove the second. The left branch is the low-degree estimate of \cref{sec:low}, and the right branch is the high-degree estimate of \cref{sec:high}. The convergence criterion is proved in \cref{app:conv}.}
\label{fig:ov-dependencies}
\end{figure}

We begin by defining the edge process. For a finite set $S$, we call a finite multiset of unordered pairs of distinct elements of $S$ a \emph{multigraph on $S$}, and we call its elements, counted with multiplicity, its \emph{occurrences}. We call a map $\sigma:S \to [q]$ a \emph{coloring} when the endpoints of every occurrence have different colors, and call the multigraph \emph{colorable} when there exists such a map. For a multigraph $G$ and a pair $e$ from its vertex set, we let $G + e$ denote the result of adding one occurrence of $e$; when $e$ occurs in $G$, we let $G - e$ denote the result of removing one occurrence. We define $\deg_G(u)$ to be the number of occurrences containing the vertex $u$. We call a uniformly chosen pair of distinct elements of the specified vertex set a \emph{uniform pair}.

Let $n \ge 2q$ be an integer. We let $e_1, e_2, \dots$ denote independent uniform pairs of $[n]$, and call this sequence the \emph{edge process}. For each integer $m \ge 0$, we define $G_m$ to be the multigraph consisting of the first $m$ occurrences of the process. We define $H$ to be the first $m$ for which $G_m$ is not colorable. We let $\cF_m$ denote the $\sigma$-algebra generated by $e_1, \dots, e_m$. The lifetime bound below implies that $H$ is finite almost surely. These definitions apply on other finite vertex sets when they are specified.

\section{Step 1: the window determines the threshold}\label{sec:ov-convergence}

We set $\gamma_q = \log q/\log(q/(q - 1))$. Throughout this overview, we take $\Lambda = \gamma_q + 2$ and set $M = \lfloor\Lambda n\rfloor$. We define $V$ to be the last colorable time capped at $M$ by setting $V = \min(H - 1, M)$. This makes $V$ a bounded function of $e_1, \dots, e_M$. For each $\eta \in (0, 1/2]$, we define $L_n(\eta)$ to be the smallest nonnegative integer $m$ with $\Pp(V \le m) \ge \eta$, and define $U_n(\eta)$ to be the largest integer $m \in \{0, \dots, M\}$ with $\Pp(V \ge m) \ge \eta$. We call $w_n(\eta) = U_n(\eta) - L_n(\eta)$ the \emph{window} at level $\eta$. These definitions give
\begin{equation}\label{eq:ov-quantiles}
\begin{aligned}
L_n(\eta) &= \min\{m \ge 0 : \Pp(V \le m) \ge \eta\}, \\
U_n(\eta) &= \max\{m \in \{0, \dots, M\} : \Pp(V \ge m) \ge \eta\}.
\end{aligned}
\end{equation}
Here the extrema are over nonnegative integers. By \cref{lem:pre-quantile}, $G_m$ is colorable with probability greater than $1 - \eta$ for $0 \le m \le L_n(\eta)$, and with probability smaller than $\eta$ for $U_n(\eta) < m \le M$. Applying \cref{lem:pre-centre}(b) with $c_0 = 1/4$ and part (a) with $c_1 = (\gamma_q + \Lambda)/2$, we have $n/4 < L_n(\eta) \le U_n(\eta) < M$ for all large $n$. The lower bound uses $q \ge 3$ to color components having at most one cycle.

We define $\eta_0$ to be the small quantile cutoff given by
\[
\eta_0 = \frac{e^{-2\Lambda}}8.
\]
We will use the following convergence criterion, proved in \cref{thm:red-convergence}. Let $(\delta_n)_{n \ge 1}$ be a nonincreasing sequence in $(0, 1]$ such that $\sum_{k \ge 1}\delta_{\lfloor e^k\rfloor} < \infty$. Suppose that, for all $\eta \in (0, \eta_0]$, there exists $K'(\eta) > 0$ such that
\begin{equation}\label{eq:ov-window-hyp}
w_n(\eta) \le K'(\eta) \, n \delta_n \qquad \text{for all large } n.
\end{equation}
Then there exists $\alpha_q \in [1/2, \gamma_q]$ such that $L_n(\eta)/n \to \alpha_q$ and $U_n(\eta)/n \to \alpha_q$ for all $\eta \in (0, \eta_0]$. For every fixed real $c \ge 0$ with $c \ne \alpha_q$,
\[
\lim_{n \to \infty}\Pp(H > \lfloor cn\rfloor) =
\begin{cases}
1 & \text{if } c < \alpha_q, \\
0 & \text{if } c > \alpha_q.
\end{cases}
\]
The same limits hold for the probability of colorability of a uniformly random simple graph on $[n]$ with $\lfloor cn\rfloor$ edges, and of a simple graph on $[n]$ in which each pair is an edge independently with probability $2c/n$. These models are defined for all sufficiently large $n$. In Steps 2--4, we prove \eqref{eq:ov-window-hyp} with $\delta_n = (\log n)^{-3/2}$ for $n \ge 3$ and $\delta_1 = \delta_2 = 1$. This sequence is admissible, since $\delta_{\lfloor e^k \rfloor} \le (k - 1)^{-3/2}$ for $k \ge 2$, which proves \cref{thm:intro-main}. For the error probability $\epsilon \in (0, 1)$ of \cref{thm:intro-window}, we set $\eta_1 = \min(\epsilon/2, \eta_0)$ and define $m_n = L_n(\eta_1)$. Using the window bound at level $\eta_1$, we get the stated concentration around $m_n$. The convergence criterion also implies that $m_n/n \to \alpha_q$, as shown in \cref{sec:proof-main}.

To prove \cref{thm:red-convergence}, we compare quantiles at different numbers of vertices through an auxiliary Poisson model. For every integer $n \ge 1$ and real $c \ge 0$, we let $A_n(c)$ denote the multiset formed by first sampling a $\Pois(cn)$ number of occurrences and, independently of that count, sampling ordered endpoints independently and uniformly in $[n] \times [n]$ for those occurrences. We call an occurrence whose two endpoints coincide a \emph{loop}. On a finite vertex set, we call a finite multiset of unordered pairs, allowing coincident endpoints, a \emph{multigraph with loops}. We require different colors at the endpoints of every occurrence; a loop prevents colorability. We let $\rho_n(c)$ denote the probability that $A_n(c)$ is colorable. \Cref{lem:conv-interp} states that
\begin{equation}\label{eq:ov-product}
\rho_{N_1 + N_2}(c) \ge \rho_{N_1}(c) \, \rho_{N_2}(c) \qquad \text{for all integers } N_1, N_2 \ge 1 \text{ and all } c \ge 0.
\end{equation}

Let $N = N_1 + N_2$, and set $S_1 = [N_1]$ and $S_2 = [N] \setminus [N_1]$. For a colorable multigraph $G$ on a finite vertex set and a vertex $u$, we define the \emph{frozen class} $C_G(u)$ to be the set of vertices that receive the same color as $u$ in every coloring, following \cite{culberson2001frozen}. For $0 \le \upsilon \le c$, we define $A^{(\upsilon)}$ to be the union of three independent families. These consist of $\Pois(N\upsilon)$ occurrences with independent uniform ordered endpoints in $[N]^2$, and, for $j = 1, 2$, $\Pois(N_j(c - \upsilon))$ occurrences with independent uniform ordered endpoints in $S_j^2$. The counts and endpoints are independent. We set $\chi(\upsilon) = \Pp(A^{(\upsilon)}\text{ is colorable})$. Adding an occurrence destroys colorability exactly when its endpoints lie in the same frozen class. On the event that $A^{(\upsilon)}$ is colorable, let $C_1, \dots, C_k$ denote its distinct frozen classes. We define the entire integrand below to be zero otherwise. Differentiating in the Poisson means, we find
\[
\chi'(\upsilon) = \E\!\left[\ind\{A^{(\upsilon)}\text{ is colorable}\}\sum_{i = 1}^k\left(\frac{|C_i \cap S_1|^2}{N_1} + \frac{|C_i \cap S_2|^2}{N_2} - \frac{|C_i|^2}{N}\right)\right].
\]
For each such class,
\begin{equation}\label{eq:ov-square}
\frac{|C_i \cap S_1|^2}{N_1} + \frac{|C_i \cap S_2|^2}{N_2} - \frac{|C_i|^2}{N_1 + N_2} = \frac{(N_2 |C_i \cap S_1| - N_1 |C_i \cap S_2|)^2}{N_1 N_2 (N_1 + N_2)} \ge 0.
\end{equation}
It follows that $\rho_N(c) = \chi(c) \ge \chi(0) = \rho_{N_1}(c)\rho_{N_2}(c)$. This is the coloring interpolation of \cite[Theorem~5 and Proposition~2]{bayati2013combinatorial}, in Poisson form as in \cite[Lemma~5.2]{openai2026limiting}.

Fix $\eta \in (0, \eta_0]$, and set $\xi_n = L_n(\eta)/n$ and $\Delta_n = \lceil\sqrt n\log n\rceil$. For $n \ge 2q$ and $c \ge 0$, let $J$ have the $\Pois(c(n - 1))$ law independently of $H$. Splitting loops from other occurrences gives
\[
\rho_n(c) = e^{-c}\Pp(H > J).
\]
For $0 < c \le \Lambda$, \cref{lem:conv-quantile}(b) gives, for all large $n$,
\[
\begin{aligned}
c(n - 1) \le L_n(\eta) - \Delta_n&\ \Longrightarrow\ \rho_n(c) \ge e^{-\Lambda}(1 - \eta)(1 - e^{-(\log n)^2/(2\Lambda)}), \\
c(n - 1) \ge U_n(\eta) + \Delta_n&\ \Longrightarrow\ \rho_n(c) \le \eta + e^{-(\log n)^2/(2\Lambda)}.
\end{aligned}
\]
For large integers $N_1, N_2$, we apply these comparisons at sizes $N_1, N_2, N = N_1 + N_2$ with
\[
c = \min_{j \in \{1, 2\}}\left(\xi_{N_j} - \frac{2\Delta_{N_j}}{N_j}\right).
\]
By the product inequality, we have $\rho_N(c) > e^{-2\Lambda}/2 > 2\eta$. The upper comparison then forces $c(N - 1) < U_N(\eta) + \Delta_N$. We can now use the bound on $w_N(\eta)$ to bound the lower quantile. For real $u \ge 1$, we define $g(u) = K'(\eta)\delta_{\lfloor u\rfloor} + (\Lambda + 6)(\log u)/\sqrt u$. By \Cref{lem:conv-superadd}, for all sufficiently large integers $N_1, N_2$, we have
\begin{equation}\label{eq:ov-superadd}
\xi_{N_1 + N_2} \ge \min(\xi_{N_1}, \xi_{N_2}) - g\bigl(\min(N_1, N_2)\bigr), \qquad \xi_{N_1 + 1} \ge \xi_{N_1} - g(N_1).
\end{equation}
For the adjacent-size inequality, we apply the product inequality with $N_2 = 1$ and $c = \xi_{N_1} - 2\Delta_{N_1}/N_1$, using $\rho_1(c) = e^{-c}$. This is a separate one-vertex comparison; the preceding large-size hypothesis does not allow direct substitution of $N_2 = 1$.

Fix a sufficiently large integer $\ell$ and an integer $n \ge \ell^2$. We write $n$ as a sum of blocks of sizes $\ell$ and $\ell + 1$. We form a binary tree by repeatedly dividing each collection of blocks into two collections whose numbers of blocks differ by at most one. At each internal node, we apply the first inequality of \eqref{eq:ov-superadd} with $(N_1, N_2)$ equal to the total vertex counts in its two children. For a leaf of size $\ell + 1$, we use the second inequality with $N_1 = \ell$. Along each path toward the root, the vertex counts grow by a factor of at least $4/3$. The function $g$ is positive, nonincreasing on $[e^2, \infty)$, and satisfies $\sum_{k \ge 1}g(e^k) < \infty$. Applying \cref{lem:conv-tree} to this $g$, this sequence $(\xi_n)$, and $u_0 = e^2$, we find
\begin{equation}\label{eq:ov-tree}
\xi_n \ge \xi_\ell - g(\ell) - \sum_{i \ge 1} g\Bigl(\frac{\ell}{4} \Bigl(\frac{4}{3}\Bigr)^i\Bigr) \qquad \text{for all } n \ge \ell^2.
\end{equation}
The right side differs from $\xi_\ell$ by an amount that tends to zero as $\ell \to \infty$, since $\sum_k g(e^k) < \infty$. We conclude that $\liminf_n \xi_n \ge \limsup_\ell \xi_\ell$, and $\xi_n$ converges. The summability condition is essential for this convergence argument: \cref{rem:conv-necessity} constructs a bounded, nonconvergent sequence satisfying the same comparisons with a nonsummable error. In particular, a window of order $n/\log n$ would not suffice for this argument.

Finally, for $0 < \eta_- < \eta_+ \le \eta_0$ we have $L_n(\eta_-) \le L_n(\eta_+) \le U_n(\eta_-)$. All levels therefore share one limit $\alpha_q$. Letting $\eta \to 0$ after $n \to \infty$, we get the zero-one law for $\Pp(H > \lfloor cn \rfloor)$, and we apply \cref{lem:conv-transfer}(a) with $m = \lfloor cn\rfloor$ for each fixed $0 \le c < \Lambda$; larger $c$ follow by monotonicity. For a fixed $d \ge 0$, in part (b) we set $p = d/n$ and $m = \lfloor cn\rfloor$, choosing $d/2 < c < \alpha_q$ below the threshold and $\alpha_q < c < d/2$ above it. These substitutions give \cref{thm:intro-main}. Since the window bound holds at every small level, we do not need the sharp threshold theorem of Achlioptas and Friedgut (see \cref{rem:red-af}). This step adapts \cite[Section~5]{openai2026limiting}, which compares expectations rather than quantiles and has polynomially small errors. Summable errors appear in the classical theory of nearly superadditive sequences \cite[Theorem~23]{debruijn1952some}, \cite[arXiv:1006.3786v1, Section~4.3, Lemma~8 and Remark~3]{abbe2014concentration}.

\section{Step 2: the window and the last colorable graph}\label{sec:ov-lifetime}

We fix $\eta \in (0, \eta_0]$. We write $L$, $U$, and $w$ for $L_n(\eta)$, $U_n(\eta)$, and $w_n(\eta)$, respectively. We call $\Gamma = G_{H - 1}$ the \emph{last colorable graph} and call $e_H$ the \emph{final edge}. For a multigraph $G$ on a finite set with at least two vertices, we define $T_G$ to be the number of independent uniform pair additions on that same vertex set needed to destroy colorability, counting the destroying addition and taking $T_G = 0$ if $G$ is already not colorable and $T_G = \infty$ if no addition destroys colorability. For integers $b \ge 0$, we define the \emph{expected lifetime} and its cap by
\[
r(G) = \E T_G, \qquad r_b(G) = \E\min(T_G, b).
\]

We define $Y = \min(U, \max(L, V))$ to be $V$ clipped to the interval $[L, U]$. Then
\[
Y = \min(U, \max(L, V)) = L + \sum_{m = L + 1}^{U} \ind\{G_m \text{ is colorable}\}.
\]
For $1 \le i \le U$, we define $Y^{(-i)}$ by the same sum with $e_i$ omitted from every $G_m$ and all remaining indices unchanged. The omission form of the Efron--Stein inequality \cite[(3.6)]{boucheron2005moment} states that, for a square-integrable function $Z$ of independent random variables $X_1, \dots, X_U$ and square-integrable functions $Z_i$ of the coordinates other than $X_i$, we have
\[
\Var(Z) \le \sum_{i = 1}^U \E[(Z - Z_i)^2].
\]
The variables $Y$ and $Y^{(-i)}$ are bounded functions of the required coordinates. Substituting $X_i = e_i$, $Z = Y$, and $Z_i = Y^{(-i)}$ gives
\[
\Var(Y) \le \sum_{i = 1}^U\E[(Y^{(-i)} - Y)^2].
\]
We evaluate this sum exactly in \Cref{lem:red-deletion}:
\begin{equation}\label{eq:ov-deletion}
\sum_{i = 1}^{U} \E\bigl[(Y^{(-i)} - Y)^2\bigr] = \E\bigl[H \bigl(2 r_{U - H + 1}(\Gamma) - 1\bigr); L < H \le U\bigr].
\end{equation}

The difference $Y^{(-i)} - Y$ counts the times $m \in (L, U]$ at which $G_m$ is not colorable but $G_m$ without the occurrence $e_i$ is colorable. We call a pair $\{u, v\}$ of distinct vertices a \emph{frozen pair} of a colorable multigraph $G$ if $u$ and $v$ lie in the same frozen class. For a colorable multigraph $G$ on a finite vertex set $S$ with $|S| \ge 2$, we define $\kappa(G)$ to be the number of its frozen pairs divided by $\binom{|S|}{2}$. We set $\kappa(G) = 0$ if $G$ is not colorable. Conditionally on the other occurrences up to time $m$, the occurrence $e_i$ is a uniform pair that makes their multigraph $G$ non-colorable with probability $\kappa(G)$. The number of later times at which the difference persists is governed by an independent continuation of $G$, which leads to the weight $2r_{U - m + 1}(G) - 1$ after squaring. The factor $H$ comes from exchangeability.

For every nonnegative function $h$ on multigraphs on $[n]$ and every integer $m \ge 0$, we have the terminal identity
\begin{equation}\label{eq:ov-terminal}
\E\bigl[h(\Gamma); H = m + 1\bigr] = \E\bigl[\kappa(G_m) h(G_m)\bigr].
\end{equation}

More generally, let $\cE$ be a set of pairs $(G, e)$ with $G$ a colorable multigraph on $[n]$ and $e$ a frozen pair of $G$. We define $\kappa_\cE(G)$ to be the fraction of all pairs $e$ for which $(G, e) \in \cE$. For this restriction on the edge, the identity becomes
\[
\E[h(\Gamma);H = m + 1, \ (\Gamma, e_H) \in \cE]
= \E[\kappa_\cE(G_m)h(G_m)].
\]
Taking $\cE$ to contain all frozen pairs gives \eqref{eq:ov-terminal}. By \cref{lem:pre-hazard,lem:pre-terminal}, this identity follows from the fact that $\Pp(H = m + 1 \mid \cF_m) = \kappa(G_m)$. We will use \eqref{eq:ov-terminal} repeatedly, also for final edges of a prescribed kind, to convert statements about $\Gamma$ into statements about the graphs $G_m$ at deterministic times. This step follows the deletion argument of \cite[Lemmas~3.1 and~3.2]{openai2026limiting}, and the role of the lifetime of the last state goes back to Aldous \cite[Theorem~1(b)]{aldous1991threshold} and its application by Carenini \cite[Theorem~4.1]{carenini2026polynomial}. The related argument with weighted deletion appears in \cite[Section~3.3, Lemma~3.4]{openai2026variance}; see \cref{sec:intro-related}.

\Cref{lem:red-varlower} uses an independent copy of $Y$ to show that $\Var(Y) \ge \Pp(Y = L) \Pp(Y = U) w^2 \ge \eta^2 w^2$. On the event $\{L < H \le U\}$, we have $r_{U - H + 1}(\Gamma) \le \min(w, r(\Gamma))$. Inserting this into \eqref{eq:ov-deletion}, we find
\begin{equation}\label{eq:ov-window-lifetime}
\eta^2 w^2 \le \Var(Y) \le 2U \, \E\bigl[\min(w, r(\Gamma)); L < H \le U\bigr].
\end{equation}
To bound the window, it suffices to bound the expected lifetime at final edges inside it. We set $\eps_n = (\log n)^{-3/2}$ for $n \ge 3$ and $\eps_1 = \eps_2 = 1$, and define $x_n = n\eps_n^2$. In the present range $n \ge 2q$, we have $x_n = n(\log n)^{-3}$. For every constant $K_0 \ge 0$, we have by \cref{prop:red-window}
\[
\E[\min(n\eps_n, r(\Gamma));L < H \le U] \le K_0n\eps_n^2
\quad\Longrightarrow\quad
w \le \max(1, 2\Lambda K_0/\eta^2)n\eps_n.
\]
Since
\[
\min(n\eps_n, r(\Gamma)) \le x_n + n\eps_n\ind\{r(\Gamma) > x_n\},
\]
it suffices to prove
\begin{equation}\label{eq:ov-target}
\Pp\bigl(L < H \le U, \ r(\Gamma) > x_n\bigr) \le 2 \eps_n,
\end{equation}
and \cref{prop:red-split} gives the conclusion with $K_0 = 3$.

However, the elementary lifetime bound for a single graph does not give \eqref{eq:ov-target}. For a colorable multigraph $G$ on $[n]$, waiting for a uniform pair that is a frozen pair gives $\kappa(G) \le 1/r(G)$ by \cref{lem:pre-lifetime}. We apply \eqref{eq:ov-terminal} with $h(G) = \ind\{r(G) > x_n\}$ and sum over $L \le m < U$. Since $\kappa(G)r(G) \le 1$, each such time contributes at most $1/x_n$ to the left side of \eqref{eq:ov-target}, and summing over the $M$ times up to the cap gives only the bound $M/x_n \approx \Lambda (\log n)^3$, which is far too large. We instead bound the sums over time directly in Steps 3 and 4. We set $D_n = \lfloor(\log n)^{1/4}\rfloor$. We call an endpoint \emph{low-degree} if its degree in $\Gamma$ is at most $D_n$, and \emph{high-degree} otherwise. We split according to these degrees and prove that, for all large $n$,
\begin{align}
\Pp\Bigl(H < \infty, \ \min_{v \in e_H} \deg_\Gamma(v) \le D_n, \ r(\Gamma) > x_n\Bigr) &\le \eps_n, \label{eq:ov-low} \\
\Pp\Bigl(H \le M, \ \max_{v \in e_H} \deg_\Gamma(v) > D_n\Bigr) &\le \eps_n. \label{eq:ov-high}
\end{align}
These are \cref{cor:low,thm:high-main}. Each final edge in the event of \eqref{eq:ov-target} belongs to one of these events, since $U \le M$. The estimates hold uniformly over the location of the window, which is not known in advance. The first concerns final edges at all times, and the second concerns all final edges up to time $M$, without any restriction on the lifetime. Together with \cref{prop:red-split}, they give $w_n(\eta) \le (6\Lambda/\eta^2) \, n (\log n)^{-3/2}$ for all large $n$, which is \cref{thm:proof-window}.

\section{Step 3: final edges at a vertex of low degree}\label{sec:ov-low}

We now explain the proof of \eqref{eq:ov-low}. \Cref{thm:low-main} shows that there exists $C_{\rm low} > 0$, depending only on $q$, such that for all large $n$ and all $y > 0$,
\begin{equation}\label{eq:ov-low-main}
\Pp\Bigl(H < \infty, \ \min_{v \in e_H} \deg_\Gamma(v) \le D_n, \ r(\Gamma) > y\Bigr) \le C_{\rm low} y^{-1/2} n^{3/8},
\end{equation}
and with $y = x_n$ the right side is $C_{\rm low} (\log n)^{3/2} n^{-1/8}$, which is smaller than $\eps_n$ for all large $n$.

We use the restriction on the lifetime to bound the sizes of the frozen classes. We set $A_q = 2q\log q + 2q + 1$. For a colorable multigraph $G$ on $[n]$, we have $\kappa(G)r(G) \le 1$ and $r(G) \le A_qn$ by \Cref{lem:pre-lifetime}. Applying \cref{lem:pre-classmass} with vertex set $[n]$ and size $n \ge 2q$, we get the following bound whenever $y > 0$ and $r(G) > y$. Let $C_1, \dots, C_k$ denote the frozen classes of $G$.
\begin{equation}\label{eq:ov-classmass}
\max_{1 \le i \le k} \frac{|C_i|^2}{n^2} \le \sum_{i = 1}^k \frac{|C_i|^2}{n^2} = \Bigl(1 - \frac{1}{n}\Bigr) \kappa(G) + \frac{1}{n} \le \frac{1 + A_q}{r(G)} < \frac{1 + A_q}{y},
\end{equation}
so every frozen class of $G$ has at most $n((1 + A_q)/y)^{1/2}$ elements. For each low-degree endpoint of a frozen pair in $\Gamma$, the number of possible partners is one less than the size of its frozen class. In the edge-restricted terminal identity, we take $\cE$ to contain the frozen pairs with an endpoint of degree at most $D_n$, take $h(G) = \ind\{r(G) > y\}$, and sum over all $m \ge 0$. We let $(\Pi_t)_{t \ge 0}$ denote a Poisson counting process of rate one, independent of the pairs and define $G(t) = G_{\Pi_t}$ for $t \ge 0$; we call this the \emph{Poissonized process}. On any other finite vertex set, the same term means the edge process evaluated at an independent Poisson clock of rate one. Applying \cref{lem:pre-poisson} with $\theta = 1$ to the resulting nonnegative function of $G_m$, and using invariance under vertex relabeling, reduces the bound to a fixed vertex. We set $a_0 = n$. We bound the left side of \eqref{eq:ov-low-main} by $(2n/(n - 1))((1 + A_q)/y)^{1/2}$ times
\begin{equation}\label{eq:ov-low-occupation}
\int_0^\infty \Pp\bigl(\deg_{G(t)}(a_0) \le D_n \text{ and } a_0 \text{ lies in a frozen pair of } G(t)\bigr) \dd t.
\end{equation}
It remains to show that this integral, the expected total time that $a_0$ spends in a frozen pair while $\deg_{G(t)}(a_0) \le D_n$, is $O(n^{3/8})$.

Let $G$ be a colorable multigraph on $[n]$, and let $a, b \in [n]$ be distinct vertices in one frozen class. We define $B$ to be the multigraph on $[n] \setminus \{a\}$ formed by the occurrences not containing $a$, and define a \emph{list} to be a finite sequence. We let $\mathbf x$ denote the list of the other endpoints of occurrences containing $a$, counting repetitions, and call it a \emph{neighbor list} of $a$ in $G$. If a coloring of $B$ left two colors unused on this list, it would extend to two colorings of $G$ agreeing at $b$ and disagreeing at $a$. For each coloring $\sigma$ of $B$, we therefore have
\[
b \in C_G(a) \setminus \{a\}
\quad\Longrightarrow\quad
\bigl|\{\sigma(v):v\text{ is an entry of }\mathbf x\}\bigr| \ge q - 1 \ge 2.
\]
We call a list \emph{never monochromatic} in a colorable multigraph if no coloring of that multigraph gives the same color to all entries of the list; a coloring is constant on an empty list by convention. The displayed inequality says that $\mathbf x$ is never monochromatic in $B$. This is the only use of $q \ge 3$ in the window estimates. For $q = 2$, take distinct vertices $a, v, b$ and define a graph by the edges $\{a, v\}$ and $\{v, b\}$. It forces equal colors at $a, b$, but its one-entry neighbor list $(v)$ is constant in every coloring.

For a multigraph $B$ on a finite set and a nonempty subset $Z$ of that set, we define $B/Z$ by replacing all vertices of $Z$ by one new vertex in every occurrence, preserving multiplicity and retaining occurrences internal to $Z$ as loops. We apply \cref{lem:low-frozen}(iv) with $F = B$. By the classical correspondence for identification \cite[Section~3, Theorem~1]{read1968introduction}, we have
\[
B/Z\text{ is colorable}\quad\Longleftrightarrow\quad
\text{some coloring of }B\text{ is constant on }Z.
\]
Taking $Z$ to be the set of entries of $\mathbf x$, we conclude that $B/Z$ is not colorable.

We call the multigraph on $[n - 1]$ formed by the occurrences of $G(t)$ not containing $a_0$ the background graph at time $t$. We let $\mathbf x$ denote the neighbor list of $a_0$ at time $t$. We apply \cref{lem:pre-split}(c) with $S = [n]$ and $a = a_0 = n$, using Poisson splitting \cite[Theorems~5.2 and~5.8]{last2018lectures}. Before conditioning on the event in \eqref{eq:ov-low-occupation}, the background graph and this list are independent. The background graph has the law of the Poissonized process on $[n - 1]$ at time $(n - 2)t/n$, and the list has a $\Pois(2t/n)$ length and, conditional on its length, independent uniform entries in $[n - 1]$. Conditional on a list length at most $D_n$, we enlarge its set of entries by choosing $D_n$ minus the number of distinct entries uniformly from the complement. The resulting set $Z$ has an unconditionally uniform law among the $D_n$-subsets, by invariance under relabeling, and is independent of the background graph. The property of being never monochromatic persists under enlargement. We now let $(B_m)_{m \ge 0}$ denote the time-$m$ graphs of a fresh edge process on $[n - 1]$, independent of a uniform $D_n$-subset $Z$. We define $H^{(n - 1)}$ and $H_Z$ to be the first non-colorable indices of $B_m$ and $B_m/Z$, respectively. By this equivalence, we have
\[
H^{(n - 1)} - H_Z = \sum_{m \ge 0}\ind\{B_m\text{ is colorable}, \ B_m/Z\text{ is not colorable}\}.
\]
We apply \cref{lem:pre-poisson} on $[n - 1]$ with $\theta = (n - 2)/n$ to the function assigning a colorable graph the probability that an independent uniform $D_n$-subset destroys colorability on identification, and assigning zero otherwise. We find
\begin{equation}\label{eq:ov-to-identify}
\int_0^\infty \Pp\bigl(\deg_{G(t)}(a_0) \le D_n \text{ and } a_0 \text{ lies in a frozen pair of } G(t)\bigr) \dd t \le \frac{n}{n - 2} \E\bigl[H^{(n - 1)} - H_Z\bigr].
\end{equation}
The right side is $n/(n - 2)$ times the expected advance of non-colorability caused by identifying $D_n$ random vertices.

\Cref{prop:low-identify} bounds this advance. Let $N \ge \lceil4q\log q\rceil + 1$ and $1 \le D \le N/4$ be integers. We set $z_0 = 4q\log q + 1$ and $C_q = 2e^{z_0}$, and let $k > eDz_0$ be an integer. In this comparison, we use $(B_m)_{m \ge 0}$ for the time-$m$ graphs of an edge process on $[N]$. We take $Z$ to be an independent uniform $D$-subset of $[N]$, and define $H^{(N)}$ and $H_Z$ as the first non-colorable indices of $B_m$ and $B_m/Z$. Then
\begin{equation}\label{eq:ov-identify}
\E\bigl[H^{(N)} - H_Z\bigr] \le C_q \Bigl(D^k + N \Bigl(\frac{e D z_0}{k}\Bigr)^k + D^2 + 1\Bigr).
\end{equation}
For all sufficiently large $n$, we set $k_n = \lceil\log n/\log\log n\rceil$ and apply \eqref{eq:ov-identify} with $(N, D, k) = (n - 1, D_n, k_n)$. These values satisfy its integer, size, and truncation hypotheses, and
\[
D_n^{k_n} \le n^{1/4}(\log n)^{1/4}, \qquad
(n - 1)\left(\frac{eD_nz_0}{k_n}\right)^{k_n} \le n^{3/8}.
\]
Using $D_n^2 \le (\log n)^{1/2}$, we get
\[
\E[H^{(n - 1)} - H_Z] \le 2C_qn^{3/8}.
\]
Together with the class-mass bound, the reduction to \eqref{eq:ov-low-occupation}, and \eqref{eq:ov-to-identify}, this proves \eqref{eq:ov-low-main}.

To prove \eqref{eq:ov-identify}, we compare the lists of neighbors before and after identification. After Poissonizing the edge process on $[N]$, we bound the part of the time integral beyond $2q(\log q)N$ by $2q$. Fix $0 \le s \le 2q(\log q)N$ and a set of $D$ vertices to identify. We call the multigraph of occurrences internal to the remaining $N - D$ vertices the background graph. Its pair counts are independent with law $\Pois(s/\binom N2)$. We define $\hat z$ to be the mean length of each neighbor list from an identified vertex into the remaining set. By Poisson splitting, we have
\[
\hat z = \frac{(N - D)s}{\binom N2} \le z_0.
\]
The neighbor lists have independent $\Pois(\hat z)$ lengths, conditionally independent uniform entries in the remaining set, and are independent of the background graph. Occurrences internal to the identified set contribute $O_q(D^2)$ to the time integral. In their absence, identification gives the background graph together with a new vertex whose neighbor list has $\Pois(D\hat z)$ independent uniform entries.

For comparison, replace the new vertex's neighbor list by an independent list of $\Pois(\hat z)$ uniform entries of the remaining set. We define $F(t)$ to be the Poissonized process on the remaining $N - D$ vertices together with that new vertex, and set $\lambda' = \binom{N - D + 1}{2}/\binom N2$. The resulting graph has the law of $F(\lambda's)$. We apply \cref{lem:low-added} at the same time $s$ with the explicit substitutions
\[
\begin{gathered}
V' = \text{remaining vertex set}, \quad N' = N - D, \quad o = \text{new vertex}, \\
\lambda = 1/\tbinom N2, \quad B' = \text{background graph}.
\end{gathered}
\]
It bounds the probability that this replacement list destroys colorability of a colorable background graph by $(e^{\hat z} - 1)\E\kappa(F(\lambda's))$. Since the process loses colorability at most once, we can integrate the probability that one further edge destroys colorability to get
\[
\int_0^\infty\E\kappa(F(\lambda's))\,\dd s \le \frac1{\lambda'} \le \frac{16}{9}.
\]
We use the convention $0^0 = 1$. For $z \ge 0$ and integers $j \ge 0$, we define $\pi_z(j) = e^{-z}z^j/j!$. To compare the two list lengths, we use
\begin{equation}\label{eq:ov-tilt}
\pi_{D\hat z}(j) = D^j e^{-(D - 1)\hat z}\pi_{\hat z}(j) \le D^k\pi_{\hat z}(j)\qquad\text{for integers }0 \le j \le k.
\end{equation}
We bound the probability that the identified vertex has more than $k$ neighbors by $(eDz_0/k)^k$. Together with \eqref{eq:ov-tilt}, this accounts for the terms in \eqref{eq:ov-identify} from the length comparison, long lists, occurrences inside the identified set, and large times. The limitation of this comparison is quantified in \cref{rem:low-limit}: at degree of order $\sqrt{\log n}$, its two principal terms cannot both be $o(\sqrt n)$, whereas the factor $((1 + A_q)/y)^{1/2}$ used before \eqref{eq:ov-low-occupation} is of order $n^{-1/2}(\log n)^{3/2}$ at $y = x_n$. This explains why we use a different argument for the remaining degrees.

\section{Step 4: final edges at vertices of high degree}\label{sec:ov-high}

We now prove \eqref{eq:ov-high}. For a colorable multigraph $G$ on $[n]$ and an integer $j \ge 0$, we define $\kappa_j(G)$ to be the number of ordered pairs $(u, v)$ of distinct vertices in one frozen class, which we call \emph{ordered frozen pairs}, with $\deg_G(u) = j$, divided by $n(n - 1)$. We set $\kappa_j(G) = 0$ for non-colorable $G$. In the edge-restricted terminal identity, we take $\cE$ to contain the frozen pairs with an endpoint of degree greater than $D_n$, take $h = 1$, and sum over $0 \le m < M$. Counting each eligible unordered pair through at least one of its endpoints gives
\[
\Pp(H \le M, \ \max_{v \in e_H}\deg_\Gamma(v) > D_n)
\le 2\sum_{m = 0}^{M - 1}\sum_{j = D_n + 1}^{m}\E\kappa_j(G_m).
\]
The identification estimate above does not cover all these degrees. We instead delete occurrences. For a colorable $G$ and an ordered frozen pair $(u, v)$, we call an occurrence $e$ containing $u$ \emph{non-pivotal} if the pair remains frozen in $G - e$, and \emph{pivotal} otherwise. We set $K = \lceil8\Lambda\rceil$. We define $\kappa_j^{\rm piv}(G)$ by restricting the count in $\kappa_j(G)$ to pairs for which at most $K - 1$ occurrences at $u$ are non-pivotal, and set it to zero for non-colorable $G$.

Let $f$ be a uniform pair independent of the process. For integers $1 \le i \le m$, the pair $(G_m - e_i, e_i)$ has the law of $(G_{m - 1}, f)$; this is the analogue of the Mecke equation for a fixed number of independent points \cite[Theorem~4.1]{last2018lectures}. Deleting one of at least $K$ occurrences at $u$ that are non-pivotal for $(u, v)$ leaves $u$ and $v$ in one frozen class and lowers the degree of $u$ by one. A uniform pair contains $u$ with probability $2/n$. Using \cref{lem:high-deletion}(a), we find
\begin{equation}\label{eq:ov-descent}
\E \kappa_j(G_m) \le \E \kappa_j^{\rm piv}(G_m) + \frac{2m}{nK} \E \kappa_{j - 1}(G_{m - 1}) \le \E \kappa_j^{\rm piv}(G_m) + \frac{1}{4} \E \kappa_{j - 1}(G_{m - 1})
\end{equation}
for $1 \le j \le m \le M$, where we used $m \le M \le \Lambda n$ and $K \ge 8\Lambda$. Let $w_0 = \lfloor D_n/2 \rfloor$. Iterating \eqref{eq:ov-descent} from the degree $j$ down to $w_0$, we find that for $w_0 < j \le m < M$,
\begin{equation}\label{eq:ov-iterate}
\E \kappa_j(G_m) \le \sum_{j' = w_0 + 1}^{j} 4^{-(j - j')} \E \kappa_{j'}^{\rm piv}(G_{m - j + j'}) + 4^{-(j - w_0)} \E \kappa_{w_0}(G_{m - j + w_0}).
\end{equation}
Summing \eqref{eq:ov-terminal} with $h = 1$ over $m$ shows that $\sum_m \E \kappa(G_m) \le 1$, since the process becomes non-colorable at most once. Therefore the last terms in \eqref{eq:ov-iterate} contribute at most $(8/3) 4^{-(D_n + 1 - w_0)}$. It remains to bound the pivotal terms, whose degrees exceed $w_0$.

If a pivotal occurrence for $(u, v)$ joins $u$ to a vertex $y$, deleting it separates $u$ and $v$ into different frozen classes. Adding it back merges those classes. For the rest of this step, we fix a vertex $u \in [n]$ and call it the \emph{root}. For a colorable multigraph $G$ on $[n]$, we set $W = [n] \setminus C_G(u)$. We define $R_G(u)$ to be the set of ordered pairs $(y, v) \in W^2$ such that every coloring gives $u$ the color of $y$ or of $v$. We define its density $\beta_G(u) = |R_G(u)|/n^2$. For every $y \in [n]$, we define the \emph{section} $R_G(u;y) = \{v \in [n]:(y, v) \in R_G(u)\}$. When $G$ and $u$ are fixed, we write $R(y)$ for $R_G(u;y)$. We set $R_G(u) = \emptyset$ and $\beta_G(u) = 0$ when $G$ is not colorable. Thus
\begin{equation}\label{eq:ov-relation}
\begin{gathered}
R_G(u) = \bigl\{(y, v) \in W \times W : \sigma(u) \in \{\sigma(y), \sigma(v)\} \text{ for every coloring } \sigma \text{ of } G\bigr\}, \\
\beta_G(u) = \frac{|R_G(u)|}{n^2}, \qquad R(y) = \{v : (y, v) \in R_G(u)\}.
\end{gathered}
\end{equation}
The relation is symmetric. For $y \in W$, \cref{lem:high-relation} shows that $G + \{u, y\}$ is colorable and
\begin{equation}\label{eq:ov-section}
C_{G + \{u, y\}}(u) = C_G(u) \cup R(y) \qquad \text{for all } y \in W.
\end{equation}
Indeed, the colorings of $G + \{u, y\}$ are the colorings $\sigma$ of $G$ with $\sigma(y) \ne \sigma(u)$, and for $v \in W$, all of them satisfy $\sigma(v) = \sigma(u)$ if and only if $(y, v) \in R_G(u)$. Hence $n \beta_G(u)$ is the average of $|R(y)|$ over $y \in [n]$, and for $y \in W$, the number $|R(y)|$ counts the vertices that join the frozen class of $u$ when $\{u, y\}$ is added. Note that a pair counted in $\kappa_j^{\rm piv}(G_m)$ has at least $j - K + 1$ pivotal occurrences at $u$. Deleting one of them and using exchangeability as before, we get the following bound from \cref{lem:high-deletion}(b) for integers $m \ge 1$, $j \ge K$, and all $u \in [n]$:
\begin{equation}\label{eq:ov-pivotal}
\E \kappa_j^{\rm piv}(G_m) \le \frac{2mn}{(n - 1)^2 (j - K + 1)} \, \E\bigl[\ind\{\deg_{G_{m - 1}}(u) = j - 1\} \beta_{G_{m - 1}}(u)\bigr].
\end{equation}
Suppose $w_0 \ge K$. Summing \eqref{eq:ov-pivotal} over $w_0 < j \le m < M$ and using the geometric descent, we get
\[
\begin{split}
\Pp(H \le M, \ \max_{v \in e_H}\deg_\Gamma(v) > D_n)
\le \frac83\,4^{-(D_n + 1 - w_0)}\\
{} + \frac{16\Lambda}{3(w_0 - K + 2)}\left(\frac n{n - 1}\right)^2
\sum_{m = 0}^{M - 1}\E[\ind\{\deg_{G_m}(u) \ge w_0\}\beta_{G_m}(u)].
\end{split}
\]
It therefore suffices to bound the sum
\begin{equation}\label{eq:ov-occupation}
\sum_{m = 0}^{M - 1}\E[\ind\{\deg_{G_m}(u) \ge w_0\}\beta_{G_m}(u)].
\end{equation}

The main difficulty is that $\beta_G(u)$ can be of order one for a single colorable graph. For the example of \cref{rem:high-killing}, take $q = 3$, $n \ge 7$, and $n \equiv 1\pmod3$. Let $G$ consist of a triangle, an edge from $u$ to one vertex of that triangle, and, for each triangle vertex, $(n - 4)/3$ additional vertices joined to the other two triangle vertices. Every additional vertex receives the color of the triangle vertex to which it is not joined. The vertex $u$ can take two colors and $C_G(u) = \{u\}$. Here
\[
\beta_G(u) = \frac{2(n - 1)^2}{9n^2}\longrightarrow\frac29.
\]
Adding occurrences from $u$ to vertices forced to have its neighbor's color can increase its degree without changing this relation. Heuristically, this obstructs a useful uniform bound on $\beta_G(u)$ based only on a large degree. For $0 \le m < M$, the elementary degree estimate gives
\[
\Pp(\deg_{G_m}(u) \ge w_0) \le \left(\frac{2e\Lambda}{w_0}\right)^{w_0}.
\]
Bounding $\beta_{G_m}(u)$ by one at each time and summing would give only an upper bound of order $n(2e\Lambda/w_0)^{w_0}$, which does not tend to zero. We instead control the total time spent in such states through growth of the frozen class. In the example, the fraction of frozen pairs tends to one third, so the expected lifetime stays bounded.

For a colorable $G$ and a pair $f = \{a, b\}$, we define its \emph{credited set} $\mathscr A_G(f)$ as follows. If $a, b \in W$, we set $\mathscr A_G(f) = R(a) \cap R(b)$. If exactly one endpoint lies in $W$, we set $\mathscr A_G(f)$ equal to the section of that endpoint. If both endpoints lie in $C_G(u)$, we set $\mathscr A_G(f) = \emptyset$. This definition includes pairs for which $G + f$ is not colorable. We let $\dot \cup$ denote a union of disjoint sets. When $G + f$ is colorable, we apply \cref{lem:high-absorb} with this $G$, the fixed root $u$, and this pair $f$ to get
\[
C_G(u)\,\dot \cup \,\mathscr A_G(f) \subseteq C_{G + f}(u), \qquad
|C_{G + f}(u)| \ge |C_G(u)| + |\mathscr A_G(f)|,
\]
For example, if $a, b \in W$ and $y \in R(a) \cap R(b)$, symmetry gives $a, b \in R(y)$. By \eqref{eq:ov-section}, adding $\{u, y\}$ puts $a, b$ into the class of $u$. The graph $G + \{u, y\} + \{a, b\}$ is then not colorable. Every coloring of $G + \{a, b\}$ must give $y$ the color of $u$.

We now take $f$ uniformly among pairs of distinct vertices of $[n]$, and let $\E_f$ denote expectation over this pair. The main combinatorial estimate is
\begin{equation}\label{eq:ov-onestep}
\beta_G(u) \le \frac{2}{\log 2} \, \E_f \log\Bigl(1 + \frac{|\mathscr A_G(f)|}{|C_G(u)|}\Bigr),
\end{equation}
which bounds the density by the expected credited logarithmic increment. The increment is bounded by the actual growth of the logarithm of the class size when the added pair preserves colorability, and remains defined when it destroys colorability.
To prove this estimate, we use the following inequality from \cref{lem:high-comb}, with the current graph $G$ and root $u$:
\[
\begin{aligned}
|R_G(u)| &\le
2\sum_{(a, b) \in W \times W}\min\left(\frac{|R(a) \cap R(b)|}{|C_G(u)|}, 1\right)\\
&\quad + 2(|C_G(u)| - 1)\sum_{a \in W}\min\left(\frac{|R(a)|}{|C_G(u)|}, 1\right).
\end{aligned}
\]
Its proof uses symmetry to count relation pairs incident to a section of size at least $|C_G(u)|$ through the first sum; pairs for which both sections are smaller are bounded by the second sum. Counting ordered pairs of distinct vertices in the definition of $\mathscr A_G(f)$ then gives
\[
\E_f\min\left(\frac{|\mathscr A_G(f)|}{|C_G(u)|}, 1\right)
\ge \frac{|R_G(u)|}{2n(n - 1)} \ge \frac{\beta_G(u)}2.
\]
Here pairs with one endpoint in $C_G(u)$ supply the second contribution and compensate for the diagonal terms in $W \times W$, which are not possible added edges. Applying $\log(1 + z) \ge (\log 2)\min(z, 1)$ for $z \ge 0$, with $z = |\mathscr A_G(f)|/|C_G(u)|$, proves \eqref{eq:ov-onestep}.

We define $\tau = \min(H, M)$. For integers $0 \le m < H$, we define the \emph{potential} $\Phi_m = \log|C_{G_m}(u)|$ and the \emph{credited logarithmic increment} $\Psi_m = \log(1 + |\mathscr A_{G_m}(e_{m + 1})|/|C_{G_m}(u)|)$. We set $\Phi_m = \Psi_m = 0$ for $m \ge H$. The potential takes values in $[0, \log n]$. By absorption, we have
\[
\Phi_{m + 1} - \Phi_m \ge \Psi_m\qquad\text{when }m + 1 < H.
\]
At the last index before $\tau$, disjointness inside $[n]$ gives
\[
\Psi_{\tau - 1} \le \log n - \Phi_{\tau - 1}.
\]
The pair $e_\tau$ is the final edge if $\tau = H$ and the edge at the cap if $\tau = M < H$. Let $\tau_0$ be a stopping time of $(\cF_m)_{m \ge 0}$ with values in $\{0, 1, \dots\} \cup \{\infty\}$. By telescoping, we have
\[
\sum_{m : \, \tau_0 \le m < \tau} \Psi_m \le (\log n) \ind\{\tau_0 < \tau\}.
\]
For $0 \le m < M$, we apply \eqref{eq:ov-onestep} conditionally on $\cF_m$ with $G = G_m$ and $f = e_{m + 1}$ on the event $m < H$. The fresh pair is uniform and independent of $\cF_m$. Multiplying by the $\cF_m$-measurable indicator of $\{\tau_0 \le m < \tau\}$, summing, and taking expectations gives
\begin{equation}\label{eq:ov-potential}
\E \sum_{m : \, \tau_0 \le m < \tau} \beta_{G_m}(u) \le \frac{2}{\log 2} \, (\log n) \, \Pp(\tau_0 < \tau).
\end{equation}
This estimate controls the total occupation of the process. The final edge must be credited. In the example above, pairs $f$ for which $G + f$ is colorable contribute only $O((\log n)/n)$ to \eqref{eq:ov-onestep}, although $\beta_G(u) \to 2/9$; see \cref{rem:high-killing}.

We now take $\tau_0 = \min\{m \ge 0:\deg_{G_m}(u) \ge w_0\}$, with value $\infty$ if the set is empty. Since degrees are nondecreasing and $\beta_{G_m}(u) = 0$ after $H$, the left side of \eqref{eq:ov-potential} is \eqref{eq:ov-occupation}. For $w_0 \ge K$, applying the binomial tail estimate with number of trials $M - 1$, success probability $2/n$, and threshold $w_0$ gives
\[
\Pp(\tau_0 < \tau) \le \Pp(\Bin(M - 1, 2/n) \ge w_0)
\le \left(\frac{2e\Lambda}{w_0}\right)^{w_0}.
\]
Inserting this bound into \eqref{eq:ov-potential} gives the estimate of \cref{cor:high-degree} with $j = w_0$. This argument uses no independence between the degree and the relation density. Combining it with the descent estimate proves, for $w_0 \ge K$,
\begin{equation}\label{eq:ov-high-final}
\begin{split}
&\Pp\Bigl(H \le M, \ \max_{v \in e_H} \deg_\Gamma(v) > D_n\Bigr) \\
&\qquad \le \frac{8}{3} \, 4^{-(D_n + 1 - w_0)} + \frac{32 \Lambda}{3 \log 2} \Bigl(\frac{n}{n - 1}\Bigr)^2 \frac{\log n}{w_0 - K + 2} \Bigl(\frac{2e\Lambda}{w_0}\Bigr)^{w_0}.
\end{split}
\end{equation}
Since $w_0 \ge (\log n)^{1/4}/2 - 1$, the bound is at most $\eps_n$ for all large $n$, proving \eqref{eq:ov-high}. The argument also holds for $q = 2$ (\cref{rem:high-q2}).

\chapter{Preliminaries}\label{sec:pre}

This section gives the definitions and elementary facts used throughout the paper. We introduce the random graph process and its quantiles, record identities for the step at which colorability is lost, and prove the elementary estimates used below. The assumption $q \ge 3$ is used to color unicyclic components in \cref{lem:pre-cycles}(a), which supplies the lower bound in \cref{lem:pre-centre}(b). All remaining statements of this section hold for $q = 2$ as well.

Throughout the paper, we fix an integer $q \ge 3$. For each integer $k \ge 1$, we define $[k] = \{1, \dots, k\}$. We let $\log$ denote the natural logarithm. The constants in our statements may depend on $q$ and on the fixed parameters named in the statement, but not on $n$. Constants denoted by $C$ or $C'$ are local to the statement or proof in which they occur. We say that a statement holds \emph{for all large $n$} if there exists an integer $n_0$, depending only on $q$ and those fixed parameters, such that it holds for every $n \ge n_0$. Each assertion involving a fixed parameter holds separately at each value of that parameter; its onset may depend on that value.

\section{The edge process}\label{sec:pre-process}

We begin with multigraphs and their colorings. We will sample pairs of vertices with replacement and count degrees with multiplicity.

\begin{definition}\label{def:pre-multigraph}
Let $S$ be a finite set. We call a set $\{u, v\}$ with $u, v \in S$ and $u \ne v$ a \emph{pair} of $S$. We define a \emph{multigraph on $S$} to be a finite multiset $G$ of pairs of $S$. We call the elements of $G$, counted with multiplicity, the \emph{occurrences} of $G$. For a pair $e$ of $S$, we let $G + e$ denote the multigraph obtained from $G$ by adding one occurrence of $e$, and, if $e$ occurs in $G$, we let $G - e$ denote the multigraph obtained from $G$ by removing one occurrence of $e$. For pairs $p_1, \dots, p_k$ of $S$, we write $G + p_1 + \dots + p_k$ for the multigraph obtained from $G$ by adding one occurrence of each $p_i$, and $p_1 + \dots + p_k$ for the multigraph with occurrences $p_1, \dots, p_k$ (the empty multigraph if $k = 0$). We write $G \subset G'$ if every pair has at most the same multiplicity in $G$ as in $G'$. For $u \in S$, we define the \emph{degree} $\deg_G(u)$ to be the number of occurrences of $G$ that contain $u$. For $S' \subset S$, we let $G[S']$ denote the multigraph on $S'$ formed by the occurrences of $G$ that are pairs of $S'$. For a bijection $\phi \colon S \to S'$, we let $\phi G$ denote the multigraph on $S'$ whose occurrences are the pairs $\{\phi(u), \phi(v)\}$, where $\{u, v\}$ runs over the occurrences of $G$. We call a map $\sigma\colon S \to [q]$ a \emph{coloring} of $G$ if $\sigma(u) \ne \sigma(v)$ for every occurrence $\{u, v\}$ of $G$. We let $\Col(G)$ denote the set of colorings of $G$, and we say that $G$ is \emph{colorable} if $\Col(G) \neq \emptyset$.
\end{definition}

A multigraph has the same colorings as the simple graph formed by its distinct pairs. Repeated occurrences therefore do not affect colorability. They do affect degrees, and in this paper degrees always count occurrences. Note also that $\Col(G') \subset \Col(G)$ whenever $G \subset G'$.

\begin{definition}\label{def:pre-process}
Let $S$ be a finite set with at least two elements. We define an \emph{edge process on $S$} to be a sequence of independent random pairs of $S$, each uniformly distributed on the set of its $\binom{|S|}{2}$ pairs. For each integer $m \ge 0$, we call the multigraph whose occurrences are the first $m$ terms its \emph{time-$m$ graph}.
\end{definition}

Let $n \ge 2q$ be an integer. Throughout the paper, we let $(e_i)_{i \ge 1}$ denote an edge process on $[n]$. For each integer $m \ge 0$, we define $G_m = e_1 + \dots + e_m$. We let $\cF_m$ denote the $\sigma$-algebra generated by $e_1, \dots, e_m$. A vertex $u \in [n]$ belongs to $e_i$ with probability $(n - 1)/\binom{n}{2} = 2/n$, and $\deg_{G_m}(u)$ has the law $\Bin(m, 2/n)$. We define the hitting time of non-colorability
\begin{equation}\label{eq:pre-H}
H = \min\{m \ge 0 : G_m \text{ is not colorable}\}.
\end{equation}
We use the convention $\min\emptyset = \infty$. We set $A_q = 2q \log q + 2q + 1$. Since $G_0$ is empty, we have $H \ge 1$, and we will show in \cref{lem:pre-lifetime}(d) that $\E H \le A_q n$; in particular, $H$ is finite almost surely. On the event $\{H < \infty\}$, we let $\Gamma = G_{H - 1}$ denote the last colorable graph of the process, and we call $e_H$ the \emph{final edge}. Note that $\Pp(H > m) = \Pp(G_m \text{ is colorable})$ for all $m \ge 0$, and that this probability is nonincreasing in $m$.

We set
\begin{equation}\label{eq:pre-gamma}
\gamma_q = \frac{\log q}{\log(q/(q - 1))}.
\end{equation}
For every real $c_1 > \gamma_q$, the first-moment bound of \cref{lem:pre-centre}(a) implies that the probability that $G_{\lfloor c_1 n\rfloor}$ is colorable tends to zero. Thus $\gamma_q$ is an upper bound for the threshold, measured in edges per vertex. We fix a real number $\Lambda > \gamma_q + 1$. Throughout the paper, we allow constants and the onset of statements for all large $n$ to depend on this fixed $\Lambda$. In the final proofs, we set $\Lambda = \gamma_q + 2$. We set $M = \lfloor \Lambda n \rfloor$ and
\begin{equation}\label{eq:pre-V}
V = \min(H - 1, M).
\end{equation}
The cap $M$ makes $V$ a bounded function of $e_1, \dots, e_M$, which is convenient for variance bounds. By \cref{lem:pre-centre}(a), the quantiles defined below are smaller than $M$ for all large $n$.

\begin{definition}\label{def:pre-quantile}
For $\eta \in (0, 1/2]$, we define the \emph{lower quantile} $L_n(\eta)$ and the \emph{upper quantile} $U_n(\eta)$ by
\begin{equation}\label{eq:pre-quantile}
\begin{aligned}
L_n(\eta) &= \min\{m \ge 0 : \Pp(V \le m) \ge \eta\}, \\
U_n(\eta) &= \max\{m \in \{0, \dots, M\} : \Pp(V \ge m) \ge \eta\},
\end{aligned}
\end{equation}
where $m$ ranges over integers, and we call $w_n(\eta) = U_n(\eta) - L_n(\eta)$ the \emph{window} at level $\eta$. We also set $\eta_0 = e^{-2\Lambda}/8$.
\end{definition}

The two sets in \eqref{eq:pre-quantile} are nonempty, since $\Pp(V \le M) = 1$ and $\Pp(V \ge 0) = 1$. The first identity also implies that $L_n(\eta) \le M$. In later sections, we restrict the quantile level to $\eta \in (0, \eta_0]$. We use the value of $\eta_0$ in \cref{app:conv}. When $n$ and $\eta$ are fixed, we write $L$, $U$, and $w$ for $L_n(\eta)$, $U_n(\eta)$, and $w_n(\eta)$.

\begin{lemma}\label{lem:pre-quantile}
Let $\eta \in (0, 1/2]$. Then $L \le U \le M$. Further, we have
\begin{equation}\label{eq:pre-quantile-tails}
\begin{aligned}
\Pp(H > m) &> 1 - \eta \qquad \text{for all integers } 0 \le m \le L, \\
\Pp(H > m) &< \eta \qquad \text{for all integers } U < m \le M.
\end{aligned}
\end{equation}
\end{lemma}

\begin{proof}
By the minimality of $L$, we have $\Pp(V \le L - 1) < \eta$ (if $L = 0$, this holds because $V \ge 0$). Hence
\[
\Pp(V \ge L) = 1 - \Pp(V \le L - 1) > 1 - \eta \ge \eta,
\]
where we used $\eta \le 1/2$ in the last step. Since also $L \le M$, the integer $L$ belongs to the set that defines $U$, and therefore $L \le U$. The bound $U \le M$ holds by the definition of $U$.

Next, we observe that for every integer $m \le M$,
\begin{equation}\label{eq:pre-V-H}
\{V \ge m\} = \{H - 1 \ge m\} = \{H > m\}.
\end{equation}
Taking $m = L$ in \eqref{eq:pre-V-H}, we find that
\[
\Pp(H > L) = \Pp(V \ge L) > 1 - \eta,
\]
and the first bound in \eqref{eq:pre-quantile-tails} follows because $\Pp(H > m)$ is nonincreasing in $m$. Now let $U < m \le M$. By the maximality of $U$, the integer $m$ does not belong to the set that defines $U$, and \eqref{eq:pre-V-H} implies that $\Pp(H > m) = \Pp(V \ge m) < \eta$.
\end{proof}

\section{Frozen pairs}\label{sec:pre-frozen}

The final edge of the process joins two vertices that receive the same color in every coloring of the last colorable graph. We now introduce the corresponding notions.

\begin{definition}\label{def:pre-frozen}
Let $S$ be a finite set with $N \ge 2$ elements, and let $G$ be a multigraph on $S$. For $u, v \in S$, we write $u \sim _G v$ if $G$ is colorable and $\sigma(u) = \sigma(v)$ for all $\sigma \in \Col(G)$. We call a pair $\{u, v\}$ of $S$ a \emph{frozen pair} of $G$ if $u \sim _Gv$. If $G$ is colorable, we call the set $C_G(u) = \{v \in S : u \sim _G v\}$ the \emph{frozen class} of $u \in S$. If $G$ is colorable, let $C_1, \dots, C_k$ denote its distinct frozen classes. We define
\begin{equation}\label{eq:pre-kappa}
\kappa(G) = \binom{N}{2}^{-1} \#\{\text{frozen pairs of } G\}, \qquad Q(G) = \sum_{i = 1}^k \Big(\frac{|C_i|}{N}\Big)^2,
\end{equation}
and we set $Q(G) = 0$ if $G$ is not colorable.
\end{definition}

The relation $\sim _G$ is an equivalence relation on $S$ whenever $G$ is colorable, since it is defined by equalities of colors. Its classes are the frozen classes, and $u \in C_G(u)$ for all $u \in S$. We say that $u$ \emph{lies in a frozen pair} of $G$ if $|C_G(u)| \ge 2$. We have $\kappa(G) = 0$ when $G$ is not colorable, since $\sim _G$ is then empty. Observe that $\kappa(G)$ is the probability that a uniformly random pair of $S$ is a frozen pair of $G$. A frozen pair of $G$ is not an occurrence of $G$, since every coloring of $G$ gives different colors to the endpoints of each occurrence.

Throughout the paper, we define numerical statistics constructed from colorings or frozen classes to be $0$ when the multigraph is not colorable, and set-valued statistics to be the empty set. We evaluate expressions involving frozen classes, including logarithms and quotients, on the colorable event, and define the entire expression to be $0$ otherwise. In particular, $\Col(G) = \emptyset$ when $G$ is not colorable. We use this convention in every identity involving these statistics.

Culberson and Gent studied frozen pairs in random graph coloring and the graph obtained by merging their equivalence classes \cite[Sections~2.1 and~3.2]{culberson2001frozen}. Their tests for frozen and free pairs use edge addition and vertex identification, respectively \cite[Section~2.2]{culberson2001frozen}. Our relation $\sim _G$ quantifies over all colorings of $G$.

\begin{lemma}\label{lem:pre-frozen}
Let $S$ be a finite set with $N \ge 2$ elements, and let $G$ be a multigraph on $S$.
\begin{enumerate}[label=(\alph*)]
\item For all distinct $u, v \in S$, the multigraph $G + \{u, v\}$ is colorable if and only if $G$ is colorable and $u \not \sim _G v$. If $f$ is a uniformly random pair of $S$, then
\[
\Pp(G \text{ is colorable and } G + f \text{ is not colorable}) = \kappa(G).
\]
\item Let $G'$ be a colorable multigraph on $S$ with $G \subset G'$. Then $G$ is colorable, $u \sim _G v$ implies $u \sim _{G'} v$ for all $u, v \in S$, $C_G(u) \subset C_{G'}(u)$ for all $u \in S$, and $\kappa(G) \le \kappa(G')$.
\item If $G$ is colorable, then
\[
Q(G) = \Big(1 - \frac 1N\Big) \kappa(G) + \frac 1N,
\]
and $|C_G(u)| \le N Q(G)^{1/2}$ for all $u \in S$.
\end{enumerate}
\end{lemma}

\begin{proof}
We begin with (a). Observe that
\[
\Col(G + \{u, v\}) = \{\sigma \in \Col(G) : \sigma(u) \neq \sigma(v)\}.
\]
This set is empty if and only if either $\Col(G) = \emptyset$, or $\Col(G) \neq \emptyset$ and $\sigma(u) = \sigma(v)$ for all $\sigma \in \Col(G)$. By the definition of $\sim _G$, the second alternative means that $u \sim _G v$. The second statement of (a) follows by averaging over the $\binom{N}{2}$ values of $f$.

To prove (b), we use the inclusion $\Col(G') \subset \Col(G)$, which implies that $G$ is colorable. Suppose that $u \sim _G v$. Then every $\sigma \in \Col(G')$ is a coloring of $G$ and therefore satisfies $\sigma(u) = \sigma(v)$. This means that $u \sim _{G'} v$. The remaining statements of (b) follow.

Finally, let $C_1, \dots, C_k$ be the distinct frozen classes of $G$. A frozen pair is a pair of distinct elements of one class, so
\[
\sum_{i = 1}^k |C_i|^2 = \sum_{i = 1}^k |C_i| (|C_i| - 1) + \sum_{i = 1}^k |C_i| = 2 \#\{\text{frozen pairs of } G\} + N = N (N - 1) \kappa(G) + N.
\]
Dividing by $N^2$, we obtain the formula for $Q(G)$ in (c). The bound on $|C_G(u)|$ follows from $(|C_G(u)|/N)^2 \le Q(G)$.
\end{proof}

\Cref{lem:pre-frozen}(a) is the elementary case of the addition--identification relation for chromatic polynomials \cite[Section~3, Theorem~1]{read1968introduction}. For nonadjacent $u$ and $v$, this relation splits the colorings of $G$ into the colorings of $G + \{u, v\}$ and the colorings that are constant on $\{u, v\}$. The latter correspond to the colorings of the graph obtained from $G$ by identifying $u$ and $v$.

\section{The final edge}\label{sec:pre-hazard}

We now express the law of the final edge $e_H$ through the frozen pairs of the graphs $G_m$ at deterministic times. Culberson and Gent used frozen-pair counts to compute conditional one-edge survival probabilities in a process without replacement \cite[Section~3.3]{culberson2001frozen}; here independent sampling with replacement gives the normalization $\binom{n}{2}^{-1}$. For random $2$-SAT, a related conditional probability is computed in \cite[Section~3.3]{bollobas2001scaling}; our proofs follow the same idea. They are short applications of \cref{lem:pre-frozen}(a).

\begin{lemma}\label{lem:pre-hazard}
For all integers $m \ge 0$, we have
\begin{equation}\label{eq:pre-hazard}
\Pp(H = m + 1 \mid \cF_m) = \kappa(G_m) \qquad \text{almost surely}.
\end{equation}
For every $\cF_m$-measurable random variable $\zeta \ge 0$, we have $\E[\zeta \kappa(G_m)] = \E[\zeta \ind\{H = m + 1\}]$. In particular, if $\mathcal J$ is a set of nonnegative integers and $\zeta_m$ is an $\cF_m$-measurable random variable with values in $[0, 1]$ for each $m \in \mathcal J$, then
\begin{equation}\label{eq:pre-hazard-sum}
\sum_{m \in \mathcal J} \E[\zeta_m \kappa(G_m)] \le \Pp(H - 1 \in \mathcal J) \le 1.
\end{equation}
\end{lemma}

\begin{proof}
We have
\[
\{H = m + 1\} = \{G_m \text{ is colorable and } G_m + e_{m + 1} \text{ is not colorable}\},
\]
where $G_m$ is a function of $e_1, \dots, e_m$, and $e_{m + 1}$ is a uniformly random pair of $[n]$ independent of $\cF_m$. Using \cref{lem:pre-frozen}(a) with $G = G_m$ and $f = e_{m + 1}$, we obtain \eqref{eq:pre-hazard}. Multiplying \eqref{eq:pre-hazard} by $\zeta$ and taking expectations gives the second statement. To prove \eqref{eq:pre-hazard-sum}, we use the second statement for each $m \in \mathcal J$ and the bound $\zeta_m \le 1$, and we find
\[
\sum_{m \in \mathcal J} \E[\zeta_m \kappa(G_m)] = \sum_{m \in \mathcal J} \E[\zeta_m \ind\{H = m + 1\}] \le \sum_{m \in \mathcal J} \Pp(H = m + 1) = \Pp(H - 1 \in \mathcal J),
\]
where the last equality holds because the events $\{H = m + 1\}$ are disjoint.
\end{proof}

To isolate final edges with prescribed properties, we use sets of frozen pairs. Let $\cE$ be a set whose elements are of the form $(G, e)$, where $G$ is a colorable multigraph on $[n]$ and $e$ is a frozen pair of $G$. For a multigraph $G$ on $[n]$, we define
\begin{equation}\label{eq:pre-kappaE}
\kappa_{\cE}(G) = \binom{n}{2}^{-1} \#\{e : (G, e) \in \cE\}.
\end{equation}
Then $0 \le \kappa_{\cE}(G) \le \kappa(G)$, and $\kappa_{\cE}(G) = 0$ if $G$ is not colorable. For a permutation $\phi$ of $[n]$ and a pair $e = \{u, v\}$, we write $\phi(e) = \{\phi(u), \phi(v)\}$. We say that $\cE$ is \emph{invariant under relabeling} if, for every permutation $\phi$ of $[n]$, we have $(\phi G, \phi(e)) \in \cE$ if and only if $(G, e) \in \cE$. We call a function $h$ on multigraphs on $[n]$ \emph{invariant under relabeling} if $h(\phi G) = h(G)$ for every multigraph $G$ and permutation $\phi$ of $[n]$. We call the law of a random multigraph $G$ on $[n]$ \emph{invariant under relabeling} if $\phi G$ has the law of $G$ for every permutation $\phi$ of $[n]$. For example, the law of $G_m$ is invariant under relabeling for every $m \ge 0$.

\begin{lemma}\label{lem:pre-terminal}
Let $\cE$ be a set whose elements are of the form $(G, e)$, where $G$ is a colorable multigraph on $[n]$ and $e$ is a frozen pair of $G$.
\begin{enumerate}[label=(\alph*)]
\item Let $h_0, h_1, \dots$ be functions with values in $[0, \infty)$ on the set of multigraphs on $[n]$, let $L' \ge 0$ be an integer, and let $U'$ be an integer with $U' > L'$ or let $U' = \infty$. When $U' = \infty$, we define $\{L' < H \le U'\}$ to mean $\{L' < H < \infty\}$. Then
\begin{equation}\label{eq:pre-terminal}
\E\big[h_{H - 1}(\Gamma) \ind\{L' < H \le U', \, (\Gamma, e_H) \in \cE\}\big] = \sum_{m = L'}^{U' - 1} \E\big[\kappa_{\cE}(G_m) h_m(G_m)\big].
\end{equation}
If $U' = \infty$, the sum runs over all $m \ge L'$. In particular, \eqref{eq:pre-terminal} holds with $h_m = h$ for all $m$, for every function $h \ge 0$ on multigraphs on $[n]$.
\item Suppose that $\cE$ and a function $h \ge 0$ on multigraphs on $[n]$ are invariant under relabeling, and let $G$ be a random multigraph on $[n]$ whose law is invariant under relabeling. Then for all distinct $a, b \in [n]$, we have
\begin{equation}\label{eq:pre-fixed-pair}
\E\big[\kappa_{\cE}(G) h(G)\big] = \E\big[\ind\{(G, \{a, b\}) \in \cE\} h(G)\big].
\end{equation}
\end{enumerate}
\end{lemma}

\begin{proof}
We begin with (a). Fix $m \ge 0$, and suppose that $(G_m, e_{m + 1}) \in \cE$. Then $G_m$ is colorable and $e_{m + 1}$ is a frozen pair of $G_m$, so that $G_m + e_{m + 1}$ is not colorable by \cref{lem:pre-frozen}(a), and $H = m + 1$. Since $\Gamma = G_m$ and $e_H = e_{m + 1}$ on the event $\{H = m + 1\}$, it follows that
\begin{equation}\label{eq:pre-terminal-event}
\{H = m + 1, \, (\Gamma, e_H) \in \cE\} = \{(G_m, e_{m + 1}) \in \cE\}.
\end{equation}

The pair $e_{m + 1}$ is uniformly distributed on the pairs of $[n]$ and independent of $\cF_m$. Therefore the conditional probability of the event \eqref{eq:pre-terminal-event} given $\cF_m$ is $\kappa_{\cE}(G_m)$. Since $h_m(G_m)$ is nonnegative and $\cF_m$-measurable, we obtain
\[
\begin{aligned}
\E\big[h_{H - 1}(\Gamma) \ind\{H = m + 1, \, (\Gamma, e_H) \in \cE\}\big] &= \E\big[h_m(G_m) \ind\{(G_m, e_{m + 1}) \in \cE\}\big]\\
&= \E\big[\kappa_{\cE}(G_m) h_m(G_m)\big].
\end{aligned}
\]
The event $\{L' < H \le U'\}$ is the disjoint union of the events $\{H = m + 1\}$ over $L' \le m \le U' - 1$. Summing the previous display over these $m$, we obtain \eqref{eq:pre-terminal}. (All terms are nonnegative, and the exchange of sum and expectation is justified even when the sum is infinite.)

We now prove (b). By the definition \eqref{eq:pre-kappaE}, we have
\[
\kappa_{\cE}(G) h(G) = \binom{n}{2}^{-1} \sum_{e} \ind\{(G, e) \in \cE\} h(G),
\]
where the sum runs over the pairs $e$ of $[n]$. For each pair $e$, we choose a permutation $\phi$ of $[n]$ with $\phi(e) = \{a, b\}$. The invariance of $\cE$ and $h$ gives
\[
\ind\{(G, e) \in \cE\} h(G) = \ind\{(\phi G, \{a, b\}) \in \cE\} h(\phi G),
\]
and since $\phi G$ has the law of $G$, the expectation of the right side is the right side of \eqref{eq:pre-fixed-pair}. Averaging over the $\binom{n}{2}$ pairs $e$ completes the proof.
\end{proof}

\Cref{lem:pre-terminal}(a) converts statements about the final edge into statements about frozen pairs at deterministic times. We use this identity to estimate probabilities involving the final edge.

\section{Lifetime}\label{sec:pre-lifetime}

We will control the window by the expected number of further uniformly random pairs that the last colorable graph $\Gamma$ can accept. We now define this quantity for a general multigraph.

\begin{definition}\label{def:pre-lifetime}
Let $S$ be a finite set with $N \ge 2$ elements, let $G$ be a multigraph on $S$, and let $f_1, f_2, \dots$ be an edge process on $S$. We define
\begin{equation}\label{eq:pre-T}
T_G = \min\{k \ge 0 : G + f_1 + \dots + f_k \text{ is not colorable}\} \in \{0, 1, \dots\} \cup \{\infty\}.
\end{equation}
We call $r(G) = \E T_G \in [0, \infty]$ the \emph{expected lifetime} of $G$, and for integers $b \ge 0$, we set $r_b(G) = \E \min(T_G, b)$.
\end{definition}

In words, $r(G)$ is the expected number of independent uniformly random pairs that must be added to $G$ to make it non-colorable. The law of $T_G$ is the same for every edge process on $S$, and therefore $r(G)$ and $r_b(G)$ are functions of $G$. Since $\min(T_G, b) = \sum_{k = 0}^{b - 1} \ind\{T_G > k\}$, we have
\begin{equation}\label{eq:pre-rb}
r_b(G) = \sum_{k = 0}^{b - 1} \Pp(T_G > k).
\end{equation}

\begin{lemma}\label{lem:pre-lifetime}
Let $S$ be a finite set with $N \ge 2$ elements, and let $G$ and $G'$ be multigraphs on $S$.
\begin{enumerate}[label=(\alph*)]
\item If $G$ is not colorable, then $r(G) = 0$ and $r_b(G) = 0$ for all $b \ge 0$. If $G$ is colorable, then $1 \le r_b(G) \le \min(b, r(G))$ for all integers $b \ge 1$, and $r_b(G)$ is nondecreasing in $b$. If $G$ is colorable and $N \ge 2q$, then
\begin{equation}\label{eq:pre-Aq}
r(G) \le A_q N.
\end{equation}
\item If $G \subset G'$, then $r(G') \le r(G)$ and $r_b(G') \le r_b(G)$ for all $b \ge 0$. If $\phi \colon S \to S'$ is a bijection, then $r(\phi G) = r(G)$ and $r_b(\phi G) = r_b(G)$ for all $b \ge 0$.
\item If $\kappa(G) > 0$, then $r(G) \le 1/\kappa(G)$.
\item For the edge process on $[n]$ and all integers $m \ge 0$, we have
\begin{equation}\label{eq:pre-r-conditional}
r(G_m) = \E\big[(H - m) \ind\{H > m\} \mid \cF_m\big] \qquad \text{almost surely}.
\end{equation}
In particular, $\E H = r(G_0) \le A_q n$.
\end{enumerate}
\end{lemma}

\begin{proof}
We begin with (a). We have $T_G = 0$ if $G$ is not colorable, and $T_G \ge 1$ if $G$ is colorable, so that $\min(T_G, b) \ge 1$ for $b \ge 1$ in the latter case. The bounds $\min(T_G, b) \le b$ and $\min(T_G, b) \le T_G$, and the monotonicity in $b$, hold pointwise.

We now prove \eqref{eq:pre-Aq}. Let $\emptyset_S$ denote the empty multigraph on $S$, and compute $T_G$ and $T_{\emptyset_S}$ from the same edge process. If $G + f_1 + \dots + f_k$ is colorable, then so is $f_1 + \dots + f_k$. This implies that $T_G \le T_{\emptyset_S}$, and it suffices to bound $r(\emptyset_S)$.

Let $\sigma\colon S \to [q]$ be a map. For $i \in [q]$, we define $b_i = |\sigma^{-1}(\{i\})|$. Using the convexity of $x \mapsto x(x - 1)/2$ and Jensen's inequality, we have
\[
\#\{\text{pairs } \{u, v\} \text{ of } S : \sigma(u) = \sigma(v)\} = \sum_{i = 1}^q \frac{b_i (b_i - 1)}{2} \ge \frac{N (N/q - 1)}{2}.
\]
The fraction of pairs of $S$ on which $\sigma$ is constant is at least $(N/q - 1)/(N - 1)$, which is at least $1/(2q)$ when $N \ge 2q$.

It follows that $\sigma$ is a coloring of $f_1 + \dots + f_k$ with probability at most $(1 - 1/(2q))^k \le e^{-k/(2q)}$. Taking a union bound over the $q^N$ maps $\sigma$, we obtain
\[
\Pp(T_{\emptyset_S} > k) = \Pp(f_1 + \dots + f_k \text{ is colorable}) \le \min\big(1, q^N e^{-k/(2q)}\big).
\]
Let $k_0 = \lceil 2 q N \log q \rceil$. For $k \ge k_0$, we have $q^N e^{-k/(2q)} \le e^{-(k - k_0)/(2q)}$, and therefore
\[
\begin{aligned}
r(\emptyset_S) = \sum_{k \ge 0} \Pp(T_{\emptyset_S} > k) &\le k_0 + \sum_{j \ge 0} e^{-j/(2q)} = k_0 + \frac{1}{1 - e^{-1/(2q)}}\\
&\le 2 q N \log q + 2q + 2 \le A_q N.
\end{aligned}
\]
In the second line, we used $1/(1 - e^{-x}) \le 1 + 1/x$ for $x > 0$ and $2q + 2 \le (2q + 1) N$, which holds since $N \ge 2$.

For (b), we compute $T_G$ and $T_{G'}$ from the same edge process. If $G' + f_1 + \dots + f_k$ is colorable, then so is $G + f_1 + \dots + f_k$. We conclude that $T_{G'} \le T_G$ pointwise, which gives the inequalities.

Now let $\phi \colon S \to S'$ be a bijection. Then $(\phi(f_k))_{k \ge 1}$ is an edge process on $S'$, and $\sigma \mapsto \sigma \circ \phi^{-1}$ is a bijection from $\Col(G + f_1 + \dots + f_k)$ to $\Col(\phi G + \phi(f_1) + \dots + \phi(f_k))$. The random variable $T_{\phi G}$, computed from $(\phi(f_k))$, coincides with $T_G$, which gives the invariance.

Next, we prove (c). Let $\mathcal Y$ denote the set of frozen pairs of $G$. Then $|\mathcal Y| = \kappa(G)\binom{N}{2} > 0$. We set
\[
T' = \min\{k \ge 1 : f_k \in \mathcal Y\}.
\]
Then $T'$ has the geometric law with parameter $\kappa(G)$ on $\{1, 2, \dots\}$, and $\E T' = 1/\kappa(G)$. It remains to show that $T_G \le T'$.

This holds if $G + f_1 + \dots + f_{T' - 1}$ is not colorable. Otherwise, every coloring of $G + f_1 + \dots + f_{T' - 1}$ is a coloring of $G$ and is constant on the frozen pair $f_{T'}$, so $G + f_1 + \dots + f_{T'}$ is not colorable.

Finally, we prove (d). For a multigraph $G$ on $[n]$ and a sequence $(p_k)_{k \ge 1}$ of pairs of $[n]$, let
\[
\Theta\big(G, (p_k)_{k \ge 1}\big) = \min\{k \ge 0 : G + p_1 + \dots + p_k \text{ is not colorable}\}.
\]
By the definition of $T_G$, we have $\E \Theta(G, (f_k)_{k \ge 1}) = r(G)$ for every multigraph $G$ on $[n]$ and every edge process $(f_k)$ on $[n]$.

If $H > m$, then $H - m = \Theta(G_m, (e_{m + k})_{k \ge 1})$. If $H \le m$, then $G_m$ is not colorable and $\Theta(G_m, (e_{m + k})_{k \ge 1}) = 0$. Hence
\[
(H - m) \ind\{H > m\} = \Theta\big(G_m, (e_{m + k})_{k \ge 1}\big).
\]
The sequence $(e_{m + k})_{k \ge 1}$ is an edge process on $[n]$ independent of $\cF_m$, and $G_m$ is a function of $e_1, \dots, e_m$. Since $\Theta \ge 0$, the conditional expectation of the right side given $\cF_m$ is $r(G_m)$, which is \eqref{eq:pre-r-conditional}. Taking $m = 0$ and expectations, and using (a) with $N = n \ge 2q$, we obtain $\E H = r(G_0) \le A_q n$.
\end{proof}

The identity \eqref{eq:pre-r-conditional} holds at deterministic times $m$. Since $H - 1$ is not a stopping time, we always pass from $\Gamma$ to the graphs $G_m$ through \cref{lem:pre-terminal} when we bound $r(\Gamma)$.

\begin{lemma}\label{lem:pre-classmass}
Let $S$ be a finite set with $N \ge 2q$ elements, let $G$ be a colorable multigraph on $S$, and let $y > 0$. If $r(G) > y$, then
\begin{equation}\label{eq:pre-classmass}
Q(G) \le \frac{1 + A_q}{y}, \qquad |C_G(u)| \le N \Big(\frac{1 + A_q}{y}\Big)^{1/2} \quad \text{for all } u \in S.
\end{equation}
\end{lemma}

\begin{proof}
By \cref{lem:pre-lifetime}(c), we have $\kappa(G) r(G) \le 1$ (this is trivial if $\kappa(G) = 0$). \Cref{lem:pre-lifetime}(a) gives $r(G)/N \le A_q$. Using \cref{lem:pre-frozen}(c), we obtain
\[
Q(G) r(G) = \Big(1 - \frac 1N\Big) \kappa(G) r(G) + \frac{r(G)}{N} \le 1 + A_q.
\]
Since $r(G) > y$, this gives the first bound in \eqref{eq:pre-classmass}, and the second follows from the bound $|C_G(u)| \le N Q(G)^{1/2}$ of \cref{lem:pre-frozen}(c).
\end{proof}

\section{Poissonization and splitting}\label{sec:pre-poisson}

When the number of sampled pairs is Poisson, splitting the possible pairs into disjoint sets gives independent sublists. We record the facts that we use.

\begin{definition}\label{def:pre-poissonized}
Let $S$ be a finite set with at least two elements, let $f_1, f_2, \dots$ be an edge process on $S$, and let $(\Pi_t)_{t \ge 0}$ be a rate-one Poisson counting process independent of $(f_i)$. We define a \emph{list} to be a finite sequence. We call the multigraph-valued process $t \mapsto f_1 + \dots + f_{\Pi_t}$ the \emph{Poissonized process on $S$}. We call its value at time $t$ the Poissonized process on $S$ at time $t$. For $a \in S$, $S' \subset S \setminus \{a\}$, and $t \ge 0$, we define the \emph{neighbor list of $a$ in $S'$ at time $t$} by including one entry $x$ for every index $i \le \Pi_t$ with $f_i = \{a, x\}$ and $x \in S'$, in increasing order of $i$.
\end{definition}

A neighbor list may have repeated entries. For $\mu \ge 0$ and a finite set $E'$ with a probability measure $\nu'$, we call a random list a \emph{list of $\Pois(\mu)$ independent elements of $E'$ with law $\nu'$} if its length has the $\Pois(\mu)$ law and, conditionally on its length, its entries are independent with law $\nu'$. Throughout the paper, we let $(\Pi_t)_{t \ge 0}$ denote a rate-one Poisson counting process independent of $(e_i)$, and we set $G(t) = G_{\Pi_t}$ for $t \ge 0$. Thus $(G(t))_{t \ge 0}$ is the Poissonized process on $[n]$, and the law of $G(t)$ is invariant under relabeling for every $t \ge 0$. For the rest of this section, let $S$ be a finite set with at least two elements. We let $(G^S_m)_{m \ge 0}$ denote the time-$m$ graphs of an edge process on $S$, and $(G^S(t))_{t \ge 0}$ the Poissonized process on $S$ built from it.

\begin{lemma}\label{lem:pre-poisson}
Let $S$ be a finite set with at least two elements, and let $h \ge 0$ be a function on multigraphs on $S$. For every real $\theta > 0$, we have
\begin{equation}\label{eq:pre-poisson}
\int_0^\infty \E h\big(G^S(\theta t)\big) \dd t = \frac{1}{\theta} \sum_{m \ge 0} \E h(G^S_m),
\end{equation}
where both sides may be infinite.
\end{lemma}

\begin{proof}
For $s \ge 0$, the random variable $\Pi_s$ is independent of the edge process and has the $\Pois(s)$ law, so that
\[
\E h\big(G^S(s)\big) = \sum_{m \ge 0} e^{-s} \frac{s^m}{m!} \E h(G^S_m).
\]
We insert this identity with $s = \theta t$ into the left side of \eqref{eq:pre-poisson}, and we exchange sum and integral, which is allowed by Tonelli's theorem since all terms are nonnegative. Using that
\[
\int_0^\infty e^{-\theta t} \frac{(\theta t)^m}{m!} \dd t = \frac{1}{\theta}
\]
for every integer $m \ge 0$, we obtain \eqref{eq:pre-poisson}.
\end{proof}

\begin{lemma}\label{lem:pre-split}
\begin{enumerate}[label=(\alph*)]
\item Let $E$ be a finite set, let $\nu$ be a probability measure on $E$, let $X_1, X_2, \dots$ be independent random elements of $E$ with law $\nu$, and let $I$ be an independent random variable with the $\Pois(\mu)$ law, where $\mu \ge 0$. Let $E_1, \dots, E_{J}$ be a partition of $E$ into sets with $\nu(E_j) > 0$. Then the subsequences of $(X_1, \dots, X_I)$ formed by the terms in $E_1, \dots, E_{J}$ are independent, and for each $j$, the subsequence formed by the terms in $E_j$ is a list of $\Pois(\mu \nu(E_j))$ independent elements of $E_j$ with law $\nu(\, \cdot \mid E_j)$.
\item Let $S$ be a finite set with $N \ge 3$ elements, let $S_0 \subset S$, and set $S' = S \setminus S_0$. Suppose that $S'$ has at least two elements, and let $t \ge 0$. Then the multigraphs $G^S(t)[S']$ and $G^S(t)[S_0]$ and the neighbor lists of the vertices $a \in S_0$ in $S'$ at time $t$ are independent. Further, $G^S(t)[S']$ has the law of the Poissonized process on $S'$ at time $t \binom{|S'|}{2}/\binom{N}{2}$, and each of these neighbor lists is a list of $\Pois(t |S'|/\binom{N}{2})$ independent uniformly random elements of $S'$.
\item Let $S$ be a finite set with $N \ge 3$ elements, let $a \in S$, and let $t \ge 0$. Then $\deg_{G^S(t)}(a)$ is the length of the neighbor list of $a$ in $S \setminus \{a\}$ at time $t$, and $G^S(t)$ is obtained from $G^S(t)[S \setminus \{a\}]$ by adding one occurrence $\{a, x\}$ for each entry $x$ of this list. The list is independent of $G^S(t)[S \setminus \{a\}]$, it is a list of $\Pois(2t/N)$ independent uniformly random elements of $S \setminus \{a\}$, and $G^S(t)[S \setminus \{a\}]$ has the law of the Poissonized process on $S \setminus \{a\}$ at time $(N - 2) t/N$.
\end{enumerate}
\end{lemma}

\begin{proof}
We begin with (a). Fix integers $k_1, \dots, k_{J} \ge 0$, let $k = k_1 + \dots + k_{J}$, and fix sequences
\[
y^{(j)} = (y^{(j)}_1, \dots, y^{(j)}_{k_j}) \in E_j^{k_j}, \qquad 1 \le j \le J.
\]
The subsequences of $(X_1, \dots, X_I)$ formed by the terms in $E_1, \dots, E_{J}$ equal $y^{(1)}, \dots, y^{(J)}$ if and only if $I = k$ and $(X_1, \dots, X_k)$ is one of the $k!/(k_1! \cdots k_{J}!)$ sequences obtained by interleaving $y^{(1)}, \dots, y^{(J)}$. Each such sequence has probability $\prod_{j = 1}^{J} \prod_{i = 1}^{k_j} \nu(y^{(j)}_i)$, so the probability of this event is
\[
e^{-\mu} \frac{\mu^k}{k!} \cdot \frac{k!}{k_1! \cdots k_{J}!} \prod_{j = 1}^{J} \prod_{i = 1}^{k_j} \nu(y^{(j)}_i) = \prod_{j = 1}^{J} \bigg( e^{-\mu \nu(E_j)} \frac{(\mu \nu(E_j))^{k_j}}{k_j!} \prod_{i = 1}^{k_j} \frac{\nu(y^{(j)}_i)}{\nu(E_j)} \bigg),
\]
where we used that $\nu(E_1) + \dots + \nu(E_{J}) = 1$. The $j$-th factor on the right is the probability that a list of $\Pois(\mu \nu(E_j))$ independent elements of $E_j$ with law $\nu(\, \cdot \mid E_j)$ coincides with $y^{(j)}$. Since all lists take values in countable sets, (a) follows.

For (b), let $(f_i)_{i \ge 1}$ denote the edge process underlying $G^S(t)$. We apply (a) with $E$ the set of pairs of $S$, $\nu$ the uniform law, $X_i = f_i$, $I = \Pi_t$, and $\mu = t$. We partition $E$ into the pairs of $S'$, the pairs of $S_0$ (if $|S_0| \ge 2$), and, for each $a \in S_0$, the $|S'|$ pairs $\{a, x\}$ with $x \in S'$. The multigraph $G^S(t)[S']$ is formed by the subsequence of the pairs of $S'$, which is a list of $\Pois(t \binom{|S'|}{2}/\binom{N}{2})$ independent uniformly random pairs of $S'$. Let $\Pi'$ be an independent rate-one Poisson counting process, and set $s = t\binom{|S'|}{2}/\binom N2$. Then this is the law of the multigraph formed by the first $\Pi'_s$ terms of an edge process on $S'$.

The multigraph $G^S(t)[S_0]$ is a function of the subsequence of the pairs of $S_0$. The neighbor list of $a \in S_0$ in $S'$ is obtained from the subsequence of the pairs $\{a, x\}$ with $x \in S'$ by replacing each pair $\{a, x\}$ by $x$, and the uniform law on these pairs corresponds to the uniform law on $S'$. By (a), these subsequences are independent and have the stated laws, which proves (b).

To prove (c), we take $S_0 = \{a\}$ in (b), so that $|S'| = N - 1$, and we note that $\binom{N - 1}{2}/\binom{N}{2} = (N - 2)/N$ and $(N - 1)/\binom{N}{2} = 2/N$. The occurrences of $G^S(t)$ that contain $a$ are the pairs $\{a, x\}$ with $x$ an entry of the neighbor list of $a$ in $S \setminus \{a\}$, which gives the first two statements of (c).
\end{proof}

\Cref{lem:pre-split} is a special case of the restriction and marking theorems for Poisson processes; see \cite[Theorems~5.2, 5.6, and~5.8, Corollary~5.9]{last2018lectures}. We use it only at fixed times, and we have included the short computation for this case.

\section{Tail bounds}\label{sec:pre-tails}

We use the convention $0^0 = 1$. For $z \ge 0$ and integers $k \ge 0$, we define $\pi_z(k) = e^{-z} z^k/k!$. We also set $\pi_z(-1) = 0$. For $k \ge 0$, the value $\pi_z(k)$ is the probability that a $\Pois(z)$ random variable takes the value $k$. Inside a probability expression, we let $\Pois(z)$ denote a random variable with the Poisson law of mean $z$. We use the analogous convention for $\Bin(m, p)$. We will use the following elementary tail bounds.

\begin{lemma}\label{lem:pre-tails}
\begin{enumerate}[label=(\alph*)]
\item Let $X$ have the $\Bin(m, p)$ law, where $m \ge 1$ is an integer and $p \in [0, 1]$. For all integers $k \ge 1$, we have
\[
\Pp(X \ge k) \le \binom{m}{k} p^k \le \Big(\frac{e m p}{k}\Big)^k.
\]
\item Let $X$ have the $\Pois(\mu)$ law, where $\mu \ge 0$. For all real $k > \mu$, we have $\Pp(X \ge k) \le (e \mu/k)^k$.
\item Let $X$ have the $\Pois(\mu)$ law, where $\mu > 0$. For all real $t > 0$, we have
\[
\Pp(X \ge \mu + t) \le \exp\Big(-\frac{t^2}{2(\mu + t)}\Big), \qquad \Pp(X \le \mu - t) \le \exp\Big(-\frac{t^2}{2\mu}\Big).
\]
\item For all real $z \ge 0$ and integers $i \ge 1$, we have
\[
\Pp\big(\Pois(z) \ge i\big) \le (e^z - 1) \pi_z(i - 1).
\]
\end{enumerate}
\end{lemma}

\begin{proof}
We begin with (a), where we may assume $k \le m$, since otherwise the probability is zero. If $X \ge k$, then some set of $k$ of the $m$ trials consists of successes, and the union bound over these sets gives the first inequality. The second follows from $\binom{m}{k} \le m^k/k!$ and $k! \ge (k/e)^k$.

To prove (b) and (c), we use the Cram\'er--Chernoff method \cite[Section~2.2]{boucheron2013concentration}. If $\mu = 0$, then $X = 0$ and (b) holds. We now suppose that $\mu > 0$. We have $\E e^{\varphi X} = \exp(\mu (e^\varphi - 1))$ for every real $\varphi$. Applying Markov's inequality, we find that for $\varphi \ge 0$ and every real $x$,
\begin{equation}\label{eq:pre-chernoff}
\Pp(X \ge x) \le \exp\big(-\varphi x + \mu (e^\varphi - 1)\big).
\end{equation}
For (b), choosing $e^\varphi = k/\mu$ gives
\[
\Pp(X \ge k) \le e^{-\mu} \Big(\frac{e\mu}{k}\Big)^k \le \Big(\frac{e\mu}{k}\Big)^k.
\]

For the first bound in (c), we set $x = t/\mu$ and
\[
\psi(v) = (1 + v) \log(1 + v) - v, \qquad v > -1.
\]
Choosing $e^\varphi = 1 + x$ in \eqref{eq:pre-chernoff}, we obtain
\[
\Pp(X \ge \mu + t) \le \exp\big(-\varphi (\mu + t) + \mu (e^\varphi - 1)\big) = \exp\big(-\mu \psi(x)\big).
\]
Since $\psi(0) = 0$, $\psi'(0) = 0$, and $\psi''(v) = 1/(1 + v)$, we have
\[
\psi(x) = \int_0^x \frac{x - v}{1 + v} \dd v \ge \frac{1}{1 + x} \int_0^x (x - v) \dd v = \frac{x^2}{2(1 + x)},
\]
and therefore $\mu \psi(x) \ge t^2/(2(\mu + t))$.

For the second bound in (c), observe that the event is empty if $t > \mu$, and that its probability is $e^{-\mu} \le e^{-\mu/2}$ if $t = \mu$. Suppose that $0 < t < \mu$, and set $x = t/\mu \in (0, 1)$. For $\varphi \le 0$, Markov's inequality applied to $e^{\varphi X}$ gives
\[
\Pp(X \le \mu - t) \le \exp\big(-\varphi (\mu - t) + \mu (e^\varphi - 1)\big),
\]
and the choice $e^\varphi = 1 - x$ yields $\Pp(X \le \mu - t) \le \exp(-\mu \psi(-x))$, by the same computation. As before, we have
\[
\psi(-x) = \int_0^x \frac{x - v}{1 - v} \dd v \ge \int_0^x (x - v) \dd v = \frac{x^2}{2},
\]
so that $\mu \psi(-x) \ge t^2/(2\mu)$.

It remains to prove (d). Both sides vanish if $z = 0$. If $z > 0$, then
\[
\frac{\Pp(\Pois(z) \ge i)}{\pi_z(i - 1)} = \sum_{j \ge 1} \frac{z^j (i - 1)!}{(i - 1 + j)!} \le \sum_{j \ge 1} \frac{z^j}{j!} = e^z - 1,
\]
where we used that
\[
\frac{(i - 1 + j)!}{(i - 1)!} = \prod_{l = 1}^{j} (i - 1 + l) \ge j!. \qedhere
\]
\end{proof}

\section{Bounds on the quantiles}\label{sec:pre-centre}

We now bound the quantiles $L_n(\eta)$ and $U_n(\eta)$ from above and below. The upper bound is the first-moment bound. The lower bound rests on the classical fact that, for every fixed $c_0 \in (0, 1/2)$, with probability tending to one every component of $G(n, \lfloor c_0 n \rfloor)$ is a tree or has exactly one cycle \cite[Theorem~5e]{erdos1960evolution}. A quantitative form, proved by counting witness subgraphs, appears in \cite[Theorem~5.5 and its proof, pp.~112--113]{janson2000random}. We use the same witness types below, with a count adapted to our process, which samples pairs with replacement.

We will need the following lemma on cycles in multigraphs. For an integer $j \ge 0$, we define a \emph{path of length $j$} to consist of distinct vertices $v_0, \dots, v_j$ and distinct occurrences $p_1, \dots, p_j$ with $p_i = \{v_{i - 1}, v_i\}$. For an integer $j \ge 2$, we define a \emph{cycle of length $j$} to consist of distinct vertices $v_1, \dots, v_j$ and distinct occurrences $p_1, \dots, p_j$ with $p_i = \{v_i, v_{i + 1}\}$ for $i < j$ and $p_j = \{v_j, v_1\}$. (A cycle of length $2$ consists of two occurrences of the same pair.) We call a multigraph a \emph{theta} if it consists of three paths of lengths at least $1$ between the same two distinct end vertices, with no other vertices or occurrences in common. We call those two end vertices its \emph{branch vertices}. We call a multigraph a \emph{dumbbell} if it consists of two cycles without common occurrences, which either have exactly one common vertex, or have no common vertex and are joined by a path of length at least $1$ whose only vertices on the two cycles are its end vertices, one on each cycle. A theta or a dumbbell with $k$ vertices has $k + 1$ occurrences. For a multigraph $F$, we let $v(F)$ denote the number of vertices that lie in some occurrence of $F$, and $|F|$ the number of occurrences of $F$.

\begin{lemma}\label{lem:pre-cycles}
Let $F$ be a connected multigraph with at least two vertices.
\begin{enumerate}[label=(\alph*)]
\item If $|F| \le v(F)$, then $F$ is colorable.
\item If $|F| \ge v(F) + 1$, then $F$ contains a theta or a dumbbell.
\end{enumerate}
\end{lemma}

\begin{proof}
Let $\mathcal T$ be a spanning tree of $F$. It consists of $v(F) - 1$ occurrences of $F$, no two of them occurrences of the same pair. For an occurrence $p = \{x, y\}$ of $F$ outside $\mathcal T$, the occurrence $p$ and the path in $\mathcal T$ from $x$ to $y$ form a cycle of length at least $2$, and $p$ is the only occurrence of this cycle outside $\mathcal T$.

For (a), we use the assumption $q \ge 3$. If $|F| = v(F) - 1$, then $F = \mathcal T$ is a tree, which has a coloring with two colors. Otherwise, $F$ has exactly one occurrence $p$ outside $\mathcal T$, and we let $\mathcal C$ be the cycle formed by $p$ as above. The cycle $\mathcal C$ has a coloring with three colors.

We now extend this coloring to $F$. Since $\mathcal T$ is a tree and $p$ is the only occurrence of $F$ outside $\mathcal T$, the only occurrences of $F$ between vertices of $\mathcal C$ are those of $\mathcal C$. We color the remaining vertices in order of increasing distance in $\mathcal T$ from the vertex set of $\mathcal C$. Each vertex $y$ outside $\mathcal C$ is joined in $F$ to exactly one vertex that is colored before it, namely its neighbor on the path in $\mathcal T$ from $y$ to $\mathcal C$, since its other neighbors in $\mathcal T$ are farther from $\mathcal C$. We give $y$ a color different from the color of that neighbor, and we obtain a coloring of $F$.

We now prove (b). Let $p$ and $p'$ be two distinct occurrences of $F$ outside $\mathcal T$, and let $\mathcal C$ and $\mathcal C'$ be the cycles that they form. The only occurrence of $\mathcal C$ outside $\mathcal T$ is $p$, so $p'$ does not belong to $\mathcal C$.

Suppose first that $\mathcal C$ and $\mathcal C'$ have at most one common vertex. Then they have no common occurrence. If they have one common vertex, they form a dumbbell. If they have none, we let $\mathcal P$ be a shortest path in $F$ from a vertex of $\mathcal C$ to a vertex of $\mathcal C'$. By minimality, the only vertices of $\mathcal P$ on $\mathcal C$ or $\mathcal C'$ are its end vertices, and $\mathcal C$, $\mathcal C'$, and $\mathcal P$ form a dumbbell.

Suppose now that $\mathcal C$ and $\mathcal C'$ have at least two common vertices. We write the vertices of $\mathcal C'$ in cyclic order as $v_1, \dots, v_j$, with $p' = \{v_j, v_1\}$, and we let $i_1$ and $i_2$ be the smallest and the largest index $i$ such that $v_i$ is a vertex of $\mathcal C$. Then $i_1 < i_2$, and the path $v_{i_2}, v_{i_2 + 1}, \dots, v_j, v_1, \dots, v_{i_1}$ along $\mathcal C'$ contains $p'$, joins the distinct vertices $v_{i_2}$ and $v_{i_1}$ of $\mathcal C$, and has no other vertex on $\mathcal C$. Its occurrences are not occurrences of $\mathcal C$.

Indeed, $p'$ is not, and if the path has length at least $2$, then each of its occurrences contains a vertex outside $\mathcal C$. Together with the two paths along $\mathcal C$ between $v_{i_1}$ and $v_{i_2}$, this path forms a theta.
\end{proof}

\begin{lemma}\label{lem:pre-centre}
Let $\eta \in (0, 1/2]$.
\begin{enumerate}[label=(\alph*)]
\item Let $c_1 > \gamma_q$ be real, and let $n \ge 2q$ be an integer. We set
\[
C_{c_1} = \frac{q}{q - 1} e^{2c_1}.
\]
Then
\begin{equation}\label{eq:pre-first-moment}
\Pp\big(H > \lfloor c_1 n \rfloor\big) \le C_{c_1} \exp\Big(-n (c_1 - \gamma_q) \log \frac{q}{q - 1}\Big).
\end{equation}
For every $c_1 > \gamma_q$, we have $L_n(\eta) \le U_n(\eta) \le c_1 n$ for all large $n$. Moreover, $U_n(\eta) < M$ for all large $n$.
\item For every $c_0 \in (0, 1/2)$, there exists $C(c_0)$ such that $\Pp(H \le \lfloor c_0 n \rfloor) \le C(c_0)/n$ for all $n \ge 2q$. Further, $L_n(\eta) > c_0 n$ for all large $n$.
\end{enumerate}
\end{lemma}

\begin{proof}
We begin with (a). Let $\sigma \colon [n] \to [q]$ be a map. As in the proof of \cref{lem:pre-lifetime}(a), the fraction of pairs of $[n]$ on which $\sigma$ is constant is at least $(n/q - 1)/(n - 1)$. This implies that $\sigma$ is a coloring of the single occurrence $e_1$ with probability at most
\[
1 - \frac{n/q - 1}{n - 1} = \frac{q - 1}{q} \cdot \frac{n}{n - 1}.
\]
Since the pairs $e_i$ are independent, a union bound over the $q^n$ maps $\sigma$ gives, for all $m \ge 0$,
\begin{equation}\label{eq:pre-union}
\Pp(H > m) = \Pp(G_m \text{ is colorable}) \le q^n \Big(\frac{q - 1}{q}\Big)^m \Big(\frac{n}{n - 1}\Big)^m.
\end{equation}
We set $a = \log(q/(q - 1))$ and $m = \lfloor c_1 n\rfloor$. Then $\log q = \gamma_q a$ and $c_1 n - 1 < m \le c_1 n$. Taking logarithms in \eqref{eq:pre-union}, we find
\[
\begin{aligned}
\log \Pp(H > m) &\le n \gamma_q a - m a + \frac{m}{n - 1} \\
&\le n \gamma_q a - (c_1 n - 1) a + 2c_1 \\
&= -n (c_1 - \gamma_q) a + a + 2c_1,
\end{aligned}
\]
where we used $\log(n/(n - 1)) \le 1/(n - 1)$ and $m/(n - 1) \le 2c_1$ for $n \ge 2$. Since $e^{a} = q/(q - 1)$, this is \eqref{eq:pre-first-moment}.

To deduce the bounds on the quantiles, fix $c_1 > \gamma_q$. For every integer $m$ with $\lfloor c_1 n \rfloor < m \le M$, we have by \eqref{eq:pre-V-H} that
\[
\Pp(V \ge m) = \Pp(H > m) \le \Pp\big(H > \lfloor c_1 n \rfloor\big),
\]
and by \eqref{eq:pre-first-moment} the right side is smaller than $\eta$ for all large $n$. For such $n$, no integer $m > \lfloor c_1 n \rfloor$ belongs to the set that defines $U_n(\eta)$. Therefore $U_n(\eta) \le c_1 n$, and $L_n(\eta) \le U_n(\eta)$ by \cref{lem:pre-quantile}. To prove that $U_n(\eta) < M$, we set $c_1' = (\gamma_q + \Lambda)/2$. Then $c_1' > \gamma_q$ and $c_1' n < \lfloor\Lambda n\rfloor = M$ for all large $n$. We apply the preceding bound with this value of $c_1'$.

We now prove (b). Suppose that $G_m$ is not colorable. Then some connected component $F$ of $G_m$ is not colorable, and \cref{lem:pre-cycles}(a) implies that $|F| \ge v(F) + 1$. By \cref{lem:pre-cycles}(b), it follows that $G_m$ contains a theta or a dumbbell. We now count these multigraphs.

Let $k \ge 2$. A theta with $k$ vertices is determined by the lengths $j_1, j_2, j_3 \ge 1$ of its three paths, which satisfy $j_1 + j_2 + j_3 = k + 1$, together with the list of its two branch vertices followed by the interior vertices of the three paths in order. A dumbbell with $k$ vertices is determined by the lengths $j_1, j_2 \ge 2$ of its cycles and the length $j_3 \ge 0$ of the joining path, which satisfy $j_1 + j_2 + j_3 = k + 1$, together with a list of its vertices.

We list the vertices of the first cycle in cyclic order, starting from its attachment vertex. If $j_3 = 0$, the two attachment vertices coincide, and we next list the other $j_2 - 1$ vertices of the second cycle in cyclic order from that shared vertex. If $j_3 \ge 1$, we next list all $j_2$ vertices of the second cycle in cyclic order, starting from its attachment vertex, and finally the $j_3 - 1$ interior vertices of the joining path in order. In the first case the list has $j_1 + j_2 - 1 = k$ entries, and in the second it has $j_1 + j_2 + j_3 - 1 = k$ entries.

In both cases, the list consists of $k$ distinct elements of $[n]$, and two of the three lengths, each in $\{0, \dots, k + 1\}$, determine the third. There are at most $(k + 2)^2$ choices of lengths and at most $n^k$ lists. The number of multigraphs on $[n]$ that are thetas or dumbbells with $k$ vertices is therefore at most $2 (k + 2)^2 n^k$.

Let $F$ be one of these multigraphs, and let $p_1, \dots, p_{k + 1}$ denote its occurrences. If $F \subset G_m$, then there exist distinct indices $i_1, \dots, i_{k + 1} \le m$ with $e_{i_j} = p_j$ for all $j$. There are at most $m^{k + 1}$ choices of these indices, and for each choice this event has probability $\binom{n}{2}^{-(k + 1)}$, so that
\[
\Pp(F \subset G_m) \le \bigg(m \binom{n}{2}^{-1}\bigg)^{k + 1}.
\]

Now let $m \le c_0 n$, so that $m \binom{n}{2}^{-1} \le 2 c_0/(n - 1)$. A union bound over $k$ and $F$ gives
\[
\Pp(G_m \text{ is not colorable}) \le \sum_{k = 2}^n 2 (k + 2)^2 n^k \Big(\frac{2 c_0}{n - 1}\Big)^{k + 1} \le \frac{2 e^2}{n - 1} \sum_{k \ge 2} (k + 2)^2 (2 c_0)^{k + 1},
\]
where we used that for $k \le n$,
\[
\Big(\frac{n}{n - 1}\Big)^k \le \Big(1 + \frac{1}{n - 1}\Big)^n \le e^{n/(n - 1)} \le e^2.
\]
The last series converges since $2 c_0 < 1$. Taking $m = \lfloor c_0n\rfloor$ and using $1/(n - 1) \le 2/n$ gives the required bound. We define the constant by
\[
C(c_0) = 4 e^2 \sum_{k \ge 2} (k + 2)^2 (2 c_0)^{k + 1}.
\]

Finally, we fix $c_0' \in (c_0, 1/2)$. For all large $n$, we have $c_0 n + 1 \le c_0' n < M$. For such $n$ and every integer $m$ with $0 \le m \le c_0 n$, we have $m < M$, hence $\{V \le m\} = \{H \le m + 1\}$, and therefore
\[
\Pp(V \le m) = \Pp(H \le m + 1) \le \Pp\big(H \le \lfloor c_0' n \rfloor\big) \le \frac{C(c_0')}{n}.
\]
The right side is smaller than $\eta$ for all large $n$. For such $n$, no integer $m \le c_0 n$ belongs to the set that defines $L_n(\eta)$, so $L_n(\eta) > c_0 n$.
\end{proof}

The assumption $q \ge 3$ is needed in (b), since for $q = 2$ a component with exactly one cycle is not colorable when the cycle has odd length.

\chapter{Reduction to two estimates}\label{sec:red}

In this section, we reduce the existence of the threshold to two estimates on the final edge $e_H$. The reduction has three steps. In \cref{sec:red-convergence}, we state our convergence criterion, \cref{thm:red-convergence}. It asserts that $\alpha_q$ exists if, at every small quantile level $\eta$, the window satisfies $w_n(\eta) = O(n \delta_n)$ for a sequence $\delta_n$ that is summable along $n = \lfloor e^k \rfloor$. We prove it in \cref{app:conv}. In \cref{sec:red-deletion}, we bound the variance of a capped version of $V$ by the expected lifetime of $\Gamma$, restricted to final edges inside the window. We conclude in \cref{prop:red-window} that $w_n(\eta) = O(n \eps_n)$ whenever this expected lifetime, capped at $n \eps_n$, is $O(n \eps_n^2)$. In \cref{sec:red-split}, we split the last expectation according to the degrees in $\Gamma$ of the endpoints of $e_H$, and we state the two estimates that remain. The heuristics behind these steps are explained in \cref{sec:overview}.

From \cref{sec:pre}, we use the quantiles of \cref{def:pre-quantile,lem:pre-quantile}, the frozen pairs of \cref{def:pre-frozen,lem:pre-frozen}, the identity \eqref{eq:pre-terminal}, and the expected lifetime of \cref{def:pre-lifetime}. We now fix the sequences that are used in the rest of the paper. We set $\eps_n = (\log n)^{-3/2}$ for $n \ge 3$ and $\eps_1 = \eps_2 = 1$, and we define
\begin{equation}\label{eq:red-parameters}
x_n = n \eps_n^2, \qquad D_n = \lfloor (\log n)^{1/4} \rfloor.
\end{equation}
We call $x_n$ the \emph{lifetime level} and $D_n$ the \emph{degree cut}.

\section{The convergence criterion}\label{sec:red-convergence}

Recall the level $\eta_0 = e^{-2 \Lambda}/8$ and the window $w_n(\eta) = U_n(\eta) - L_n(\eta)$ of \cref{def:pre-quantile}. The following theorem reduces the existence of the threshold to a bound on the window. Its proof, in \cref{app:conv}, uses the coloring interpolation of Bayati, Gamarnik, and Tetali \cite[Theorem~5 and Proposition~2]{bayati2013combinatorial} and adapts the comparisons between system sizes and the balanced-tree argument of the $k$-SAT preprint \cite[Section~5]{openai2026limiting} to quantiles and geometrically summable errors. The earlier additive form of near-superadditivity appears in \cite[arXiv:1006.3786v1, Section~4.3, Lemma~8 and Remark~3]{abbe2014concentration}.

\begin{theorem}\label{thm:red-convergence}
Let $(\delta_n)_{n \ge 1}$ be a nonincreasing sequence in $(0, 1]$ such that $\sum_{k \ge 1} \delta_{\lfloor e^k \rfloor} < \infty$. Suppose that for every $\eta \in (0, \eta_0]$ there exist $K'(\eta) > 0$ and an integer $n_1(\eta)$ such that
\begin{equation}\label{eq:red-window-hyp}
w_n(\eta) \le K'(\eta) n \delta_n \qquad \text{for all } n \ge n_1(\eta).
\end{equation}
Then there exists $\alpha_q \in [1/2, \gamma_q]$ with the following properties.
\begin{enumerate}[label=(\arabic*)]
\item For every $\eta \in (0, \eta_0]$, we have $L_n(\eta)/n \to \alpha_q$ and $U_n(\eta)/n \to \alpha_q$ as $n \to \infty$.
\item For every $c \in [0, \alpha_q)$, we have $\Pp(H > \lfloor cn \rfloor) \to 1$ and $\Pp(G(n, \lfloor cn \rfloor) \text{ is colorable}) \to 1$ as $n \to \infty$. For every $c > \alpha_q$, both probabilities tend to $0$.
\item For every $d \in [0, 2 \alpha_q)$, we have $\Pp(G(n, d/n) \text{ is colorable}) \to 1$ as $n \to \infty$. For every $d > 2 \alpha_q$, this probability tends to $0$.
\end{enumerate}
\end{theorem}

Note that the statements on $\Pp(H > \lfloor cn \rfloor)$ in (2) determine $\alpha_q$, so $\alpha_q$ depends only on $q$ (and not on $\Lambda$ or on the sequence $\delta_n$). We will apply \cref{thm:red-convergence} with $\delta_n = \eps_n$. The remaining argument verifies \eqref{eq:red-window-hyp} for this sequence.

\begin{remark}\label{rem:red-af}
The window bound at every level $\eta \in (0, \eta_0]$ makes the sharp-threshold theorem of Achlioptas and Friedgut \cite[Theorem~1.1]{achlioptas1999sharp} unnecessary for the proof of \cref{thm:red-convergence}. In \cref{app:conv}, we first prove convergence of the normalized quantiles at each level, use their nesting to identify a common limit, and then let $\eta$ tend to zero to obtain the zero-one laws. The summability of $(\delta_{\lfloor e^k \rfloor})$ is needed for the comparisons used there; see \cref{rem:conv-necessity}.
\end{remark}

\section{Deletion and the expected lifetime of the last colorable graph}\label{sec:red-deletion}

Throughout \cref{sec:red-deletion,sec:red-split}, we fix $\eta \in (0, \eta_0]$ and an integer $n \ge 2q$. We write $L$, $U$, and $w$ for $L_n(\eta)$, $U_n(\eta)$, and $w_n(\eta)$. Recall from \cref{def:pre-quantile,lem:pre-quantile} that $0 \le L \le U \le M$. For an event $A$ and a nonnegative or integrable expression $X$ defined on $A$, we define $\E[X; A]$ to be the integral of $X$ over $A$. When we write this as $\E[X\ind_A]$, we first extend $X$ by zero outside $A$. We define the capped variable
\begin{equation}\label{eq:red-Y}
Y = \min(U, \max(L, V)).
\end{equation}
Since $U \le M$ and $V = \min(H - 1, M)$, we have $Y = \min(U, \max(L, H - 1))$. Checking the three cases $H \le L + 1$, $L + 1 < H \le U$, and $H > U$ separately, we find
\begin{equation*}
\#\{m \in \{L + 1, \dots, U\} : m < H\} = \min(U, \max(L, H - 1)) - L.
\end{equation*}
Since $G_m$ is colorable if and only if $m < H$, this implies
\begin{equation}\label{eq:red-Ysum}
Y = L + \sum_{m = L + 1}^{U} \ind\{G_m \text{ is colorable}\}.
\end{equation}
For $1 \le i \le U$ and $m \ge 0$, we let $G_m^{(-i)}$ denote the multiset $\{e_j : 1 \le j \le m, \ j \ne i\}$. We have $G_m^{(-i)} = G_m$ for $m < i$, and for $m \ge i$ the multigraph $G_m^{(-i)}$ is obtained from $G_m$ by omitting the occurrence $e_i$ (the other occurrences keep their indices). Following \eqref{eq:red-Ysum}, we define
\begin{equation}\label{eq:red-Ydel}
Y^{(-i)} = L + \sum_{m = L + 1}^{U} \ind\{G_m^{(-i)} \text{ is colorable}\}.
\end{equation}
The variable $Y^{(-i)}$ is a function of the pairs $e_j$ with $j \le U$ and $j \ne i$.

The next lemma gives an exact identity for the sum of the squared differences $Y^{(-i)} - Y$. Bounds of this kind have antecedents. In \cite[Theorem~1(b)]{aldous1991threshold}, Aldous considered the first time $T$ at which independent copies of a random subset of a finite set cover that set. He bounded the variance of $T$ by an absolute constant times $\E T$ times the expected time that a fresh independent sequence of such subsets needs to cover the set left uncovered by the first $T - 1$ of them. This theorem applies to $H$. Indeed, let each pair $\{u, v\}$ cover the maps $\sigma \colon [n] \to [q]$ with $\sigma(u) = \sigma(v)$; since $n > q$, every map is covered by some pair. Then $H$ is the first time at which $e_1, \dots, e_H$ cover all maps, the maps left uncovered by $e_1, \dots, e_{H - 1}$ form the set $\Col(\Gamma)$, and the expected time to cover this set is $r(\Gamma)$. The theorem therefore bounds $\Var(H)$ by an absolute constant times $\E H \, \E r(\Gamma)$.

For random $k$-SAT, Carenini's 2026 preprint applies Aldous's bound to the expected number of further independent clauses needed to make the final satisfiable formula unsatisfiable \cite[Theorem~4.1 and equation~(4.1)]{carenini2026polynomial}. The limiting-threshold preprint \cite[Lemmas~3.1 and~3.2]{openai2026limiting} applies the deletion form of the Efron--Stein inequality to a capped hitting time, retains the original clause indices, and expands the squares of deletion delays over pairs of times. The variance preprint \cite[Section~3.3, Lemma~3.4]{openai2026variance} also combines these delays with expected lifetimes. \Cref{lem:red-deletion} is the analogous computation for colorability. Instead of the bound for $\Var(H)$, we use the exact identity \eqref{eq:red-deletion}, in which only final edges $e_H$ with $L < H \le U$ appear, and we give its short proof.

\begin{lemma}\label{lem:red-deletion}
Let $\eta \in (0, \eta_0]$ and let $n \ge 2q$ be an integer. Then
\begin{equation}\label{eq:red-deletion}
\sum_{i = 1}^{U} \E\bigl[(Y^{(-i)} - Y)^2\bigr] = \E\bigl[H \bigl(2 r_{U - H + 1}(\Gamma) - 1\bigr); L < H \le U\bigr]
\end{equation}
and
\begin{equation}\label{eq:red-var-upper}
\Var(Y) \le 2U \, \E\bigl[\min(w, r(\Gamma)); L < H \le U\bigr].
\end{equation}
\end{lemma}

\begin{proof}
We begin by computing the differences $Y^{(-i)} - Y$. Fix $1 \le i \le U$, and for $m \ge i$ let
\begin{equation*}
\cX_{m, i} = \{G_m^{(-i)} \text{ is colorable and } G_m \text{ is not colorable}\}.
\end{equation*}
For $m < i$, we have $G_m^{(-i)} = G_m$. For $m \ge i$, we have $G_m = G_m^{(-i)} + e_i$, and every coloring of $G_m$ is a coloring of $G_m^{(-i)}$, so that
\begin{equation*}
\ind\{G_m^{(-i)} \text{ is colorable}\} - \ind\{G_m \text{ is colorable}\} = \ind_{\cX_{m, i}}.
\end{equation*}
Subtracting \eqref{eq:red-Ysum} from \eqref{eq:red-Ydel}, we have
\begin{equation}\label{eq:red-diff}
Y^{(-i)} - Y = \sum_{m = \max(i, L + 1)}^{U} \ind_{\cX_{m, i}}.
\end{equation}

Next, we square \eqref{eq:red-diff}. Let $m < l$ be indices in the range of the sum. The multigraph $G_l$ contains $G_m$, so it is not colorable on $\cX_{m, i}$, and hence
\begin{equation*}
\cX_{m, i} \cap \cX_{l, i} = \cX_{m, i} \cap \{G_l^{(-i)} \text{ is colorable}\}.
\end{equation*}
Expanding the square gives
\begin{equation}\label{eq:red-square}
(Y^{(-i)} - Y)^2 = \sum_{m = \max(i, L + 1)}^{U} \Bigl(\ind_{\cX_{m, i}} + 2 \sum_{l = m + 1}^{U} \ind_{\cX_{m, i} \cap \{G_l^{(-i)} \text{ is colorable}\}}\Bigr).
\end{equation}

The main step is to compute the conditional probabilities of these events given $G_m^{(-i)}$. Fix $m \ge i$ and $l > m$. The multigraph $G_m^{(-i)}$ is a function of the pairs $e_j$ with $j \le m$ and $j \ne i$, and we have
\begin{equation*}
G_m = G_m^{(-i)} + e_i, \qquad G_l^{(-i)} = G_m^{(-i)} + e_{m + 1} + \dots + e_l.
\end{equation*}
The pairs $e_i, e_{m + 1}, \dots, e_l$ are independent, uniform, and independent of $G_m^{(-i)}$.

Let $G$ be a colorable multigraph on $[n]$ with $\Pp(G_m^{(-i)} = G) > 0$. By \cref{lem:pre-frozen}(a), the multigraph $G + e_i$ is not colorable if and only if $e_i$ is a frozen pair of $G$, and this has conditional probability $\kappa(G)$.

The definition of $T_G$ in \cref{def:pre-lifetime} shows that $G + e_{m + 1} + \dots + e_l$ is colorable with conditional probability $\Pp(T_G > l - m)$. Using the independence of $e_i$ from $e_{m + 1}, \dots, e_l$, we find that
\begin{equation}\label{eq:red-cond}
\begin{gathered}
\Pp\bigl(\cX_{m, i} \bigm| G_m^{(-i)} = G\bigr) = \kappa(G), \\
\Pp\bigl(\cX_{m, i} \cap \{G_l^{(-i)} \text{ is colorable}\} \bigm| G_m^{(-i)} = G\bigr) = \kappa(G) \Pp(T_G > l - m).
\end{gathered}
\end{equation}
The two probabilities in \eqref{eq:red-cond} vanish if $G$ is not colorable, and so does $\kappa(G)$.

For $m \in \{L + 1, \dots, U\}$ and a multigraph $G$ on $[n]$, we set
\begin{equation*}
h'_m(G) = \ind\{G \text{ is colorable}\} \bigl(2 r_{U - m + 1}(G) - 1\bigr).
\end{equation*}
For colorable $G$, we have $\Pp(T_G > 0) = 1$, and \eqref{eq:pre-rb} and \cref{lem:pre-lifetime}(a) give
\begin{equation}\label{eq:red-weight}
1 + 2 \sum_{l = m + 1}^{U} \Pp(T_G > l - m) = 2 r_{U - m + 1}(G) - 1 = h'_m(G) \ge 1.
\end{equation}
The indicator makes $h'_m$ nonnegative, as \cref{lem:pre-terminal} requires. Taking expectations in \eqref{eq:red-square} and inserting \eqref{eq:red-cond} and \eqref{eq:red-weight}, we obtain
\begin{equation}\label{eq:red-one-i}
\E\bigl[(Y^{(-i)} - Y)^2\bigr] = \sum_{m = \max(i, L + 1)}^{U} \E\bigl[\kappa(G_m^{(-i)}) h'_m(G_m^{(-i)})\bigr].
\end{equation}

We now sum over $i$. For fixed $m \in \{L + 1, \dots, U\}$, the terms of \eqref{eq:red-one-i} with index $m$ occur exactly for $1 \le i \le m$. For each such $i$, $G_m^{(-i)}$ is the multiset of $m - 1$ independent uniform pairs, and it has the law of $G_{m - 1}$. Exchanging the two sums, we find
\begin{equation}\label{eq:red-sum-i}
\sum_{i = 1}^{U} \E\bigl[(Y^{(-i)} - Y)^2\bigr] = \sum_{m = L + 1}^{U} m \, \E\bigl[\kappa(G_{m - 1}) h'_m(G_{m - 1})\bigr].
\end{equation}

To evaluate the right side of \eqref{eq:red-sum-i}, let $\cE$ denote the set of all pairs $(G, e)$ such that $G$ is a colorable multigraph on $[n]$ and $e$ is a frozen pair of $G$. Then $\kappa_\cE = \kappa$. By \cref{lem:pre-frozen}(a), the final edge $e_H$ is a frozen pair of $\Gamma$ whenever $H < \infty$, and then $(\Gamma, e_H) \in \cE$. Applying \cref{lem:pre-terminal}(a) with this $\cE$, with $L' = m - 1$ and $U' = m$, and with $h_{m - 1} = h'_m$, we have
\begin{equation*}
\E\bigl[\kappa(G_{m - 1}) h'_m(G_{m - 1})\bigr] = \E\bigl[h'_m(\Gamma); H = m\bigr].
\end{equation*}
Inserting this into \eqref{eq:red-sum-i}, and using that $\Gamma$ is colorable, we arrive at \eqref{eq:red-deletion}. This is an identity of expectations. An individual deletion delay need not vanish on a trajectory with $H \le L$.

To prove \eqref{eq:red-var-upper}, we use the Efron--Stein inequality \cite{efron1981jackknife,steele1986efron} in the following form, stated in \cite[(3.6)]{boucheron2005moment} and \cite[Theorem~3.1]{boucheron2013concentration}. Let $X_1, \dots, X_k$ be independent random variables, let $Z$ be a square-integrable function of $(X_1, \dots, X_k)$, and for each $i$ let $Z_i$ be a square-integrable random variable that is measurable with respect to $\sigma(X_j : j \ne i)$. Then
\begin{equation*}
\Var(Z) \le \sum_{i = 1}^{k} \E\bigl[(Z - Z_i)^2\bigr].
\end{equation*}

By \eqref{eq:red-Ysum}, the variable $Y$ is a bounded function of the independent pairs $e_1, \dots, e_U$, and $Y^{(-i)}$ is a bounded function of the pairs $e_j$ with $j \le U$ and $j \ne i$. Applying the inequality with $Z = Y$ and $Z_i = Y^{(-i)}$, and using \eqref{eq:red-deletion}, we find
\begin{equation*}
\Var(Y) \le \E\bigl[H \bigl(2 r_{U - H + 1}(\Gamma) - 1\bigr); L < H \le U\bigr].
\end{equation*}
We have $H \le U$ and $1 \le U - H + 1 \le w$ on the event $\{L < H \le U\}$, and therefore
\begin{equation*}
r_{U - H + 1}(\Gamma) \le r_w(\Gamma) \le \min\bigl(w, r(\Gamma)\bigr),
\end{equation*}
where we used that $\Gamma$ is colorable and that, by \cref{lem:pre-lifetime}(a), $r_b(G)$ is nondecreasing in $b$ and bounded by $\min(b, r(G))$ for colorable $G$ and integers $b \ge 1$. Inserting this into the previous bound and dropping the term $-1$ gives \eqref{eq:red-var-upper}.
\end{proof}

We will also need the following elementary lower bound on the variance of $Y$.

\begin{lemma}\label{lem:red-varlower}
Let $\eta \in (0, \eta_0]$ and let $n \ge 2q$ be an integer. Then $\Var(Y) \ge \eta^2 w^2$.
\end{lemma}

\begin{proof}
If $w = 0$, there is nothing to prove. We may therefore assume that $L < U$. Let $Y'$ be an independent copy of $Y$. We have $|Y - Y'| = w$ on each of the disjoint events $\{Y = L, \ Y' = U\}$ and $\{Y = U, \ Y' = L\}$, and therefore
\begin{equation*}
\Var(Y) = \frac{1}{2} \E\bigl[(Y - Y')^2\bigr] \ge \Pp(Y = L) \Pp(Y = U) w^2.
\end{equation*}
Because $L < U$, the definition \eqref{eq:red-Y} gives
\begin{equation*}
\{Y = L\} = \{V \le L\}, \qquad \{Y = U\} = \{V \ge U\}.
\end{equation*}
Recall from the definitions of $L$ and $U$ in \cref{def:pre-quantile} that $\Pp(V \le L) \ge \eta$ and $\Pp(V \ge U) \ge \eta$, which completes the proof.
\end{proof}

Combining the two lemmas, we bound the window by the expected lifetime of $\Gamma$ at final edges inside the window, capped at $n \eps_n$.

\begin{proposition}\label{prop:red-window}
Let $\eta \in (0, \eta_0]$, let $n \ge 2q$ be an integer, and let $K_0 \ge 0$. Suppose that
\begin{equation}\label{eq:red-hyp}
\E\bigl[\min(n \eps_n, r(\Gamma)); L < H \le U\bigr] \le K_0 n \eps_n^2.
\end{equation}
Then $w \le \max(1, 2 \Lambda K_0/\eta^2) n \eps_n$.
\end{proposition}

\begin{proof}
If $w < n \eps_n$, there is nothing to prove. We therefore suppose that $w \ge n \eps_n$. Then $w > 0$, and for all real $y \ge 0$ we have
\begin{equation}\label{eq:red-minscale}
\min(w, y) \le \frac{w}{n \eps_n} \min(n \eps_n, y).
\end{equation}
Indeed, if $y \le n \eps_n$, the right side is at least $y$, and if $y > n \eps_n$, it equals $w$. Combining \cref{lem:red-varlower}, \eqref{eq:red-var-upper}, \eqref{eq:red-minscale}, and \eqref{eq:red-hyp}, we find
\begin{equation*}
\begin{split}
\eta^2 w^2 \le \Var(Y) &\le 2U \, \E\bigl[\min(w, r(\Gamma)); L < H \le U\bigr] \\
&\le \frac{2Uw}{n \eps_n} \E\bigl[\min(n \eps_n, r(\Gamma)); L < H \le U\bigr] \le 2 \Lambda K_0 n \eps_n w,
\end{split}
\end{equation*}
where in the last step we also used $U \le M \le \Lambda n$. Dividing by $\eta^2 w$ gives $w \le (2 \Lambda K_0/\eta^2) n \eps_n$.
\end{proof}

\section{Splitting by the degrees of the endpoints}\label{sec:red-split}

We now split the expectation in \eqref{eq:red-hyp} according to the degrees in $\Gamma$ of the two endpoints of $e_H$. Recall that degrees count occurrences with multiplicity. When some endpoint has degree at most $D_n$, we keep the restriction $r(\Gamma) > x_n$; it is used in \cref{sec:low}, through \cref{lem:pre-classmass}, to bound the size of the frozen class of that endpoint. When both endpoints have degree larger than $D_n$, however, the estimate of \cref{sec:high} holds without a restriction on the lifetime.

\begin{proposition}\label{prop:red-split}
Let $\eta \in (0, \eta_0]$ and let $n \ge 2q$ be an integer. Suppose that
\begin{equation}\label{eq:red-low}
\Pp\Bigl(L < H \le U, \ \min_{v \in e_H} \deg_\Gamma(v) \le D_n, \ r(\Gamma) > x_n\Bigr) \le \eps_n
\end{equation}
and
\begin{equation}\label{eq:red-high}
\Pp\Bigl(L < H \le U, \ \min_{v \in e_H} \deg_\Gamma(v) > D_n, \ r(\Gamma) > x_n\Bigr) \le \eps_n.
\end{equation}
Then
\begin{equation*}
\E\bigl[\min(n \eps_n, r(\Gamma)); L < H \le U\bigr] \le 3 n \eps_n^2 \qquad \text{and} \qquad w \le \frac{6 \Lambda}{\eta^2} n \eps_n.
\end{equation*}
\end{proposition}

\begin{proof}
For every real $y \ge 0$, we have
\begin{equation*}
\min(n \eps_n, y) \le x_n + n \eps_n \ind\{y > x_n\},
\end{equation*}
since the left side is bounded by $y \le x_n$ if $y \le x_n$, and by $n \eps_n$ otherwise. Exactly one of the two degree conditions in \eqref{eq:red-low} and \eqref{eq:red-high} holds on the event $\{L < H \le U, \ r(\Gamma) > x_n\}$. Using this fact, \eqref{eq:red-low}, \eqref{eq:red-high}, and $x_n = n \eps_n^2$, we find that
\begin{equation*}
\begin{split}
\E\bigl[\min(n \eps_n, r(\Gamma)); L < H \le U\bigr] &\le x_n \Pp(L < H \le U) + n \eps_n \Pp\bigl(L < H \le U, \ r(\Gamma) > x_n\bigr) \\
&\le n \eps_n^2 + 2 n \eps_n^2.
\end{split}
\end{equation*}
This is the first claim. The second claim follows from \cref{prop:red-window} with $K_0 = 3$, since $6 \Lambda/\eta^2 \ge 1$.
\end{proof}

The hypotheses of \cref{prop:red-split} are verified in the next two sections, by estimates that hold for all large $n$. In \cref{sec:low}, we prove \cref{thm:low-main}, and \cref{cor:low} gives, for all large $n$,
\begin{equation}\label{eq:red-low-target}
\Pp\Bigl(H < \infty, \ \min_{v \in e_H} \deg_\Gamma(v) \le D_n, \ r(\Gamma) > x_n\Bigr) \le \eps_n.
\end{equation}
This bound concerns final edges at all times. It implies \eqref{eq:red-low} for every $\eta \in (0, \eta_0]$, since the event in \eqref{eq:red-low} is contained in the event in \eqref{eq:red-low-target}. In \cref{sec:high}, we prove \cref{thm:high-main}, which gives, for all large $n$ and all integers $0 \le L' \le U' \le M$,
\begin{equation}\label{eq:red-high-target}
\Pp\Bigl(L' < H \le U', \ \min_{v \in e_H} \deg_\Gamma(v) > D_n\Bigr) \le \eps_n.
\end{equation}
Taking $L' = L$ and $U' = U$, we see that \eqref{eq:red-high-target} implies \eqref{eq:red-high}. The two estimates can then be combined with \cref{prop:red-split,thm:red-convergence}.

\chapter{Final edges at a vertex of low degree}\label{sec:low}

In this section, we prove \cref{thm:low-main}. It bounds the probability that the final edge $e_H$ has an endpoint of degree at most $D_n$ in $\Gamma$ and that $r(\Gamma)$ exceeds a prescribed positive level. At the level $x_n$, it gives \eqref{eq:red-low-target}, which implies the first hypothesis of \cref{prop:red-split}; this is \cref{cor:low}. The heuristics are explained in \cref{sec:overview}.

The proof has a structural part and a probabilistic part. For a multigraph $G$ and a vertex $a$, we call any list of the other endpoints of its occurrences at $a$, with one entry per occurrence, a \emph{neighbor list of $a$ in $G$}. The structural result, \cref{lem:low-frozen}(iii), shows that if $a$ lies in a frozen pair of $G$, then every coloring of the multigraph obtained by removing $a$ uses at least two colors on this list. This is the only step of the window estimates that uses $q \ge 3$, and \cref{rem:low-q2} shows that it fails for $q = 2$. (The other use of $q \ge 3$ in the paper is \cref{lem:pre-centre}(b), which is used in \cref{app:conv}.)

Let $B$ be a colorable multigraph, and let $\mathbf x$ be a nonempty list of its vertices. Then every coloring of $B$ uses at least two colors on $\mathbf x$ exactly when the multigraph obtained by identifying its set of entries is not colorable. This is the classical correspondence between the colorings of an identification and the colorings that are constant on the identified set; for two nonadjacent vertices, Read uses it in the proof of \cite[Section~3, Theorem~1]{read1968introduction} to split the chromatic polynomial. Further, Poisson splitting at a vertex of degree at most $D_n$ makes its neighbor list a short list of independent uniform vertices, independent of the rest of the graph. We therefore compare the hitting time of non-colorability of an edge process with the hitting time of non-colorability after a few random vertices are identified. This comparison, \cref{prop:low-identify}, is the main estimate of this section. In its proof, we compare the vertex obtained by identification with the vertex added to an edge process by the construction of \cref{lem:low-added}. In \cref{sec:low-main}, we combine these results with the bound on frozen classes of \cref{lem:pre-classmass} to prove \cref{thm:low-main}.

\section{Frozen pairs and neighbor lists}\label{sec:low-frozen}

We first fix notation for lists and identifications. Let $S$ be a finite set. Recall from \cref{sec:pre-poisson} that a list of elements of $S$ is a finite sequence of elements of $S$, possibly with repeated entries. We denote a list of length $k$ by $\mathbf x = (\mathbf x(1), \dots, \mathbf x(k))$. For $0 \le i \le k$, we let $\mathbf x_{[i]} = (\mathbf x(1), \dots, \mathbf x(i))$ denote its prefix of length $i$; the prefix $\mathbf x_{[0]}$ is the empty list. We use the same notation for infinite sequences.

For a map $\sigma \colon S \to [q]$, we let $\sigma(\mathbf x) = \{\sigma(\mathbf x(i)) : 1 \le i \le k\}$ denote the set of colors on $\mathbf x$. We say that $\sigma$ is constant on $\mathbf x$ if $|\sigma(\mathbf x)| \le 1$, and that $\sigma$ is constant on a set $Z \subset S$ if it takes a single value on $Z$. In particular, every map is constant on the empty list.

For a vertex $o \notin S$, we let $\operatorname{star}(o, \mathbf x)$ denote the multigraph on $S \cup \{o\}$ with the $k$ occurrences $\{o, \mathbf x(1)\}, \dots, \{o, \mathbf x(k)\}$. For a multigraph $B$ on $S$, we let $B + \operatorname{star}(o, \mathbf x)$ denote the multigraph on $S \cup \{o\}$ whose occurrences are those of $B$ and those of $\operatorname{star}(o, \mathbf x)$.

\begin{definition}\label{def:low-nm}
Let $B$ be a colorable multigraph on a finite set $S$. We say that a list $\mathbf x$ of elements of $S$ is \emph{never monochromatic} in $B$ if no coloring of $B$ is constant on $\mathbf x$, that is, if $|\sigma(\mathbf x)| \ge 2$ for every $\sigma \in \Col(B)$.
\end{definition}

Identifying a set of vertices that contains a pair occurring in a multigraph produces a loop, so we allow loops in the following definition.

\begin{definition}\label{def:low-identification}
Let $B$ be a multigraph on a finite set $S$, let $Z \subset S$ be nonempty, and let $o \notin S$. Let $\iota_Z \colon S \to (S \setminus Z) \cup \{o\}$ be the map with $\iota_Z(v) = o$ for $v \in Z$ and $\iota_Z(v) = v$ for $v \notin Z$. We define the \emph{identification} $B/Z$ to be the multiset that contains, for each occurrence $\{u, v\}$ of $B$, one occurrence $\{\iota_Z(u), \iota_Z(v)\}$. For a finite set $T$, we define a \emph{multigraph with loops on $T$} to be a finite multiset of sets $\{u, v\}$ with $u, v \in T$, allowing $u = v$. We call an occurrence $\{u, u\}$ a \emph{loop}. We call a map $\sigma'\colon T \to [q]$ a \emph{coloring} of such a multigraph if $\sigma'(u) \ne \sigma'(v)$ for every occurrence $\{u, v\}$. We let $\Col$ denote its set of colorings, and call it \emph{colorable} if that set is nonempty. We view $B/Z$ as a multigraph with loops on $(S \setminus Z) \cup \{o\}$.
\end{definition}

The identification $B/Z$ keeps all multiplicities. An occurrence of $B$ with exactly one endpoint in $Z$ becomes an occurrence at $o$, and an occurrence that is a pair of $Z$ becomes a loop. A multigraph with a loop is not colorable, and a loopless $B/Z$ is a multigraph on $(S \setminus Z) \cup \{o\}$ in the sense of \cref{def:pre-multigraph}. (The same convention for loops is used in \cref{app:conv}.)

The following lemma collects the facts about frozen pairs and identifications that we need. Part (iv) is the correspondence mentioned at the start of this section. The argument for part (iii) recolors a single vertex and is elementary, but its conclusion requires $q \ge 3$.

\begin{lemma}\label{lem:low-frozen}
Let $S$ be a finite set, let $G$ be a multigraph on $S$, and let $a \in S$. Let $B = G[S \setminus \{a\}]$ be the multigraph of the occurrences of $G$ that do not contain $a$, and let $\mathbf x$ be a list of the other endpoints of the occurrences of $G$ at $a$, with one entry for each occurrence, in an arbitrary order. We let $k$ denote the length of $\mathbf x$. Then $k = \deg_G(a)$ and $G = B + \operatorname{star}(a, \mathbf x)$.
\begin{enumerate}[label=(\roman*)]
\item Let $\sigma'\colon S \to [q]$, and let $\sigma$ denote its restriction to $S \setminus \{a\}$. Then $\sigma'$ is a coloring of $G$ if and only if $\sigma$ is a coloring of $B$ and $\sigma'(a) \notin \sigma(\mathbf x)$.
\item Suppose that $G$ is colorable, and let $b \in S \setminus \{a\}$. Then $a \sim_G b$ if and only if $B + \operatorname{star}(a, (\mathbf x(1), \dots, \mathbf x(k), b))$ is not colorable.
\item Suppose that $a \sim_G b$ for some $b \in S \setminus \{a\}$. Then $B$ is colorable, and $|\sigma(\mathbf x)| \ge q - 1$ for every coloring $\sigma$ of $B$. In particular, if $q \ge 3$, then $\mathbf x$ is never monochromatic in $B$.
\item Let $F$ be a colorable multigraph on $S$, and let $Z \subset S$ be nonempty. Then $F/Z$ is colorable if and only if some coloring of $F$ is constant on $Z$. In particular, a nonempty list of elements of $S$ whose set of entries is $Z$ is never monochromatic in $F$ if and only if $F/Z$ is not colorable.
\end{enumerate}
\end{lemma}

\begin{proof}
For (i), note that an occurrence of $G$ either does not contain $a$, in which case it is an occurrence of $B$, or it has the form $\{a, \mathbf x(i)\}$ for some $i$. Therefore $\sigma'$ gives distinct colors to the endpoints of all occurrences of $G$ if and only if $\sigma \in \Col(B)$ and $\sigma'(a) \ne \sigma(\mathbf x(i))$ for all $i$.

To prove (ii), observe that
\begin{equation*}
B + \operatorname{star}\bigl(a, (\mathbf x(1), \dots, \mathbf x(k), b)\bigr) = G + \{a, b\}.
\end{equation*}
Since $G$ is colorable and $a \ne b$, \cref{lem:pre-frozen}(a) shows that $G + \{a, b\}$ is not colorable if and only if $a \sim_G b$.

For (iii), recall that $a \sim_G b$ requires $G$ to be colorable. By (i), $B$ is then colorable.

Suppose, for contradiction, that some coloring $\sigma$ of $B$ satisfies $|\sigma(\mathbf x)| \le q - 2$. Then there exist two colors $a' \ne b'$ outside $\sigma(\mathbf x)$. We define $\sigma'$ and $\sigma''$ to agree with $\sigma$ on $S \setminus \{a\}$ and to satisfy $\sigma'(a) = a'$ and $\sigma''(a) = b'$. By (i), both maps are colorings of $G$.

Using $b \ne a$ and $a \sim_G b$, we find
\begin{equation*}
a' = \sigma'(a) = \sigma'(b) = \sigma(b) = \sigma''(b) = \sigma''(a) = b',
\end{equation*}
a contradiction. We conclude that $|\sigma(\mathbf x)| \ge q - 1$ for all $\sigma \in \Col(B)$. For $q \ge 3$, this gives $|\sigma(\mathbf x)| \ge 2$, and $\mathbf x$ is never monochromatic in $B$.

For (iv), suppose first that $\sigma'$ is a coloring of $F/Z$. Then $F/Z$ has no loop, and for all occurrences $\{u, v\}$ of $F$, the pair $\{\iota_Z(u), \iota_Z(v)\}$ is an occurrence of $F/Z$. This shows that $\sigma' \circ \iota_Z$ is a coloring of $F$ that is constant on $Z$.

Conversely, let $\sigma$ be a coloring of $F$ that takes the value $a'$ on $Z$. Then no occurrence of $F$ is a pair of $Z$, and $F/Z$ has no loop. We define $\sigma'(o) = a'$ and $\sigma'(v) = \sigma(v)$ for $v \in S \setminus Z$. Then $\sigma = \sigma'\circ\iota_Z$. For all occurrences $\{u, v\}$ of $F$, we have
\begin{equation*}
\sigma'(\iota_Z(u)) = \sigma(u) \ne \sigma(v) = \sigma'(\iota_Z(v)),
\end{equation*}
so $\sigma'$ is a coloring of $F/Z$. The last statement of (iv) follows, since a map is constant on a list if and only if it is constant on the set of entries of the list.
\end{proof}

The proof of (iv) shows that $\sigma' \mapsto \sigma' \circ \iota_Z$ is a bijection from $\Col(F/Z)$ onto the set of colorings of $F$ that are constant on $Z$. Its first half applies to every multigraph $F$ on $S$, colorable or not, and shows that $F$ is colorable whenever $F/Z$ is colorable.

\begin{remark}\label{rem:low-q2}
The assumption $q \ge 3$ in the last statement of \cref{lem:low-frozen}(iii) cannot be removed. Suppose that $q = 2$. Let $S$ be a finite set containing distinct vertices $a, v, b$, and define $G$ on $S$ to have the occurrences $\{a, v\}$ and $\{v, b\}$. We set $\mathbf x = (v)$ and $B = G[S \setminus \{a\}]$. Every coloring of $G$ gives $a$ and $b$ the color that differs from the color of $v$, so $a \sim_G b$. However, $\mathbf x$ has length one, and every coloring of $B$ is constant on it. More generally, for $q = 2$, the restriction of a coloring of a colorable $G$ to $S \setminus \{a\}$ is a coloring of $B$ that gives the same color to all neighbors of $a$. Thus, for $q = 2$, the conclusion of (iii) fails for every vertex $a$ of positive degree in every colorable $G$. The first assertion of (iii) remains true for $q = 2$, but it only says that $\mathbf x$ uses at least one color. The argument of \cref{thm:low-main} therefore fails for $q = 2$ at the step where it uses \cref{lem:low-frozen}(iii).
\end{remark}

\section{Identification of random vertices}\label{sec:low-identify}

Let $N \ge 2q$ be an integer. Let $(B_m)_{m \ge 0}$ be the time-$m$ graphs of an edge process on $[N]$, as in \cref{def:pre-process}. We let $H^{(N)} = \min\{m \ge 0 : B_m \text{ is not colorable}\}$ denote its hitting time of non-colorability. By \cref{def:pre-lifetime}, the variable $H^{(N)}$ has the law of $T_G$ for the empty multigraph $G$ on $[N]$, and \cref{lem:pre-lifetime}(a) implies that $\E H^{(N)} \le A_q N$. In particular, $H^{(N)}$ is almost surely finite.

For $D \in \{1, \dots, N\}$, we let $Z$ be a uniform random $D$-subset of $[N]$, independent of $(B_m)$. We fix a vertex $o \notin [N]$, and we define
\begin{equation*}
H_Z = \min\{m \ge 0 : B_m/Z \text{ is not colorable}\}.
\end{equation*}
Since $B_m \subset B_{m + 1}$, we have $B_m/Z \subset B_{m + 1}/Z$, so $B_m/Z$ is not colorable for all $m \ge H_Z$. The remark after the proof of \cref{lem:low-frozen} implies that the multigraph $B_m/Z$ is not colorable whenever $B_m$ is not colorable. Thus $H_Z \le H^{(N)}$. We conclude that, almost surely,
\begin{equation}\label{eq:low-advance}
H^{(N)} - H_Z = \sum_{m \ge 0} \ind\{B_m \text{ is colorable and } B_m/Z \text{ is not colorable}\}.
\end{equation}
We set
\begin{equation}\label{eq:low-constants}
z_0 = 4q \log q + 1, \qquad C_q = 2 e^{z_0}, \qquad N_q = \lceil 4q \log q \rceil + 1.
\end{equation}
Then $N_q > 2q$.

We first treat a vertex that is added to an edge process in the same way as the other vertices. In the proof of \cref{prop:low-identify}, the vertex obtained by identification will be compared with such a vertex.

\begin{lemma}\label{lem:low-added}
Let $V'$ be a finite set with $N' \ge 2$ elements, let $o \notin V'$, and let $\lambda > 0$. For $t \ge 0$, let $F(t)$ denote the Poissonized process on $V' \cup \{o\}$ at time $t$, and set $\lambda' = \lambda \binom{N' + 1}{2}$. Let $s \ge 0$, and let $B'$ be a random multigraph on $V'$ whose numbers of occurrences on the pairs of $V'$ are independent with the $\Pois(\lambda s)$ law. Let $\mathbf x = (\mathbf x(1), \mathbf x(2), \dots)$ be a sequence of independent uniform elements of $V'$, and let $J$ have the $\Pois(\lambda N' s)$ law, with $B'$, $\mathbf x$, and $J$ independent. We define
\begin{equation*}
I_\star = \min\bigl\{i \ge 0 : B' + \operatorname{star}(o, \mathbf x_{[i]}) \text{ is not colorable}\bigr\}, \qquad B'' = B' + \operatorname{star}(o, \mathbf x_{[J]}),
\end{equation*}
with $I_\star = \infty$ if there is no such $i$. Then $B''$ has the law of $F(\lambda' s)$, and
\begin{equation}\label{eq:low-added}
\Pp(B' \text{ is colorable}, \ I_\star \le J) \le (e^{\lambda N' s} - 1) \E \kappa\bigl(F(\lambda' s)\bigr).
\end{equation}
Moreover,
\begin{equation}\label{eq:low-added-integral}
\int_0^\infty \E \kappa\bigl(F(\lambda' s)\bigr) \dd s \le \frac{1}{\lambda'}.
\end{equation}
\end{lemma}

\begin{proof}
We begin with the law of $B''$. The numbers of occurrences of $B''$ on the pairs $\{o, v\}$, $v \in V'$, are the numbers of entries of $\mathbf x_{[J]}$ equal to $v$. Applying \cref{lem:pre-split}(a) with the partition of $V'$ into single vertices, we see that these numbers are independent with the $\Pois(\lambda s)$ law. They are also independent of $B'$.

The numbers of occurrences of $B''$ on the $\binom{N' + 1}{2}$ pairs of $V' \cup \{o\}$ are independent with the $\Pois(\lambda s)$ law. Applying \cref{lem:pre-split}(a) with the partition of the pairs of $V' \cup \{o\}$ into single pairs, we see that the numbers of occurrences of $F(\lambda' s)$ on these pairs have the same joint law. Since a multigraph is determined by these numbers, $B''$ has the law of $F(\lambda' s)$.

Next, we prove \eqref{eq:low-added}. We set $\hat z = \lambda N's$. We have $I_\star \ge 1$ on the event that $B'$ is colorable. Since $J$ is independent of $(B', \mathbf x)$ and $I_\star$ is a function of $(B', \mathbf x)$, \cref{lem:pre-tails}(d) gives
\begin{equation}\label{eq:low-shift}
\begin{split}
\Pp(B' \text{ is colorable}, \ I_\star \le J) &= \sum_{i \ge 1} \Pp(B' \text{ is colorable}, \ I_\star = i) \Pp(J \ge i) \\
&\le (e^{\hat z} - 1) \sum_{i \ge 1} \Pp(B' \text{ is colorable}, \ I_\star = i) \pi_{\hat z}(i - 1).
\end{split}
\end{equation}

We now identify the last sum in \eqref{eq:low-shift}. Using once more that $J$ is independent of $(B', \mathbf x)$, we see that it equals $\Pp(B' \text{ is colorable}, \ I_\star = J + 1)$. By the definition of $I_\star$, this is the probability that $B''$ is colorable and $B' + \operatorname{star}(o, \mathbf x_{[J + 1]})$ is not colorable. We apply \cref{lem:low-frozen}(ii) to the multigraph $B''$, its vertex $o$, and $b = \mathbf x(J + 1)$, and we find that this event is the event that $B''$ is colorable and $o \sim_{B''} \mathbf x(J + 1)$.

Conditional on $J$, the entry $\mathbf x(J + 1)$ is uniform on $V'$ and independent of $B''$. This implies that
\begin{equation}\label{eq:low-class}
\Pp(B' \text{ is colorable}, \ I_\star = J + 1) = \frac{1}{N'} \E\bigl[\ind\{B'' \text{ is colorable}\} \bigl(|C_{B''}(o)| - 1\bigr)\bigr].
\end{equation}

The law of $F(\lambda' s)$ is invariant under permutations of $V' \cup \{o\}$, so the expectation in \eqref{eq:low-class} does not change if $o$ is replaced by any other vertex of $V' \cup \{o\}$. Recall from \cref{def:pre-frozen} that for a colorable multigraph $F'$ on $V' \cup \{o\}$, the number $\kappa(F')$ is the number of frozen pairs of $F'$ divided by $\binom{N' + 1}{2}$, so that
\begin{equation*}
\sum_{v \in V' \cup \{o\}} \bigl(|C_{F'}(v)| - 1\bigr) = 2 \binom{N' + 1}{2} \kappa(F') = (N' + 1) N' \kappa(F').
\end{equation*}
Averaging the right side of \eqref{eq:low-class} over the $N' + 1$ vertices in place of $o$, we obtain
\begin{equation*}
\Pp(B' \text{ is colorable}, \ I_\star = J + 1) = \E \kappa\bigl(F(\lambda' s)\bigr).
\end{equation*}
Inserting this into \eqref{eq:low-shift}, we arrive at \eqref{eq:low-added}.

It remains to prove \eqref{eq:low-added-integral}. Let $(F_m)_{m \ge 0}$ be the time-$m$ graphs of the edge process on $V' \cup \{o\}$ from which $F(t)$ is built. The $(m + 1)$-st pair of this process is uniform and independent of $F_m$, and \cref{lem:pre-frozen}(a) shows that $\E \kappa(F_m)$ is the probability that $F_m$ is colorable and $F_{m + 1}$ is not colorable. These events are disjoint for different $m$, and therefore
\begin{equation*}
\sum_{m \ge 0} \E \kappa(F_m) \le 1.
\end{equation*}
This bound and \cref{lem:pre-poisson} with $\theta = \lambda'$ imply \eqref{eq:low-added-integral}.
\end{proof}

The following proposition shows that identifying $D$ random vertices advances the hitting time of non-colorability, in expectation, by an amount that is much smaller than $N$ when $D$ is small.

\begin{proposition}\label{prop:low-identify}
Let $N \ge N_q$ be an integer, let $D$ be an integer with $1 \le D \le N/4$, and let $k$ be an integer with $k > e D z_0$. Then
\begin{equation}\label{eq:low-identify}
\E\bigl[H^{(N)} - H_Z\bigr] \le C_q \Bigl(D^k + N \Bigl(\frac{e D z_0}{k}\Bigr)^k + D^2 + 1\Bigr).
\end{equation}
\end{proposition}

For the following figure and the proof, let $Z_0$ be a $D$-subset of $[N]$. We let $(\Pi_s)_{s \ge 0}$ denote a rate-one Poisson counting process independent of $(B_m)$ and $Z$. We set $N' = N - D$, $V' = [N] \setminus Z_0$, and $\lambda = \binom{N}{2}^{-1}$. For $s \ge 0$, we define $B'_s = B_{\Pi_s}[V']$ and $\hat z(s) = \lambda N' s$.

\begin{figure}[t]
\centering
\begin{tikzpicture}[x=0.92cm, y=0.92cm,
  vtx/.style = {circle, fill = black, inner sep = 1.2pt},
  newv/.style = {circle, draw = black, fill = white, thick, inner sep = 1.6pt},
  bg/.style = {gray, thick},
  lst/.style = {black, line width = 0.9pt},
  inz/.style = {black, densely dotted, line width = 1.1pt},
  hull/.style = {gray, dash pattern = on 2pt off 2pt, rounded corners = 7pt},
  lab/.style = {font = \small},
  ttl/.style = {font = \small}]
\def\lowidentbackground{%
  \draw[thick] (0, 0) ellipse (2.05 and 0.95);
  \node[vtx] (p1) at (-1.3, 0.15) {};
  \node[vtx] (p2) at (-0.65, -0.5) {};
  \node[vtx] (p3) at (0, 0.4) {};
  \node[vtx] (p4) at (0.65, -0.45) {};
  \node[vtx] (p5) at (1.35, 0.15) {};
  \draw[bg] (p1) -- (p2);
  \draw[bg] (p2) -- (p4);
  \draw[bg] (p3) -- (p5);
  \draw[bg] (p4) -- (p5);
  \node[lab] at (1.75, -0.95) {$V'$};
  \node[lab, gray] at (-2.15, -0.75) {$B'_s$};}
\begin{scope}
  \lowidentbackground
  \draw[hull] (-1.55, 1.75) rectangle (1.55, 2.55);
  \node[vtx] (z1) at (-1, 2.15) {};
  \node[vtx] (z2) at (0, 2.15) {};
  \node[vtx] (z3) at (1, 2.15) {};
  \node[lab] at (-1.9, 2.15) {$Z_0$};
  \draw[lst] (z1) -- (p1);
  \draw[lst] (z1) -- (p3);
  \draw[lst] (z2) -- (p3);
  \draw[lst] (z2) -- (p4);
  \draw[lst] (z3) to[bend right = 10] (p5);
  \draw[lst] (z3) to[bend left = 10] (p5);
  \draw[inz] (z2) to[bend left = 35] (z3);
  \node[lab] at (0.5, 2.95) {possible occurrence inside $Z_0$};
  \node[ttl] at (0, -1.55) {$B_{\Pi_s}$ near $Z_0$};
\end{scope}
\draw[-{Stealth[length = 2.2mm]}, thick] (2.35, 1.1) -- node[lab, above] {identify}
  node[font = \scriptsize, below, align = center] {no internal\\occurrence} (3.35, 1.1);
\begin{scope}[shift = {(5.7, 0)}]
  \lowidentbackground
  \node[newv, label = {[lab]above:$o$}] (o) at (0, 2.15) {};
  \draw[lst] (o) -- (p1);
  \draw[lst] (o) to[bend right = 9] (p3);
  \draw[lst] (o) to[bend left = 9] (p3);
  \draw[lst] (o) -- (p4);
  \draw[lst] (o) to[bend right = 6] (p5);
  \draw[lst] (o) to[bend left = 6] (p5);
  \node[lab] at (0, -1.55) {$B_{\Pi_s}/Z_0$};
  \node[lab] at (0, -2.15) {list length $\Pois(D \hat z(s))$};
\end{scope}
\draw[-{Stealth[length = 2.2mm]}, thick] (8.05, 1.1) -- node[lab, above] {compare} (9.05, 1.1);
\begin{scope}[shift = {(11.4, 0)}]
  \lowidentbackground
  \node[newv, label = {[lab]above:$o$}] (o) at (0, 2.15) {};
  \draw[lst] (o) -- (p2);
  \draw[lst] (o) -- (p5);
  \node[lab] at (0, -1.55) {$B''$};
  \node[lab] at (0, -2.15) {list length $\Pois(\hat z(s))$};
\end{scope}
\end{tikzpicture}

\caption{The construction in the proof of \cref{prop:low-identify}, at a fixed time $s$, for a $D$-set $Z_0$ with $D = 3$. Left: the multigraph $B'_s$ on $V' = [N] \setminus Z_0$ (gray) and the neighbor lists in $V'$ of the vertices of $Z_0$ (black). The dotted occurrence illustrates the exceptional event bounded in \eqref{eq:low-inside}; if present, it becomes a loop under identification. Middle: the complementary case in which no occurrence has both endpoints in $Z_0$. In this case, the identification $B_{\Pi_s}/Z_0$ is $B'_s$ together with a star at the new vertex $o$, whose list is the concatenation of the neighbor lists. It has $\Pois(D \hat z(s))$ entries, and repeated entries give parallel occurrences. Right: the multigraph $B''$ of \cref{lem:low-added}, built from $B' = B'_s$, in which $o$ has a list of $\Pois(\hat z(s))$ entries. It has the law of the Poissonized process on the $N' + 1$ vertices of $V' \cup \{o\}$. The two lists are compared in \eqref{eq:low-tilt}.}
\label{fig:low-identify}
\end{figure}

\begin{proof}
We first outline the argument, which \cref{fig:low-identify} illustrates. After Poissonization and a bound for large times, we fix a time $s$ and a $D$-set $Z_0$, and we use Poisson splitting at $Z_0$. Unless some occurrence is a pair of $Z_0$, the identification $B_{\Pi_s}/Z_0$ is $B'_s$ together with a star at $o$ whose list has $\Pois(D \hat z(s))$ independent uniform entries in $V'$, independent of $B'_s$. A vertex added as in \cref{lem:low-added} has a list of $\Pois(\hat z(s))$ entries. The main step is the comparison of the two lists by a change of the Poisson mean that is truncated at length $k$, which costs the factor $D^k$; we then apply \cref{lem:low-added}.

We set
\begin{equation}\label{eq:low-prop-params}
s_0 = 2q (\log q) N.
\end{equation}
By \cref{def:pre-poissonized}, $s \mapsto B_{\Pi_s}$ is the Poissonized process on $[N]$. For a multigraph $B$ on $[N]$, we define
\begin{equation*}
h(B) = \Pp(B \text{ is colorable and } B/Z \text{ is not colorable}),
\end{equation*}
where the probability is over $Z$. Since $Z$ is independent of $(B_m)$ and $\Pi$, \eqref{eq:low-advance}, Tonelli's theorem, and \cref{lem:pre-poisson} with $\theta = 1$ give
\begin{equation}\label{eq:low-integral}
\E[H^{(N)} - H_Z] = \sum_{m \ge 0} \E h(B_m) = \int_0^\infty \E h(B_{\Pi_s}) \dd s.
\end{equation}

We begin with large times. Since $N \ge 2q$, every map $\sigma \colon [N] \to [q]$ is constant on at least a fraction $1/(2q)$ of the pairs of $[N]$, as shown in the proof of \cref{lem:pre-lifetime}(a). Hence $\sigma$ is a coloring of $B_m$ with probability at most $(1 - 1/(2q))^m$, and a union bound over the $q^N$ maps gives
\begin{equation*}
\Pp(B_m \text{ is colorable}) \le q^N \Bigl(1 - \frac{1}{2q}\Bigr)^m.
\end{equation*}
Using $h(B) \le \ind\{B \text{ is colorable}\}$ and the $\Pois(s)$ law of $\Pi_s$, we obtain
\begin{equation}\label{eq:low-large}
\int_{s_0}^\infty \E h(B_{\Pi_s}) \dd s \le q^N \int_{s_0}^\infty e^{-s/(2q)} \dd s = 2q \, q^N e^{-s_0/(2q)} = 2q,
\end{equation}
where in the last step we used $s_0 = 2q (\log q) N$.

Next, we fix $s \in [0, s_0]$ and a $D$-subset $Z_0$ of $[N]$ and bound the probability that $B_{\Pi_s}$ is colorable and $B_{\Pi_s}/Z_0$ is not. Since $Z$ is uniform and independent of $(B_m)$ and $\Pi$, a bound that holds for every $Z_0$ also holds for $\E h(B_{\Pi_s})$.
We apply \cref{lem:pre-split}(b) with $S = [N]$ and $S_0 = Z_0$. It shows that the following are independent: the multigraph $B'_s$, which has the law of the Poissonized process on $V'$ at time $\lambda \binom{N'}{2} s$; for each $v \in Z_0$, the neighbor list $\mathbf x_v$ of $v$ in $V'$ at time $s$, which is a list of $\Pois(\hat z(s))$ independent uniform elements of $V'$; and the multigraph $B_{\Pi_s}[Z_0]$.

Let $\mathbf x_{Z_0}$ be the concatenation of the lists $\mathbf x_v$ in increasing order of $v \in Z_0$. Conditional on the lengths of all the lists, the entries of $\mathbf x_{Z_0}$ are independent and uniform on $V'$. It follows that $\mathbf x_{Z_0}$ is a list of $\Pois(D \hat z(s))$ independent uniform elements of $V'$, independent of $B'_s$.

Suppose that $B_{\Pi_s}[Z_0]$ is empty. Then every occurrence of $B_{\Pi_s}$ is either an occurrence of $B'_s$ or of the form $\{v, x\}$ with $v \in Z_0$ and $x \in V'$, and the identification replaces each occurrence of the second kind by $\{o, x\}$. Therefore
\begin{equation}\label{eq:low-quotient}
B_{\Pi_s}/Z_0 = B'_s + \operatorname{star}(o, \mathbf x_{Z_0})
\end{equation}
as multigraphs on $V' \cup \{o\}$, with all multiplicities. The multigraph $B'_s$ is colorable whenever $B_{\Pi_s}$ is colorable, and $B_{\Pi_s}[Z_0]$ has $\lambda \binom{D}{2} s$ occurrences in expectation. Then \eqref{eq:low-quotient} and Markov's inequality imply that
\begin{equation}\label{eq:low-split}
\begin{aligned}
&\Pp\bigl(B_{\Pi_s} \text{ is colorable and } B_{\Pi_s}/Z_0 \text{ is not colorable}\bigr) \\
&\le \lambda \binom{D}{2} s \\
&\quad + \Pp\bigl(B'_s \text{ is colorable and } B'_s + \operatorname{star}(o, \mathbf x_{Z_0}) \text{ is not colorable}\bigr).
\end{aligned}
\end{equation}
Integrating the first term over $[0, s_0]$, we find
\begin{equation}\label{eq:low-inside}
\int_0^{s_0} \lambda \binom{D}{2} s \dd s = \frac{D(D - 1) s_0^2}{2N(N - 1)} = 2q^2 (\log q)^2 \frac{D(D - 1) N}{N - 1} \le 2q^2 (\log q)^2 D^2,
\end{equation}
where in the last step we used $D \le N$.

We now compare the list $\mathbf x_{Z_0}$, whose length has mean $D\hat z(s)$, with a list whose length has mean $\hat z(s)$. For this comparison, we abbreviate $\hat z = \hat z(s)$. Using $\lambda N' \le 2/(N - 1)$ and $N - 1 \ge 4q \log q$, we find
\begin{equation}\label{eq:low-zeta}
\hat z \le \frac{2 s_0}{N - 1} = 4q \log q + \frac{4q \log q}{N - 1} \le z_0.
\end{equation}

Let $\mathbf x = (\mathbf x(1), \mathbf x(2), \dots)$ be a sequence of independent uniform elements of $V'$, and let $J$ and $J'$ have the laws $\Pois(\hat z)$ and $\Pois(D \hat z)$, with $B'_s$, $\mathbf x$, $J$, and $J'$ independent. By the previous paragraph, the pair $(B'_s, \mathbf x_{Z_0})$ has the same law as $(B'_s, \mathbf x_{[J']})$. Let $I_\star$ be defined as in \cref{lem:low-added} with $B' = B'_s$.

Non-colorability persists as $j$ increases, so on the event that $B'_s$ is colorable, the multigraph $B'_s + \operatorname{star}(o, \mathbf x_{[j]})$ is not colorable if and only if $I_\star \le j$. Using the independence of $J'$ from $(B'_s, \mathbf x)$, we obtain
\begin{equation*}
\begin{aligned}
&\Pp\bigl(B'_s \text{ is colorable and } B'_s + \operatorname{star}(o, \mathbf x_{Z_0}) \text{ is not colorable}\bigr) \\
&= \sum_{j \ge 0} \pi_{D \hat z}(j) \Pp(B'_s \text{ is colorable}, \ I_\star \le j).
\end{aligned}
\end{equation*}

For $0 \le j \le k$, we have
\begin{equation*}
\pi_{D \hat z}(j) = D^j e^{-(D - 1) \hat z} \pi_{\hat z}(j) \le D^k \pi_{\hat z}(j).
\end{equation*}
We use this bound for $j \le k$, and we bound the probabilities by one for $j > k$. Since $J$ is independent of $(B'_s, \mathbf x)$, this gives
\begin{equation}\label{eq:low-tilt}
\sum_{j \ge 0} \pi_{D \hat z}(j) \Pp(B'_s \text{ is colorable}, \ I_\star \le j) \le D^k \Pp(B'_s \text{ is colorable}, \ I_\star \le J) + \Pp(J' > k).
\end{equation}
Further, $k > e D z_0 \ge D \hat z$, and \cref{lem:pre-tails}(b) implies that
\begin{equation*}
\Pp(J' > k) \le \Bigl(\frac{e D \hat z}{k}\Bigr)^k \le \Bigl(\frac{e D z_0}{k}\Bigr)^k.
\end{equation*}
Integrating over $[0, s_0]$, we have
\begin{equation}\label{eq:low-tail}
\int_0^{s_0} \Pp(J' > k) \dd s \le 2q (\log q) N \Bigl(\frac{e D z_0}{k}\Bigr)^k.
\end{equation}

Next, we bound the first term on the right side of \eqref{eq:low-tilt} with \cref{lem:low-added}. Its hypothesis on $B'_s$ holds by \cref{lem:pre-split}(a), applied to the Poissonized process on $V'$ at time $\lambda \binom{N'}{2} s$ with the partition of the pairs of $V'$ into single pairs. Let $\lambda' = \lambda \binom{N' + 1}{2}$ and let $F(t)$ be as in \cref{lem:low-added}. Using \eqref{eq:low-added}, \eqref{eq:low-zeta}, and \eqref{eq:low-added-integral}, we obtain
\begin{equation}\label{eq:low-main-term}
\begin{aligned}
&\int_0^{s_0} D^k \Pp(B'_s \text{ is colorable}, \ I_\star \le J) \dd s \\
&\le D^k (e^{z_0} - 1) \int_0^\infty \E \kappa\bigl(F(\lambda' s)\bigr) \dd s \\
&\le \frac{D^k (e^{z_0} - 1)}{\lambda'} \\
&\le 2 (e^{z_0} - 1) D^k.
\end{aligned}
\end{equation}
In the last step, we used $N - D \ge 3N/4$, which gives
\begin{equation*}
\frac{1}{\lambda'} = \frac{N(N - 1)}{(N - D + 1)(N - D)} \le \frac{N^2}{(3N/4)^2} = \frac{16}{9}.
\end{equation*}

Finally, we collect the bounds. We split the integral in \eqref{eq:low-integral} at $s_0$, and we use \eqref{eq:low-large} for the integral over $[s_0, \infty)$. For the integral over $[0, s_0]$, we use \eqref{eq:low-split} and \eqref{eq:low-tilt}, and we integrate with \eqref{eq:low-inside}, \eqref{eq:low-tail}, and \eqref{eq:low-main-term}. We obtain
\begin{equation*}
\E[H^{(N)} - H_Z] \le 2q + 2q^2 (\log q)^2 D^2 + 2q (\log q) N \Bigl(\frac{e D z_0}{k}\Bigr)^k + 2 (e^{z_0} - 1) D^k.
\end{equation*}
Each of the four coefficients is bounded by $C_q = 2 e^{z_0} = 2e q^{4q}$, which gives \eqref{eq:low-identify}.
\end{proof}

\section{Proof of the low-degree estimate}\label{sec:low-main}

We now prove the main result of this section. Recall that $D_n = \lfloor (\log n)^{1/4} \rfloor$, and that $r(G)$ is the expected lifetime of a multigraph $G$ on $[n]$, defined in \cref{def:pre-lifetime}.

\begin{theorem}\label{thm:low-main}
There exist $C_{\rm low}$ and $n_q$, depending only on $q$, such that for all $n \ge n_q$ and all $y > 0$,
\begin{equation}\label{eq:low-main}
\Pp\Bigl(H < \infty, \ \min_{v \in e_H} \deg_\Gamma(v) \le D_n, \ r(\Gamma) > y\Bigr) \le C_{\rm low} y^{-1/2} n^{3/8}.
\end{equation}
\end{theorem}

\begin{proof}
We first fix the parameters. We set
\begin{equation*}
k_n = \Bigl\lceil \frac{\log n}{\log \log n} \Bigr\rceil, \qquad C_{\rm low} = 8 C_q (1 + A_q)^{1/2},
\end{equation*}
and we let $n_q$ be the least integer such that every $n \ge n_q$ satisfies $n - 1 \ge N_q$ and
\begin{equation}\label{eq:low-conditions}
\log\bigl(e z_0 \log \log n\bigr) \le \frac{1}{8} \log \log n, \qquad n^{1/4} (\log n)^{1/4} + (\log n)^{1/2} + 1 \le n^{3/8}.
\end{equation}
There exists such an integer, depending only on $q$, because $\log(e z_0 \log \log n)/\log \log n \to 0$ and $n^{1/4} (\log n)^{1/4} + (\log n)^{1/2} + 1 = o(n^{3/8})$ as $n \to \infty$. Let $n \ge n_q$ and $y > 0$.

Since $n - 1 \ge N_q > 2q$ and $N_q \ge 15$, we have $n \ge 2q$, $\log \log n > 0$, and $1 \le D_n \le (n - 1)/4$.

We begin by writing the probability in \eqref{eq:low-main} as a time integral. We define the set $\cE_{\rm low}$ and the function $h$ on multigraphs on $[n]$ by
\begin{equation*}
\begin{aligned}
&\cE_{\rm low} = \Bigl\{(G, e) : G \text{ is colorable, } e \text{ is a frozen pair of } G, \ \min_{v \in e} \deg_G(v) \le D_n\Bigr\}, \\
&h(G) = \ind\{r(G) > y\}.
\end{aligned}
\end{equation*}
Recall from \cref{lem:pre-frozen}(a) that $e_H$ is a frozen pair of $\Gamma$ whenever $H < \infty$. Hence the event in \eqref{eq:low-main} is the event that $H < \infty$, $(\Gamma, e_H) \in \cE_{\rm low}$, and $r(\Gamma) > y$. \Cref{lem:pre-terminal}(a) with $L' = 0$ and $U' = \infty$, followed by \cref{lem:pre-poisson} with $\theta = 1$, shows that its probability is
\begin{equation}\label{eq:low-step1}
\sum_{m \ge 0} \E\bigl[\kappa_{\cE_{\rm low}}(G_m) h(G_m)\bigr] = \int_0^\infty \E\bigl[\kappa_{\cE_{\rm low}}(G(t)) h(G(t))\bigr] \dd t.
\end{equation}

Next, we reduce to a single vertex. Let $G$ be a colorable multigraph on $[n]$. Every frozen pair $e$ with $(G, e) \in \cE_{\rm low}$ has an endpoint $a$ with $\deg_G(a) \le D_n$, whose frozen class $C_G(a)$ contains the other endpoint. Counting each such pair at least once through such an endpoint, we obtain
\begin{equation*}
\kappa_{\cE_{\rm low}}(G) \le \binom{n}{2}^{-1} \sum_{a \in [n]} \ind\{G \text{ is colorable}, \ \deg_G(a) \le D_n\} \bigl(|C_G(a)| - 1\bigr),
\end{equation*}
and the two sides vanish if $G$ is not colorable.

The law of $G(t)$ is invariant under relabeling, as a mixture of the laws of the multigraphs $G_m$. A permutation of $[n]$ that maps $a$ to $a'$ maps $C_G(a)$ onto the frozen class of $a'$ in the relabeled multigraph and preserves degrees, and \cref{lem:pre-lifetime}(b) shows that it preserves $h$. Therefore the $n$ terms of the sum have the same expectation at $G = G(t)$. We set $a_0 = n$. We find
\begin{equation}\label{eq:low-root}
\begin{aligned}
&\E\bigl[\kappa_{\cE_{\rm low}}(G(t)) h(G(t))\bigr] \\
&\le \frac{2}{n - 1} \E\bigl[\ind\{G(t) \text{ is colorable}, \ \deg_{G(t)}(a_0) \le D_n\} \bigl(|C_{G(t)}(a_0)| - 1\bigr) h(G(t))\bigr].
\end{aligned}
\end{equation}

We now use the restriction on the lifetime. By \Cref{lem:pre-classmass}, every colorable multigraph $G$ on $[n]$ with $r(G) > y$ satisfies $|C_G(a_0)| \le n ((1 + A_q)/y)^{1/2}$. Further, $|C_G(a_0)| - 1 > 0$ if and only if $a_0$ lies in a frozen pair of $G$. Inserting these facts into \eqref{eq:low-root}, and using $h \le 1$ afterwards, we obtain
\begin{equation}\label{eq:low-classmass}
\begin{aligned}
&\E\bigl[\kappa_{\cE_{\rm low}}(G(t)) h(G(t))\bigr] \\
&\le \frac{2n}{n - 1} \Bigl(\frac{1 + A_q}{y}\Bigr)^{1/2} \Pp\bigl(\deg_{G(t)}(a_0) \le D_n \text{ and } a_0 \text{ lies in a frozen pair of } G(t)\bigr).
\end{aligned}
\end{equation}
This is the only use of the lifetime in the proof.

Next, we use \cref{lem:low-frozen}(iii). This is the step where $q \ge 3$ is used. Let $B(t) = G(t)[[n - 1]]$ be the multigraph of the occurrences of $G(t)$ that do not contain $a_0$, and let $\mathbf x$ be the neighbor list of $a_0$ in $[n - 1]$ at time $t$. Then
\begin{equation*}
G(t) = B(t) + \operatorname{star}(a_0, \mathbf x),
\end{equation*}
and $\deg_{G(t)}(a_0)$ is the length of $\mathbf x$.

We set $p = (n - 2)/n$ and $z(t) = 2t/n$. We apply \cref{lem:pre-split}(c) with $S = [n]$ and $a = a_0$. It shows that $B(t)$ and $\mathbf x$ are independent, that $B(t)$ has the law of the Poissonized process on $[n - 1]$ at time $pt$, and that $\mathbf x$ is a list of $\Pois(z(t))$ independent uniform elements of $[n - 1]$.

Suppose that $a_0$ lies in a frozen pair of $G(t)$. Then \cref{lem:low-frozen}(iii) and $q \ge 3$ show that $B(t)$ is colorable and that $\mathbf x$ is never monochromatic in $B(t)$.

For an integer $D \ge 0$ and a multigraph $B$ on $[n - 1]$, we let $\mathrm{nm}_D(B)$ denote the probability that a list of $D$ independent uniform elements of $[n - 1]$ is never monochromatic in $B$, with $\mathrm{nm}_D(B) = 0$ if $B$ is not colorable. Since every map is constant on the empty list, we have $\mathrm{nm}_0 = 0$.

A coloring that is constant on a list is constant on each prefix of the list, so a list with a prefix that is never monochromatic in $B$ is itself never monochromatic in $B$. Realizing a list of length $D \le D_n$ as the prefix of a list of length $D_n$, we find $\mathrm{nm}_D(B) \le \mathrm{nm}_{D_n}(B)$.

Conditioning on the length of $\mathbf x$ and on $B(t)$, and using $\sum_{D} \pi_{z(t)}(D) \le 1$, we obtain
\begin{equation}\label{eq:low-nm}
\begin{aligned}
&\Pp\bigl(\deg_{G(t)}(a_0) \le D_n \text{ and } a_0 \text{ lies in a frozen pair of } G(t)\bigr) \\
&\le \sum_{D = 0}^{D_n} \pi_{z(t)}(D) \E\bigl[\mathrm{nm}_D(B(t))\bigr] \\
&\le \E\bigl[\mathrm{nm}_{D_n}(B(t))\bigr].
\end{aligned}
\end{equation}

We now pass to identifications. Let $(B_m)_{m \ge 0}$ be the time-$m$ graphs of an edge process on $[n - 1]$. As in \cref{sec:low-identify}, we let $Z$ be a uniform random $D_n$-subset of $[n - 1]$, independent of $(B_m)$. We claim that for every colorable multigraph $F$ on $[n - 1]$,
\begin{equation}\label{eq:low-completion}
\mathrm{nm}_{D_n}(F) \le \Pp(F/Z \text{ is not colorable}).
\end{equation}

To prove \eqref{eq:low-completion}, let $\mathbf x'$ be a list of $D_n$ independent uniform elements of $[n - 1]$, and let $Z_1$ denote its set of entries. Given $\mathbf x'$, let $Z'$ be the union of $Z_1$ and a uniform random subset of $[n - 1] \setminus Z_1$ with $D_n - |Z_1|$ elements, chosen independently of everything else. The joint law of $\mathbf x'$ and $Z'$ is invariant under relabeling of $[n - 1]$, so the law of $Z'$ is invariant under all permutations of $[n - 1]$. The permutations act transitively on the $D_n$-subsets of $[n - 1]$, and hence $Z'$ is a uniform $D_n$-subset.

Now suppose that $\mathbf x'$ is never monochromatic in $F$. Then no coloring of $F$ is constant on $Z' \supset Z_1$, and \cref{lem:low-frozen}(iv) shows that $F/Z'$ is not colorable. This proves \eqref{eq:low-completion}.

Applying \cref{lem:pre-poisson} on $[n - 1]$ with the factor $p$, followed by \eqref{eq:low-completion} and \eqref{eq:low-advance} with $N = n - 1$ and $D = D_n$, we find
\begin{equation}\label{eq:low-to-identify}
\begin{aligned}
&\int_0^\infty \E\bigl[\mathrm{nm}_{D_n}(B(t))\bigr] \dd t \\
&= \frac{1}{p} \sum_{m \ge 0} \E\bigl[\mathrm{nm}_{D_n}(B_m)\bigr] \\
&\le \frac{1}{p} \sum_{m \ge 0} \Pp\bigl(B_m \text{ is colorable and } B_m/Z \text{ is not colorable}\bigr) \\
&= \frac{1}{p} \E\bigl[H^{(n - 1)} - H_Z\bigr].
\end{aligned}
\end{equation}

To bound the right side of \eqref{eq:low-to-identify}, we apply \cref{prop:low-identify} with $N = n - 1 \ge N_q$, $D = D_n$, and $k = k_n$. Using $\log D_n \le (1/4) \log \log n$ and $k_n \le \log n/\log \log n + 1$, we find
\begin{equation*}
D_n^{k_n} \le \exp\Bigl(\Bigl(\frac{\log n}{\log \log n} + 1\Bigr) \frac{\log \log n}{4}\Bigr) = n^{1/4} (\log n)^{1/4}.
\end{equation*}
Further, using $D_n \le (\log n)^{1/4}$, $k_n \ge \log n/\log \log n$, and the first condition in \eqref{eq:low-conditions}, we find
\begin{equation*}
\begin{aligned}
&\frac{e D_n z_0}{k_n} \\
&\le e z_0 (\log \log n) (\log n)^{-3/4} \\
&= \exp\Bigl(\log\bigl(e z_0 \log \log n\bigr) - \frac{3}{4} \log \log n\Bigr) \\
&\le \exp\Bigl(-\frac{5}{8} \log \log n\Bigr) < 1.
\end{aligned}
\end{equation*}
In particular, $k_n > e D_n z_0$, and
\begin{equation*}
(n - 1) \Bigl(\frac{e D_n z_0}{k_n}\Bigr)^{k_n} \le n \exp\Bigl(-\frac{5}{8} k_n \log \log n\Bigr) \le n \exp\Bigl(-\frac{5}{8} \log n\Bigr) = n^{3/8}.
\end{equation*}
Inserting these bounds and $D_n^2 \le (\log n)^{1/2}$ into \eqref{eq:low-identify}, we obtain
\begin{equation}\label{eq:low-advance-bound}
\E\bigl[H^{(n - 1)} - H_Z\bigr] \le C_q \bigl(n^{1/4} (\log n)^{1/4} + n^{3/8} + (\log n)^{1/2} + 1\bigr) \le 2 C_q n^{3/8},
\end{equation}
where in the last step we used the second condition in \eqref{eq:low-conditions}.

Finally, we combine \eqref{eq:low-step1}, \eqref{eq:low-classmass}, \eqref{eq:low-nm}, \eqref{eq:low-to-identify}, and \eqref{eq:low-advance-bound}. With $p^{-1} = n/(n - 2)$, this gives
\begin{equation*}
\begin{aligned}
&\Pp\Bigl(H < \infty, \ \min_{v \in e_H} \deg_\Gamma(v) \le D_n, \ r(\Gamma) > y\Bigr) \\
&\le \frac{2n^2}{(n - 1)(n - 2)} \Bigl(\frac{1 + A_q}{y}\Bigr)^{1/2} \cdot 2 C_q n^{3/8} \\
&\le 8 C_q (1 + A_q)^{1/2} y^{-1/2} n^{3/8},
\end{aligned}
\end{equation*}
where in the last step we used $2n^2 \le 4(n - 1)(n - 2)$, which holds for $n \ge 6$. This completes the proof.
\end{proof}

With $y = x_n$, we obtain the first estimate required in \cref{sec:red-split}.

\begin{corollary}\label{cor:low}
For all large $n$, we have
\begin{equation}\label{eq:low-cor}
\Pp\Bigl(H < \infty, \ \min_{v \in e_H} \deg_\Gamma(v) \le D_n, \ r(\Gamma) > x_n\Bigr) \le \eps_n.
\end{equation}
In particular, for every $\eta \in (0, \eta_0]$ and all large $n$, the hypothesis \eqref{eq:red-low} of \cref{prop:red-split} holds.
\end{corollary}

\begin{proof}
Let $C_{\rm low}$ and $n_q$ be as in \cref{thm:low-main}, and recall that $x_n = n (\log n)^{-3}$ for $n \ge 3$. For $n \ge n_q$, \cref{thm:low-main} with $y = x_n$ bounds the left side of \eqref{eq:low-cor} by
\begin{equation*}
C_{\rm low} x_n^{-1/2} n^{3/8} = C_{\rm low} (\log n)^{3/2} n^{-1/8}.
\end{equation*}
This is bounded by $(\log n)^{-3/2} = \eps_n$ as soon as $C_{\rm low} (\log n)^3 \le n^{1/8}$, which holds for all large $n$. The last statement follows from the fact that the event in \eqref{eq:red-low} is contained in the event in \eqref{eq:low-cor}.
\end{proof}

\begin{remark}\label{rem:low-limit}
The exponent $1/4$ in $D_n$ is not essential, but \cref{prop:low-identify} does not reach degrees of order $(\log n)^{1/2}$. Let $b_0 \in (0, 1)$ and $b_1 > 0$, and let $D = \lfloor (\log N)^{b_0} \rfloor$ and $k = \lceil b_1 \log N/\log \log N \rceil$. Then, as $N \to \infty$,
\begin{equation*}
D^k = N^{b_0 b_1 + o(1)}, \qquad N \Bigl(\frac{e D z_0}{k}\Bigr)^k = N^{1 - b_1 (1 - b_0) + o(1)}.
\end{equation*}
The two exponents are equal at $b_1 = 1$, where both equal $b_0$, and for every other $b_1$ the larger of them exceeds $b_0$. With $b_1 = 1$, \cref{prop:low-identify} gives $\E[H^{(N)} - H_Z] \le N^{b_0 + o(1)}$. The proof of \cref{thm:low-main} with $D_n$ replaced by $\lfloor (\log n)^{b_0} \rfloor$ then bounds the probability in \eqref{eq:low-cor} by $n^{b_0 - 1/2 + o(1)}$. For $b_0 < 1/2$, this is at most $\eps_n$ for all large $n$.

At $b_0 = 1/2$, the two terms in \cref{prop:low-identify} cannot both be $o(\sqrt N)$, regardless of the admissible choice of $k$. Indeed, set $D = \lfloor \sqrt{\log N} \rfloor$ and suppose that $D^k \le \sqrt N$. Then $k \le (\log N)/(2 \log D)$. The expression $k \log(k/(e D z_0))$ is increasing in $k > e D z_0$, and hence
\[
\begin{split}
k \log\Bigl(\frac{k}{e D z_0}\Bigr)
&\le \frac{\log N}{2 \log D}
\log\Bigl(\frac{\log N}{2 e D z_0 \log D}\Bigr) \\
&= \frac{\log N}{2}
- \frac{\log N}{2 \log D}
\log\Bigl(\frac{2 e z_0 D^2 \log D}{\log N}\Bigr)
< \frac{\log N}{2}
\end{split}
\]
for all large $N$, since $2 e z_0 D^2 \log D/\log N \sim e z_0 \log\log N \to \infty$. Thus $N(e D z_0/k)^k > \sqrt N$. If instead $D^k > \sqrt N$, the first term already exceeds $\sqrt N$. This proves the stated limitation of the bound at $b_0 = 1/2$. Vertices of degree larger than $D_n$ require a different argument.
\end{remark}

\chapter{Final edges at vertices of high degree}\label{sec:high}

This section proves \cref{thm:high-main}, the estimate for final edges at vertices of degree greater than $D_n$ required by \cref{prop:red-split}. It bounds the probability that the process becomes non-colorable by time $M$ through a final edge $e_H$ with an endpoint of degree greater than $D_n$ in $\Gamma$. The bound requires no restriction on the expected lifetime $r(\Gamma)$. Since it concerns all final edges up to time $M$, it also bounds, for all integers $0 \le L' \le U' \le M$, the probability that $L' < H \le U'$ and both endpoints of $e_H$ have degree greater than $D_n$ in $\Gamma$. The heuristics behind the argument are explained in \cref{sec:overview}. The arguments use the number of colors only through the notion of colorability, and they hold for $q = 2$ as well (see \cref{rem:high-q2}).

The argument has three steps. In \cref{sec:high-pivotal}, we delete edge occurrences and record whether two specified vertices must still receive the same color in every coloring. Iterating the estimate for deletions that preserve this equality reduces the high-degree probability to the density of pairs for which many incident deletions destroy the equality (\cref{prop:high-descent}). We express this density through a relation, defined below, that records possible additions to the frozen class of a vertex (\cref{lem:high-deletion}). In \cref{sec:high-relation}, we bound the density of this relation by the expected value of a logarithmic increment defined below (\cref{lem:high-onestep}). This increment is also defined when the added edge makes the graph non-colorable. In \cref{sec:high-potential}, we compare these increments with the growth of the logarithm of the frozen-class size before the last step and bound the last increment separately. Their sum is bounded by $\log n$, which controls the expected sum of the relation densities after a stopping time (\cref{thm:high-potential}). Finally, \cref{sec:high-main} combines this estimate with a binomial tail bound for the probability that a fixed vertex reaches the required degree by time $M$.

Throughout this section, $n \ge 3$, and $\Lambda > 0$ and $M = \lfloor \Lambda n \rfloor$ are as in \cref{sec:pre-process}; the results of this section use only $\Lambda > 0$. Recall the degree cut $D_n = \lfloor (\log n)^{1/4} \rfloor$ from \eqref{eq:red-parameters}. We set
\begin{equation}\label{eq:high-params}
K = \lceil 8 \Lambda \rceil, \qquad w_0 = \lfloor D_n / 2 \rfloor.
\end{equation}
Since $n \ge 3$, we have $D_n \ge 1$ and therefore $0 \le w_0 < D_n$. For distinct vertices $a$ and $b$, we let $ab$ denote the pair $\{a, b\}$. Recall that the degree $\deg_G(u)$ counts occurrences with multiplicity, that the relation $u \sim_G v$ includes the requirement that $G$ is colorable, and that every quantity defined through colorings or frozen classes is set to zero on non-colorable multigraphs.

\section{Pivotal occurrences}\label{sec:high-pivotal}

We first define the two kinds of occurrences at an endpoint of a frozen pair.

\begin{definition}\label{def:high-pivotal}
Let $G$ be a colorable multigraph on $[n]$, and let $(u, v)$ be an ordered pair of distinct vertices such that $u \sim_G v$. We call $(u, v)$ an \emph{ordered frozen pair} of $G$. Let $e$ be an occurrence of $G$ with $u \in e$. We say that $e$ is \emph{non-pivotal} for $(u, v)$ if $u \sim_{G - e} v$, and that $e$ is \emph{pivotal} for $(u, v)$ otherwise. For an integer $j \ge 0$, we define
\begin{equation}\label{eq:high-kappaj}
\kappa_j(G) = \frac{1}{n(n - 1)} \, \# \bigl\{ (u, v) \in [n]^2 : u \ne v, \ u \sim_G v, \ \deg_G(u) = j \bigr\},
\end{equation}
and we define $\kappa_j^{\rm np}(G)$ by the same formula, with the additional requirement that at least $K$ of the $j$ occurrences at $u$ are non-pivotal for $(u, v)$. We set $\kappa_j^{\rm piv}(G) = \kappa_j(G) - \kappa_j^{\rm np}(G)$. If $G$ is not colorable, we set $\kappa_j(G) = \kappa_j^{\rm np}(G) = \kappa_j^{\rm piv}(G) = 0$.
\end{definition}

Occurrences are counted with multiplicity, so that exactly $\deg_G(u)$ occurrences contain $u$. Note that $G - e \subset G$ is colorable. Moreover, $G - e$ is determined by the pair $e$ and not by the occurrence that is removed. All occurrences of a pair are pivotal or all are non-pivotal. An occurrence of a pair of multiplicity at least two is non-pivotal, since $G - e$ then has the same colorings as $G$.

A pair counted in $n(n - 1) \kappa_j^{\rm piv}(G)$ has at least $j - K + 1$ pivotal occurrences at $u$. By \cref{lem:pre-frozen}(a), applied to the colorable multigraph $G - e$, the occurrence $e$ is non-pivotal for $(u, v)$ if and only if $(G - e) + uv$ is not colorable. Thus $e$ is pivotal for $(u, v)$ exactly when it is pivotal, in the usual sense \cite[Section~2, p.~130]{russo1982approximate}, for the increasing event that $G + uv$ is not colorable, in the sense that $G$ belongs to this event and $G - e$ does not (see \cite[Section~4.8, Exercise~4.9]{grimmett2018probability}).

We will also need a relation that records how the frozen class of a vertex grows when an edge at that vertex is added.

\begin{definition}\label{def:high-relation}
Let $G$ be a colorable multigraph on $[n]$, let $u \in [n]$, and let $W = [n] \setminus C_G(u)$. We define $R_G(u)$ to be the set of ordered pairs $(y, v) \in W \times W$ such that
\begin{equation}\label{eq:high-R}
\sigma(u) \in \{ \sigma(y), \sigma(v) \} \qquad \text{for all } \sigma \in \Col(G).
\end{equation}
For $y \in [n]$, we define the \emph{section} $R_G(u; y) = \{ v \in [n] : (y, v) \in R_G(u) \}$, and we set
\begin{equation}\label{eq:high-beta}
\beta_G(u) = n^{-2} \abs{R_G(u)}.
\end{equation}
If $G$ is not colorable, we set $R_G(u) = \emptyset$ and $\beta_G(u) = 0$.
\end{definition}

Equivalently, $\beta_G(u)$ is the probability that two independent uniform vertices $y$ and $v$ of $[n]$ lie outside $C_G(u)$ and that every coloring of $G$ gives $u$ the color of $y$ or the color of $v$. The section $R_G(u; y)$ is empty for $y \in C_G(u)$. When $G$ and $u$ are fixed, we let $R(y)$ denote $R_G(u; y)$. In this notation, we let $W$ denote $[n] \setminus C_G(u)$.

\begin{lemma}\label{lem:high-relation}
Let $G$ be a colorable multigraph on $[n]$ and let $u \in [n]$.
\begin{enumerate}[label=(\alph*)]
\item The relation $R_G(u)$ is symmetric, and $(y, y) \notin R_G(u)$ for all $y \in [n]$.
\item For all $y \in W$, the multigraph $G + uy$ is colorable and $C_{G + uy}(u) = C_G(u) \cup R(y)$.
\end{enumerate}
\end{lemma}

\begin{proof}
For (a), note that the condition \eqref{eq:high-R} is symmetric in $y$ and $v$. For $v = y$, it states that $\sigma(u) = \sigma(y)$ for all colorings $\sigma$ of $G$, which means that $y \in C_G(u)$ and is excluded by the requirement $y \in W$.

For (b), let $y \in W$. We have $y \ne u$ and $u \not \sim_G y$, and \cref{lem:pre-frozen}(a) implies that $G + uy$ is colorable. The colorings of $G + uy$ are the colorings $\sigma$ of $G$ with $\sigma(y) \ne \sigma(u)$. Therefore $v \in C_{G + uy}(u)$ if and only if every coloring $\sigma$ of $G$ with $\sigma(y) \ne \sigma(u)$ satisfies $\sigma(v) = \sigma(u)$, which gives
\[
v \in C_{G + uy}(u) \iff \sigma(u) \in \{ \sigma(y), \sigma(v) \} \quad \text{for all } \sigma \in \Col(G).
\]
The right side holds for all $v \in C_G(u)$, and for $v \in W$ it is the condition $(y, v) \in R_G(u)$.
\end{proof}

By \cref{lem:high-relation}(b), for $y \in W$, the number $\abs{R(y)}$ counts the vertices that join the frozen class of $u$ when the pair $uy$ is added to $G$. Since $\abs{R_G(u)} = \sum_{y \in [n]} \abs{R(y)}$, the quantity $n \beta_G(u)$ is the average of $\abs{R(y)}$ over $y \in [n]$.

We now bound the densities $\kappa_j^{\rm np}$ and $\kappa_j^{\rm piv}$ of \cref{def:high-pivotal} in terms of graphs with one occurrence fewer. To that end, we remove one of the $m$ occurrences $e_1, \dots, e_m$ of $G_m$. For each $i \le m$, the multigraph $G_m - e_i$ is formed by the $m - 1$ independent uniform pairs $e_{i'}$ with $i' \ne i$, which are also independent of $e_i$. The pair $(G_m - e_i, e_i)$ therefore has the law of $(G_{m - 1}, e)$, where $e$ is a uniform pair independent of $G_{m - 1}$, and for every function $h \ge 0$ of a multigraph and a pair we have
\begin{equation}\label{eq:high-exchange}
\E \sum_{i = 1}^m h(G_m - e_i, e_i) = m \E h(G_{m - 1}, e).
\end{equation}
This is the analogue for $m$ independent points of the Mecke equation for Poisson processes \cite{mecke1967stationare}, \cite[Theorem~4.1]{last2018lectures}.

\begin{lemma}\label{lem:high-deletion}
Let $m \ge 1$ and $j \ge 1$ be integers.
\begin{enumerate}[label=(\alph*)]
\item We have
\begin{equation}\label{eq:high-deletion-np}
\E \kappa_j^{\rm np}(G_m) \le \frac{2m}{nK} \, \E \kappa_{j - 1}(G_{m - 1}).
\end{equation}
\item If $j \ge K$, then for all $u \in [n]$,
\begin{equation}\label{eq:high-deletion-piv}
\E \kappa_j^{\rm piv}(G_m) \le \frac{2mn}{(n - 1)^2 (j - K + 1)} \, \E \bigl[ \ind \{ \deg_{G_{m - 1}}(u) = j - 1 \} \, \beta_{G_{m - 1}}(u) \bigr].
\end{equation}
\end{enumerate}
\end{lemma}

\begin{proof}
We begin with (a). For a multigraph $G'$ and a pair $e$, we define $h_1(G', e)$ by
\[
h_1(G', e) = \sum_{a \in e} \sum_{v \in [n] \setminus \{a\}} \ind \bigl\{ a \sim_{G' + e} v, \ a \sim_{G'} v, \ \deg_{G' + e}(a) = j \bigr\}.
\]
Let $(a, v)$ be a pair counted in $n(n - 1) \kappa_j^{\rm np}(G_m)$. Then $a \sim_{G_m} v$, $\deg_{G_m}(a) = j$, and there exist at least $K$ indices $i \le m$ such that $a \in e_i$ and $a \sim_{G_m - e_i} v$. Each such index contributes $1$ to $h_1(G_m - e_i, e_i)$, through the endpoint $a$ and the vertex $v$. Summing over the pairs $(a, v)$ and taking expectations, we obtain
\begin{equation}\label{eq:high-deletion-1}
K n (n - 1) \E \kappa_j^{\rm np}(G_m) \le \E \sum_{i = 1}^m h_1(G_m - e_i, e_i) = m \E h_1(G_{m - 1}, e),
\end{equation}
where the identity is \eqref{eq:high-exchange}.

We now bound the right side of \eqref{eq:high-deletion-1}. For $a \in e$, we have $\deg_{G_{m - 1} + e}(a) = \deg_{G_{m - 1}}(a) + 1$. Recall from \cref{sec:pre-process} that each vertex belongs to $e$ with probability $2/n$, and note that $e$ is independent of $G_{m - 1}$. Dropping the first indicator in the definition of $h_1$, we find
\[
\begin{split}
\E h_1(G_{m - 1}, e) &\le \E \sum_{a \in e} \sum_{v \in [n] \setminus \{a\}} \ind \bigl\{ a \sim_{G_{m - 1}} v, \ \deg_{G_{m - 1}}(a) = j - 1 \bigr\}\\
&= \frac{2}{n} \sum_{a \in [n]} \sum_{v \in [n] \setminus \{a\}} \Pp \bigl( a \sim_{G_{m - 1}} v, \ \deg_{G_{m - 1}}(a) = j - 1 \bigr)\\
&= \frac{2}{n} \, n (n - 1) \E \kappa_{j - 1}(G_{m - 1}),
\end{split}
\]
where in the last step we used the definition \eqref{eq:high-kappaj}. Inserting this bound into \eqref{eq:high-deletion-1} gives \eqref{eq:high-deletion-np}.

We now prove (b). Let $(a, v)$ be a pair counted in $n(n - 1) \kappa_j^{\rm piv}(G_m)$. At most $K - 1$ of the $j$ occurrences at $a$ are non-pivotal for $(a, v)$, so there exist at least $j - K + 1$ indices $i \le m$ such that $a \in e_i$ and $a \not \sim_{G_m - e_i} v$. Further, $\deg_{G_m - e_i}(a) = j - 1$ for these indices.

We set
\[
h_2(G', e) = \sum_{a \in e} \sum_{v \in [n] \setminus \{a\}} \ind \bigl\{ a \sim_{G' + e} v, \ a \not \sim_{G'} v, \ \deg_{G'}(a) = j - 1 \bigr\}.
\]
Each index $i$ as above contributes $1$ to $h_2(G_m - e_i, e_i)$, and as in the proof of (a), we obtain
\begin{equation}\label{eq:high-deletion-2}
(j - K + 1) \, n (n - 1) \E \kappa_j^{\rm piv}(G_m) \le m \E h_2(G_{m - 1}, e).
\end{equation}

To compute the right side of \eqref{eq:high-deletion-2}, we condition on $G' = G_{m - 1}$. We may assume that $G'$ is colorable, since otherwise $h_2(G', e) = 0$. Let $a$ and $y$ be distinct vertices. We claim that
\begin{equation}\label{eq:high-deletion-3}
\sum_{v \in [n] \setminus \{a\}} \ind \bigl\{ a \sim_{G' + ay} v, \ a \not \sim_{G'} v \bigr\} = \abs{R_{G'}(a; y)}.
\end{equation}
Indeed, suppose first that $y \in C_{G'}(a)$. Then $G' + ay$ is not colorable by \cref{lem:pre-frozen}(a), the left side vanishes, and $R_{G'}(a; y) = \emptyset$ by \cref{def:high-relation}. Otherwise, \cref{lem:high-relation}(b) shows that the vertices $v$ counted on the left are exactly the elements of $R_{G'}(a; y)$.

Every pair $ay$ has probability $\binom{n}{2}^{-1}$, and summing over the two endpoints of $e$ amounts to summing over ordered pairs $(a, y)$. Using \eqref{eq:high-deletion-3} and $R_{G'}(a; a) = \emptyset$, we obtain
\[
\begin{split}
\E \bigl[ h_2(G', e) \bigm| G' \bigr] &= \binom{n}{2}^{-1} \sum_{a \in [n]} \ind \{ \deg_{G'}(a) = j - 1 \} \sum_{y \in [n] \setminus \{a\}} \abs{R_{G'}(a; y)}\\
&= \binom{n}{2}^{-1} \sum_{a \in [n]} \ind \{ \deg_{G'}(a) = j - 1 \} \, \abs{R_{G'}(a)}\\
&= \frac{2n}{n - 1} \sum_{a \in [n]} \ind \{ \deg_{G'}(a) = j - 1 \} \, \beta_{G'}(a),
\end{split}
\]
where in the last step we used \eqref{eq:high-beta}. Recall that the law of $G_{m - 1}$ is invariant under relabeling. Relabeling $G$ and $a$ by the same permutation does not change $\ind \{ \deg_G(a) = j - 1 \} \beta_G(a)$. Every term of the last sum has the same expectation. For all $u \in [n]$, we have
\[
\E h_2(G_{m - 1}, e) = \frac{2n^2}{n - 1} \, \E \bigl[ \ind \{ \deg_{G_{m - 1}}(u) = j - 1 \} \, \beta_{G_{m - 1}}(u) \bigr].
\]
Inserting this identity into \eqref{eq:high-deletion-2} and dividing by $(j - K + 1) n (n - 1)$, we arrive at \eqref{eq:high-deletion-piv}.
\end{proof}

We now iterate \cref{lem:high-deletion}(a). Recall that $K = \lceil 8 \Lambda \rceil$, so that the coefficient in \eqref{eq:high-deletion-np} satisfies $2m/(nK) \le 1/4$ for $m \le M$. Each application lowers the degree and the time by one. After $j - w_0$ applications, the term $\E \kappa_j(G_m)$ with $j > D_n$ is bounded by $4^{-(j - w_0)} \E \kappa_{w_0}(G_{m - j + w_0})$ plus a combination of the terms $\E \kappa_{j'}^{\rm piv}$ with $j' > w_0$; see \eqref{eq:high-iterate} below. The terms of the first kind are controlled by \cref{lem:pre-hazard}, which uses that the process becomes non-colorable at most once.

\begin{proposition}\label{prop:high-descent}
Let $n \ge 3$. Then
\begin{equation}\label{eq:high-descent}
\begin{split}
&\Pp \Bigl( H \le M, \ \max_{v \in e_H} \deg_\Gamma(v) > D_n \Bigr)\\
&\qquad \le \frac{8}{3} \, 4^{-(D_n + 1 - w_0)} + \frac{8}{3} \sum_{j = w_0 + 1}^{M - 1} \sum_{m = j}^{M - 1} \E \kappa_j^{\rm piv}(G_m),
\end{split}
\end{equation}
where an empty sum is zero.
\end{proposition}

\begin{proof}
We may assume that $M \ge 1$, since $H \ge 1$ and the left side of \eqref{eq:high-descent} vanishes if $M = 0$. We define $\cE$ to be the set of pairs $(G, ab)$ such that $G$ is a colorable multigraph on $[n]$, $ab$ is a frozen pair of $G$, and $\max(\deg_G(a), \deg_G(b)) > D_n$. We recall the density $\kappa_\cE$ of \eqref{eq:pre-kappaE}. We claim that
\begin{equation}\label{eq:high-descent-1}
\Pp \Bigl( H \le M, \ \max_{v \in e_H} \deg_\Gamma(v) > D_n \Bigr) = \sum_{m = 0}^{M - 1} \E \kappa_\cE(G_m) \le 2 \sum_{m = 0}^{M - 1} \sum_{j = D_n + 1}^{m} \E \kappa_j(G_m).
\end{equation}
Since $H \ge 1$, the equality follows from \cref{lem:pre-terminal}(a) with $h = 1$, $L' = 0$, and $U' = M$. For the inequality, note that each frozen pair $ab$ counted in $\kappa_\cE(G_m)$ yields at least one ordered pair, $(a, b)$ or $(b, a)$, whose first vertex has degree $j > D_n$, that the degrees in $G_m$ do not exceed $m$, and that $\binom{n}{2}^{-1} = 2/(n(n - 1))$.

Next, let $1 \le j \le m \le M$. Since $\kappa_j = \kappa_j^{\rm piv} + \kappa_j^{\rm np}$, the bound \eqref{eq:high-deletion-np} gives
\begin{equation}\label{eq:high-recursion}
\E \kappa_j(G_m) \le \E \kappa_j^{\rm piv}(G_m) + \frac{2m}{nK} \E \kappa_{j - 1}(G_{m - 1}) \le \E \kappa_j^{\rm piv}(G_m) + \frac{1}{4} \E \kappa_{j - 1}(G_{m - 1}),
\end{equation}
where in the last step we used $m \le M \le \Lambda n$ and $K \ge 8 \Lambda$.

Now let $w_0 < j \le m \le M - 1$. We apply \eqref{eq:high-recursion} successively at the degrees $j, j - 1, \dots, w_0 + 1$ and the corresponding times $m, m - 1, \dots, m - j + w_0 + 1$. At each step, the degree $j'$ and the time $m - j + j'$ satisfy $1 \le j' \le m - j + j' \le M$. By induction on $j - w_0$, we obtain
\begin{equation}\label{eq:high-iterate}
\E \kappa_j(G_m) \le \sum_{j' = w_0 + 1}^{j} 4^{-(j - j')} \, \E \kappa_{j'}^{\rm piv}(G_{m - j + j'}) + 4^{-(j - w_0)} \, \E \kappa_{w_0}(G_{m - j + w_0}).
\end{equation}

We insert \eqref{eq:high-iterate} into \eqref{eq:high-descent-1}, which is possible since $D_n > w_0$, and we consider the two kinds of terms separately. For the last term, fix $j \ge D_n + 1$. As $m$ ranges over $\{j, \dots, M - 1\}$, the times $m - j + w_0$ are distinct. Using $\kappa_{w_0} \le \kappa$ and \eqref{eq:pre-hazard-sum}, we obtain
\[
\sum_{m = j}^{M - 1} \E \kappa_{w_0}(G_{m - j + w_0}) \le 1.
\]
Summing over $j$, we find that these terms contribute at most
\[
2 \sum_{j = D_n + 1}^{\infty} 4^{-(j - w_0)} = \frac{8}{3} \, 4^{-(D_n + 1 - w_0)}
\]
to the right side of \eqref{eq:high-descent-1}.

For the first sum in \eqref{eq:high-iterate}, fix $j' \ge w_0 + 1$ and $j \ge \max(j', D_n + 1)$. As $m$ ranges over $\{j, \dots, M - 1\}$, the times $m - j + j'$ are distinct and lie in $\{j', \dots, M - 1\}$. These times need not be distinct when $j$ varies, but we account for the repetitions by the geometric factor. Since all terms are nonnegative, the first sum in \eqref{eq:high-iterate} contributes at most
\[
2 \sum_{j' = w_0 + 1}^{M - 1} \biggl( \sum_{j \ge j'} 4^{-(j - j')} \biggr) \sum_{m = j'}^{M - 1} \E \kappa_{j'}^{\rm piv}(G_m) = \frac{8}{3} \sum_{j' = w_0 + 1}^{M - 1} \sum_{m = j'}^{M - 1} \E \kappa_{j'}^{\rm piv}(G_m)
\]
to the right side of \eqref{eq:high-descent-1}. Adding the two contributions and renaming $j'$ as $j$ gives \eqref{eq:high-descent}.
\end{proof}

The descent in \cref{prop:high-descent} runs over all degrees, since the degrees in $G_m$ are at most $m$ and the time shifts are paid for by the factor $4^{-(j - j')}$.

\section{Credited sets}\label{sec:high-relation}

In this section and the next, we fix a vertex $u \in [n]$, which we call the \emph{root}. For a colorable multigraph $G$, we call $\log \abs{C_G(u)}$ its \emph{potential} at $u$. By \cref{lem:high-relation}(b), the relation $R_G(u)$ describes the growth of $C_G(u)$ when an edge at the root is added. The next definition assigns a set to every pair. When adding the pair preserves colorability, this set gives a lower bound for the growth of the root's frozen class.

\begin{definition}\label{def:high-credited}
Let $G$ be a colorable multigraph on $[n]$, and recall that $W = [n] \setminus C_G(u)$ and that $R(y)$ denotes the section $R_G(u; y)$. For a pair $f = ab$ of $[n]$, we define the \emph{credited set} of $f$ by
\begin{equation}\label{eq:high-credited}
\mathscr{A}_G(f) =
\begin{cases}
R(a) \cap R(b) & \text{if } a, b \in W, \\
R(b) & \text{if } a \in C_G(u) \text{ and } b \in W, \\
R(a) & \text{if } a \in W \text{ and } b \in C_G(u), \\
\emptyset & \text{if } a, b \in C_G(u).
\end{cases}
\end{equation}
If $G$ is not colorable, we set $\mathscr{A}_G(f) = \emptyset$.
\end{definition}

For a colorable multigraph $G$ and a pair $f$, we call
\[
\log \biggl( 1 + \frac{\abs{\mathscr{A}_G(f)}}{\abs{C_G(u)}} \biggr)
\]
the \emph{credited logarithmic increment} of $f$ at $u$.

We call the first case of \eqref{eq:high-credited} the \emph{direct case} and the second and third cases the \emph{cross case}. Note that $\mathscr{A}_G(f) \subset W$. The cross case contains the pairs $f = uy$ with $y \in W$, for which $\mathscr{A}_G(uy) = R(y)$; by \cref{lem:high-relation}(b), this is exactly the set of vertices that join the class of the root when $uy$ is added. The credited set is defined whether or not $G + f$ is colorable, and this will be important in the proof of \cref{thm:high-potential} (see \cref{rem:high-killing}).

\begin{lemma}\label{lem:high-absorb}
Let $G$ be a colorable multigraph on $[n]$, and let $f$ be a pair of $[n]$ such that $G + f$ is colorable. Then $C_G(u)$ and $\mathscr{A}_G(f)$ are disjoint subsets of $C_{G + f}(u)$. In particular, $\abs{C_{G + f}(u)} \ge \abs{C_G(u)} + \abs{\mathscr{A}_G(f)}$.
\end{lemma}

\begin{proof}
By \cref{lem:pre-frozen}(b), we have $C_G(u) \subset C_{G + f}(u)$, and $\mathscr{A}_G(f) \subset W$ is disjoint from $C_G(u)$. Let $f = ab$ and $y \in \mathscr{A}_G(f)$. Then $y \in W$, and \cref{lem:high-relation}(b) shows that $C_{G + uy}(u) = C_G(u) \cup R(y)$. We claim that $a, b \in C_{G + uy}(u)$.

In the direct case, we have $y \in R(a) \cap R(b)$, and the symmetry of $R_G(u)$ in \cref{lem:high-relation}(a) gives $a, b \in R(y)$. In the cross case, we may assume that $a \in C_G(u)$ and $y \in R(b)$, and then $b \in R(y)$ by symmetry. This proves the claim.

The claim implies that $a \sim_{G + uy} b$. Since $a \ne b$, \cref{lem:pre-frozen}(a) shows that $G + uy + f$ is not colorable. Every coloring $\sigma$ of $G + f$ therefore satisfies $\sigma(y) = \sigma(u)$, since otherwise $\sigma$ would be a coloring of $G + f + uy$. As $G + f$ is colorable, this means that $y \in C_{G + f}(u)$.
\end{proof}

\Cref{fig:high-absorb} illustrates the two cases of the proof.

\begin{figure}[ht]
\centering
\begin{tikzpicture}[x = 0.95cm, y = 0.95cm,
  vtx/.style = {circle, fill = black, inner sep = 1.3pt},
  rel/.style = {densely dotted, very thick},
  virt/.style = {dashed, thick},
  act/.style = {line width = 1.4pt},
  hull/.style = {gray, dash pattern = on 2pt off 2pt, rounded corners = 9pt},
  lab/.style = {font = \small}]
\begin{scope}
  \draw[hull] (-1.25, 1.75) -- (0.95, 1.75) -- (4.25, -0.05) -- (4.25, -2.0) -- (-1.25, -2.0) -- cycle;
  \draw[thick] (0, 0) ellipse (0.85 and 1.3);
  \node[vtx, label = {[lab]left:$u$}] (u) at (0, 0.55) {};
  \node[vtx] at (-0.3, -0.35) {};
  \node[vtx] at (0.3, -0.75) {};
  \node[lab] at (0, -1.6) {$C_G(u)$};
  \node[vtx, label = {[lab]above:$y$}] (y) at (2.75, 1.45) {};
  \node[vtx, label = {[lab]below:$a$}] (a) at (2.05, -0.65) {};
  \node[vtx, label = {[lab]below:$b$}] (b) at (3.55, -0.65) {};
  \draw[virt] (u) -- (y);
  \draw[rel] (y) -- (a);
  \draw[rel] (y) -- (b);
  \draw[act] (a) -- (b);
  \node[lab, gray] at (2.85, -1.6) {$C_{G + uy}(u)$};
  \node[lab] at (1.5, -2.45) {direct case};
\end{scope}
\begin{scope}[xshift = 6.6cm]
  \draw[hull] (-1.25, 1.75) -- (0.95, 1.75) -- (3.45, 0.15) -- (3.45, -2.0) -- (-1.25, -2.0) -- cycle;
  \draw[thick] (0, 0) ellipse (0.85 and 1.3);
  \node[vtx, label = {[lab]left:$u$}] (u) at (0, 0.55) {};
  \node[vtx, label = {[lab]left:$a$}] (a) at (0.2, -0.65) {};
  \node[vtx] at (-0.4, -0.15) {};
  \node[lab] at (0, -1.6) {$C_G(u)$};
  \node[vtx, label = {[lab]above:$y$}] (y) at (2.6, 1.45) {};
  \node[vtx, label = {[lab]below:$b$}] (b) at (2.55, -0.65) {};
  \draw[virt] (u) -- (y);
  \draw[rel] (y) -- (b);
  \draw[act] (a) -- (b);
  \node[lab, gray] at (2.2, -1.6) {$C_{G + uy}(u)$};
  \node[lab] at (1.1, -2.45) {cross case};
\end{scope}
\end{tikzpicture}

\caption{The credited set of a pair $f = ab$, with the root $u$ and its class $C_G(u)$ (black outline). Dotted lines join pairs of $R_G(u)$. The dashed line is the pair $uy$, which is not added to $G$, and the thick solid line is the added pair $ab$. Left, the direct case: $a, b \in W$ and $y \in R(a) \cap R(b)$. Right, the cross case: $a \in C_G(u)$, $b \in W$, and $y \in R(b)$. In both cases, adding $uy$ would put $a$ and $b$ into the class of $u$ (gray outline), so $G + uy + ab$ is not colorable. Hence every coloring of $G + ab$ gives $y$ the color of $u$.}
\label{fig:high-absorb}
\end{figure}

We next compare $\abs{R_G(u)}$ with two sums, which will appear as the contributions from the direct and cross cases of \eqref{eq:high-credited} to the expected credited logarithmic increment. The comparison is purely combinatorial; it uses the symmetry of $R_G(u)$ and nothing else about colorings.

For a colorable multigraph $G$ on $[n]$ and the fixed root $u$, we define
\begin{equation}\label{eq:high-ain-acr}
\begin{split}
a_1 &= \sum_{(a, b) \in W \times W} \min \biggl( \frac{\abs{R(a) \cap R(b)}}{\abs{C_G(u)}}, 1 \biggr), \\
a_2 &= 2 \bigl( \abs{C_G(u)} - 1 \bigr) \sum_{a \in W} \min \biggl( \frac{\abs{R(a)}}{\abs{C_G(u)}}, 1 \biggr).
\end{split}
\end{equation}
\begin{lemma}\label{lem:high-comb}
For every colorable multigraph $G$ on $[n]$, we have $\abs{R_G(u)} \le 2 a_1 + a_2$.
\end{lemma}

\begin{proof}
We first bound $a_1$ from below. Let $a, b \in W$. Writing the minimum as a sum over the intersection, we have
\[
\min \biggl( \frac{\abs{R(a) \cap R(b)}}{\abs{C_G(u)}}, 1 \biggr) = \sum_{y \in R(a) \cap R(b)} \frac{1}{\max \bigl( \abs{C_G(u)}, \abs{R(a) \cap R(b)} \bigr)},
\]
where both sides vanish if $R(a) \cap R(b) = \emptyset$.

By the symmetry of $R_G(u)$, we have $y \in R(a) \cap R(b)$ if and only if $a, b \in R(y)$. Exchanging the order of summation, we obtain
\[
a_1 = \sum_{y \in W} \sum_{a \in R(y)} \sum_{b \in R(y)} \frac{1}{\max \bigl( \abs{C_G(u)}, \abs{R(a) \cap R(b)} \bigr)} \ge \sum_{y \in W} \sum_{a \in R(y)} \frac{\abs{R(y)}}{\max \bigl( \abs{C_G(u)}, \abs{R(a)} \bigr)},
\]
where in the last step we used $\abs{R(a) \cap R(b)} \le \abs{R(a)}$ and summed over $b \in R(y)$.

The right side is a sum over the pairs $(y, a) \in R_G(u)$. Since $(y, a) \mapsto (a, y)$ is a bijection of $R_G(u)$, the same lower bound holds with the roles of $y$ and $a$ exchanged. Adding the two bounds, we find
\begin{equation}\label{eq:high-comb-1}
2 a_1 \ge \sum_{(y, a) \in R_G(u)} \biggl( \frac{\abs{R(y)}}{\max \bigl( \abs{C_G(u)}, \abs{R(a)} \bigr)} + \frac{\abs{R(a)}}{\max \bigl( \abs{C_G(u)}, \abs{R(y)} \bigr)} \biggr).
\end{equation}

We now split $R_G(u)$ into the pairs $(y, a)$ with $\max(\abs{R(y)}, \abs{R(a)}) \ge \abs{C_G(u)}$ and the remaining pairs. For a pair of the first kind, the summand in \eqref{eq:high-comb-1} is at least $1$. Indeed, by the symmetry of the summand, we may assume that $\abs{R(y)} \ge \abs{R(a)}$. Then $\abs{R(y)} \ge \max(\abs{C_G(u)}, \abs{R(a)})$, and the first fraction in the summand is at least $1$. The bound \eqref{eq:high-comb-1} therefore implies that there are at most $2 a_1$ pairs of the first kind.

It remains to count the pairs $(y, a) \in R_G(u)$ with $\abs{R(y)} < \abs{C_G(u)}$ and $\abs{R(a)} < \abs{C_G(u)}$. When $\abs{C_G(u)} = 1$, there are none, since every pair $(y, a) \in R_G(u)$ satisfies $a \in R(y)$ and hence $\abs{R(y)} \ge 1$. Suppose that $\abs{C_G(u)} \ge 2$.

By symmetry, the vertices $y$ with $(y, a) \in R_G(u)$ form the set $R(a)$, so the number of pairs $(y, a) \in R_G(u)$ with $\abs{R(a)} < \abs{C_G(u)}$ satisfies
\[
\sum_{\substack{a \in W \\ \abs{R(a)} < \abs{C_G(u)}}} \abs{R(a)} \le \abs{C_G(u)} \sum_{a \in W} \min \biggl( \frac{\abs{R(a)}}{\abs{C_G(u)}}, 1 \biggr) \le a_2,
\]
where we used $\abs{C_G(u)} \le 2 (\abs{C_G(u)} - 1)$ in the last step. Adding the two counts gives $\abs{R_G(u)} \le 2 a_1 + a_2$.
\end{proof}

The following lemma is the main step of this section. It bounds $\beta_G(u)$ by the expected credited logarithmic increment when one uniform edge is added to $G$.

\begin{lemma}\label{lem:high-onestep}
Let $G$ be a colorable multigraph on $[n]$, and let $f$ be a uniform pair of $[n]$. Then
\begin{equation}\label{eq:high-onestep}
\beta_G(u) \le \frac{2}{\log 2} \, \E_f \log \biggl( 1 + \frac{\abs{\mathscr{A}_G(f)}}{\abs{C_G(u)}} \biggr),
\end{equation}
where $\E_f$ denotes the expectation with respect to $f$.
\end{lemma}

\begin{proof}
For $z \ge 0$, we have $\log(1 + z) \ge (\log 2) \min(z, 1)$. For $z \in [0, 1]$, this follows from the concavity of $z \mapsto \log(1 + z)$, and for $z > 1$ from its monotonicity. It suffices to show that
\begin{equation}\label{eq:high-onestep-1}
\E_f \min \biggl( \frac{\abs{\mathscr{A}_G(f)}}{\abs{C_G(u)}}, 1 \biggr) \ge \frac{\abs{R_G(u)}}{2 n (n - 1)},
\end{equation}
since the right side of \eqref{eq:high-onestep-1} is bounded below by $\abs{R_G(u)} / (2n^2) = \beta_G(u)/2$.

Since each pair has probability $\binom{n}{2}^{-1} = 2/(n(n - 1))$, we have
\begin{equation}\label{eq:high-onestep-2}
\E_f \min \biggl( \frac{\abs{\mathscr{A}_G(f)}}{\abs{C_G(u)}}, 1 \biggr) = \frac{1}{n(n - 1)} \sum_{\substack{(a, b) \in [n]^2 \\ a \ne b}} \min \biggl( \frac{\abs{\mathscr{A}_G(ab)}}{\abs{C_G(u)}}, 1 \biggr).
\end{equation}
We split the sum in \eqref{eq:high-onestep-2} according to the cases of \eqref{eq:high-credited}. The ordered pairs of distinct vertices of $W$ contribute
\[
a_1 - \sum_{a \in W} \min \biggl( \frac{\abs{R(a)}}{\abs{C_G(u)}}, 1 \biggr),
\]
since the terms of $a_1$ with $a = b$ are $\min(\abs{R(a)} / \abs{C_G(u)}, 1)$.

The ordered pairs with one vertex in $C_G(u)$ and one in $W$ contribute
\[
2 \abs{C_G(u)} \sum_{a \in W} \min \biggl( \frac{\abs{R(a)}}{\abs{C_G(u)}}, 1 \biggr),
\]
since there are $\abs{C_G(u)}$ choices for the vertex in $C_G(u)$ and two orders. The pairs inside $C_G(u)$ contribute zero.

Adding these contributions, we find that the sum in \eqref{eq:high-onestep-2} satisfies
\[
\begin{split}
\sum_{\substack{(a, b) \in [n]^2 \\ a \ne b}} \min \biggl( \frac{\abs{\mathscr{A}_G(ab)}}{\abs{C_G(u)}}, 1 \biggr) &= a_1 + \bigl( 2 \abs{C_G(u)} - 1 \bigr) \sum_{a \in W} \min \biggl( \frac{\abs{R(a)}}{\abs{C_G(u)}}, 1 \biggr)\\
&\ge a_1 + a_2 \ge \frac{1}{2} \abs{R_G(u)},
\end{split}
\]
where we used $2 \abs{C_G(u)} - 1 \ge 2 (\abs{C_G(u)} - 1)$ in the first inequality, and $a_2 \ge 0$ and \cref{lem:high-comb} in the second. Dividing by $n(n - 1)$ gives \eqref{eq:high-onestep-1}.
\end{proof}

The pairs $f$ containing the root compensate for the diagonal terms $a = b$ of $a_1$, which do not correspond to pairs of distinct vertices. Without them, \eqref{eq:high-onestep} would carry an additive error of order $1/n$, and the sum of these errors over the $M$ steps of the process would be of order $\Lambda$.

\section{The potential}\label{sec:high-potential}

We now sum the one-step bound of \cref{lem:high-onestep} over time. The potential $\log \abs{C_{G_m}(u)}$ takes values in $[0, \log n]$. By \cref{lem:high-absorb}, it increases at each step at least by the increment appearing in \eqref{eq:high-onestep}, as long as the graph remains colorable. However, \cref{lem:high-absorb} does not apply at the step at which the graph becomes non-colorable. The increment is still defined at that step, and we bound it using that $C_G(u)$ and $\mathscr{A}_G(f)$ are disjoint subsets of $[n]$.

For each integer $m \ge 0$ with $m < H$, we define
\[
\Phi_m = \log \abs{C_{G_m}(u)}, \qquad \Psi_m = \log \biggl( 1 + \frac{\abs{\mathscr{A}_{G_m}(e_{m + 1})}}{\abs{C_{G_m}(u)}} \biggr).
\]
For $m \ge H$, we set $\Phi_m = \Psi_m = 0$. For $m < H$, the variable $\Psi_m$ is the credited logarithmic increment of $e_{m + 1}$ at $u$.

\begin{theorem}\label{thm:high-potential}
Let $u \in [n]$, let $\tau = \min(H, M)$, and let $\tau_0$ be a stopping time of the filtration $(\cF_m)_{m \ge 0}$ with values in $\{0, 1, 2, \dots\} \cup \{\infty\}$. Then
\begin{equation}\label{eq:high-potential}
\E \sum_{m : \, \tau_0 \le m < \tau} \beta_{G_m}(u) \le \frac{2}{\log 2} \, (\log n) \, \Pp(\tau_0 < \tau).
\end{equation}
\end{theorem}

\begin{proof}
For $m < H$, the multigraph $G_m$ is colorable, so $0 \le \Phi_m \le \log n$ and $\Psi_m \ge 0$. The variable $\Psi_m$ is determined by $G_m$ and $e_{m + 1}$, whether or not $G_{m + 1}$ is colorable.

We first prove a pathwise bound. Let $m + 1 < H$. Then $G_{m + 1} = G_m + e_{m + 1}$ is colorable, and \cref{lem:high-absorb} gives
\begin{equation}\label{eq:high-step}
\abs{C_{G_{m + 1}}(u)} \ge \abs{C_{G_m}(u)} + \abs{\mathscr{A}_{G_m}(e_{m + 1})}, \qquad \text{and hence} \qquad \Psi_m \le \Phi_{m + 1} - \Phi_m.
\end{equation}

Suppose now that $\tau_0 < \tau$. Then $\tau - 1 < H$. The sets $C_{G_{\tau - 1}}(u)$ and $\mathscr{A}_{G_{\tau - 1}}(e_\tau)$ are disjoint subsets of $[n]$, so we have
\[
\Psi_{\tau - 1} = \log \bigl( \abs{C_{G_{\tau - 1}}(u)} + \abs{\mathscr{A}_{G_{\tau - 1}}(e_\tau)} \bigr) - \Phi_{\tau - 1} \le \log n - \Phi_{\tau - 1}.
\]
For $\tau_0 \le m \le \tau - 2$, we have $m + 1 < \tau \le H$, so \eqref{eq:high-step} applies. Telescoping, we obtain
\[
\sum_{m = \tau_0}^{\tau - 1} \Psi_m \le \sum_{m = \tau_0}^{\tau - 2} (\Phi_{m + 1} - \Phi_m) + \log n - \Phi_{\tau - 1} = \log n - \Phi_{\tau_0} \le \log n.
\]

The sum is empty if $\tau_0 \ge \tau$. Together, these cases give
\begin{equation}\label{eq:high-telescope}
\sum_{m : \, \tau_0 \le m < \tau} \Psi_m \le (\log n) \ind \{ \tau_0 < \tau \}.
\end{equation}

Next, we compare the expectations. For $0 \le m < M$, we define $I_m$ by
\[
I_m = \ind \{ \tau_0 \le m < \tau \}.
\]
Since $\tau_0$ is a stopping time and, for $m < M$, the event $\{ m < \tau \}$ is the event that $G_m$ is colorable, the variable $I_m$ is $\cF_m$-measurable.

The pair $e_{m + 1}$ is uniform and independent of $\cF_m$. On the event $\{ m < H \}$, the conditional expectation of $\Psi_m$ given $\cF_m$ is the expectation in \eqref{eq:high-onestep} with $G = G_m$. Using \cref{lem:high-onestep} and $I_m \le \ind \{ m < H \}$, we obtain
\[
\E \bigl[ I_m \beta_{G_m}(u) \bigr] \le \frac{2}{\log 2} \, \E \bigl[ I_m \, \E [ \Psi_m \mid \cF_m ] \bigr] = \frac{2}{\log 2} \, \E \bigl[ I_m \Psi_m \bigr].
\]
Since $\tau \le M$, summing over $0 \le m < M$ and using \eqref{eq:high-telescope}, we conclude that
\[
\begin{split}
\E \sum_{m : \, \tau_0 \le m < \tau} \beta_{G_m}(u) &= \sum_{m = 0}^{M - 1} \E \bigl[ I_m \beta_{G_m}(u) \bigr] \le \frac{2}{\log 2} \, \E \sum_{m : \, \tau_0 \le m < \tau} \Psi_m\\
&\le \frac{2}{\log 2} \, (\log n) \, \Pp(\tau_0 < \tau).
\end{split}
\]
This completes the proof.
\end{proof}

The passage from the one-step bound to \eqref{eq:high-potential} is a standard compensator argument, in the spirit of the drift bounds for hitting and occupation times in \cite[Sections~2--3]{hajek1982hitting}. A related use of a bounded monotone potential to control occupation, retaining terminal increments, appears in \cite[Lemmas~2.2 and~4.1]{openai2026linear}. Here the potential is the logarithm of a frozen-class size, and the input specific to colorings is \cref{lem:high-onestep}.

\begin{remark}\label{rem:high-killing}
The proof of \cref{thm:high-potential} credits the edge $e_\tau$ with the increment $\Psi_{\tau - 1}$, which is defined from $G_{\tau - 1}$ and $e_\tau$ even when $G_\tau$ is not colorable. If $\tau = H$, this edge is the final edge; if $\tau = M < H$, it is the edge at the cap. Correspondingly, the average in \cref{lem:high-onestep} runs over all pairs $f$, including those for which $G + f$ is not colorable. This cannot be avoided, as the following example shows. Let $q = 3$, let $n \ge 7$ with $n - 1$ divisible by $3$, and let the root be $u = n$. Let $G$ consist of the triangle on the vertices $1, 2, 3$, of $(n - 4)/3$ vertices joined to $2$ and $3$, of $(n - 4)/3$ vertices joined to $1$ and $3$, of $(n - 4)/3$ vertices joined to $1$ and $2$, and of the edge joining $u$ to $3$. In every coloring, the triangle receives three distinct colors, and every other vertex except $u$ receives the color of the vertex of the triangle to which it is not joined. Thus $G$ has three frozen classes besides $\{u\}$, each of size $(n - 1)/3$. For $i \in \{1, 2, 3\}$, we let $C_i$ denote the frozen class that contains the vertex $i$. The root can take the color of $C_1$ or that of $C_2$, so $C_G(u) = \{u\}$, and \eqref{eq:high-R} shows that $R_G(u) = (C_1 \times C_2) \cup (C_2 \times C_1)$. Therefore
\[
\beta_G(u) = \frac{2 (n - 1)^2}{9 n^2},
\]
which tends to $2/9$ as $n \to \infty$.

On the other hand, a pair $f$ with $\mathscr{A}_G(f) \ne \emptyset$ either lies inside $C_1$ or inside $C_2$, in which case $G + f$ is not colorable, or joins $u$ to a vertex of $C_1 \cup C_2$. There are $2(n - 1)/3$ pairs of the last kind, each with a credited set of size $(n - 1)/3$, so
\[
\E_f \biggl[ \log \biggl( 1 + \frac{\abs{\mathscr{A}_G(f)}}{\abs{C_G(u)}} \biggr) \ind \{ G + f \text{ is colorable} \} \biggr] = \binom{n}{2}^{-1} \frac{2(n - 1)}{3} \log \frac{n + 2}{3} \le \frac{4 \log n}{3n}.
\]
Let $C, C' \ge 0$ be constants independent of $n$. Then a bound of the form
\[
\beta_G(u) \le C \, \E_f \biggl[ \log \biggl( 1 + \frac{\abs{\mathscr{A}_G(f)}}{\abs{C_G(u)}} \biggr) \ind \{ G + f \text{ is colorable} \} \biggr] + C' \frac{\log n}{n},
\]
fails for large $n$. However, this does not contradict \cref{thm:high-potential}. About a third of all pairs lie inside one of the classes $C_1$, $C_2$, $C_3$, and $\kappa(G)$ is close to $1/3$. By \cref{lem:pre-lifetime}(c) and (d), on the event $\{G_m = G\}$, the conditional expectation of $H - m$ given $\cF_m$ is $r(G) \le 1/\kappa(G)$, which is bounded.
\end{remark}

We apply \cref{thm:high-potential} at the time at which the root reaches a given degree.

\begin{corollary}\label{cor:high-degree}
For all $u \in [n]$ and all integers $j \ge 1$,
\begin{equation}\label{eq:high-degree}
\sum_{m = 0}^{M - 1} \E \bigl[ \ind \{ \deg_{G_m}(u) \ge j \} \, \beta_{G_m}(u) \bigr] \le \frac{2}{\log 2} \, (\log n) \biggl( \frac{2 e \Lambda}{j} \biggr)^j.
\end{equation}
\end{corollary}

\begin{proof}
We may assume that $M \ge 1$, since otherwise the left side vanishes. We set $\tau = \min(H, M)$ and let $\tau_0$ denote the first time $m \ge 0$ such that $\deg_{G_m}(u) \ge j$, with $\tau_0 = \infty$ if there is no such time. Then $\tau_0$ is a stopping time.

Since degrees do not decrease in $m$, we have $\deg_{G_m}(u) \ge j$ if and only if $\tau_0 \le m$. Further, $\beta_{G_m}(u) = 0$ for $m \ge H$. The left side of \eqref{eq:high-degree} is the left side of \eqref{eq:high-potential}. By \cref{thm:high-potential}, it is bounded by $(2 / \log 2) (\log n) \Pp(\tau_0 < \tau)$.

On the event $\{ \tau_0 < \tau \}$, we have $\tau_0 \le M - 1$ and hence $\deg_{G_{M - 1}}(u) \ge j \ge 1$; in particular, this event is empty if $M = 1$. Suppose that $M \ge 2$. Recall that $\deg_{G_{M - 1}}(u)$ has the law $\Bin(M - 1, 2/n)$. Using \cref{lem:pre-tails}(a) and $M - 1 \le \Lambda n$, we obtain
\[
\Pp(\tau_0 < \tau) \le \Pp \bigl( \Bin(M - 1, 2/n) \ge j \bigr) \le \biggl( \frac{2 e (M - 1)}{n j} \biggr)^j \le \biggl( \frac{2 e \Lambda}{j} \biggr)^j,
\]
which completes the proof.
\end{proof}

\section{Combining the estimates}\label{sec:high-main}

We now combine \cref{prop:high-descent}, \cref{lem:high-deletion}(b), and \cref{cor:high-degree}.

\begin{theorem}\label{thm:high-main}
Let $n \ge 3$ be such that $w_0 \ge K$, where $K$ and $w_0$ are given by \eqref{eq:high-params}. Then
\begin{equation}\label{eq:high-main}
\begin{split}
&\Pp \Bigl( H \le M, \ \max_{v \in e_H} \deg_\Gamma(v) > D_n \Bigr)\\
&\qquad \le \frac{8}{3} \, 4^{-(D_n + 1 - w_0)} + \frac{32 \Lambda}{3 \log 2} \biggl( \frac{n}{n - 1} \biggr)^2 \frac{\log n}{w_0 - K + 2} \biggl( \frac{2 e \Lambda}{w_0} \biggr)^{w_0}.
\end{split}
\end{equation}
Further, there exists $n_0$, depending only on $\Lambda$, such that for all $n \ge n_0$ the right side of \eqref{eq:high-main} is at most $(\log n)^{-3/2}$. In particular, for all $n \ge n_0$ and all integers $0 \le L' \le U' \le M$,
\begin{equation}\label{eq:high-main-window}
\Pp \Bigl( L' < H \le U', \ \min_{v \in e_H} \deg_\Gamma(v) > D_n \Bigr) \le (\log n)^{-3/2}.
\end{equation}
\end{theorem}

\begin{proof}
We bound the pivotal sum in \eqref{eq:high-descent}. Fix $u \in [n]$ and let $w_0 + 1 \le j \le m \le M - 1$. Then $j \ge K + 1$, so that \eqref{eq:high-deletion-piv} applies, and $j - K + 1 \ge w_0 - K + 2$. Using \eqref{eq:high-deletion-piv} and $m \le M \le \Lambda n$, we find
\[
\E \kappa_j^{\rm piv}(G_m) \le \frac{2 \Lambda}{w_0 - K + 2} \biggl( \frac{n}{n - 1} \biggr)^2 \E \bigl[ \ind \{ \deg_{G_{m - 1}}(u) = j - 1 \} \, \beta_{G_{m - 1}}(u) \bigr].
\]

For fixed $m$, the events $\{ \deg_{G_{m - 1}}(u) = j - 1 \}$ with $w_0 + 1 \le j \le m$ are disjoint and contained in $\{ \deg_{G_{m - 1}}(u) \ge w_0 \}$. Exchanging the order of summation, we obtain
\[
\begin{split}
\sum_{j = w_0 + 1}^{M - 1} \sum_{m = j}^{M - 1} \E \bigl[ \ind \{ \deg_{G_{m - 1}}(u) = j - 1 \} \, \beta_{G_{m - 1}}(u) \bigr] &\le \sum_{m = 0}^{M - 1} \E \bigl[ \ind \{ \deg_{G_m}(u) \ge w_0 \} \, \beta_{G_m}(u) \bigr]\\
&\le \frac{2}{\log 2} \, (\log n) \biggl( \frac{2 e \Lambda}{w_0} \biggr)^{w_0},
\end{split}
\]
where the last inequality is \cref{cor:high-degree} with $j = w_0 \ge 1$.

Inserting the last two displays into \eqref{eq:high-descent} and using $(8/3) \cdot 2 \Lambda \cdot (2 / \log 2) = 32 \Lambda / (3 \log 2)$, we obtain \eqref{eq:high-main}.

We now show that the right side of \eqref{eq:high-main} is small for large $n$. We have
\[
D_n + 1 - w_0 \ge \frac{D_n}{2} + 1 \ge \frac{(\log n)^{1/4}}{2}, \qquad w_0 \ge \frac{D_n - 1}{2} \ge \frac{(\log n)^{1/4}}{2} - 1,
\]
where we used $D_n \ge (\log n)^{1/4} - 1$ and $(D_n - 1)/2 \le w_0 \le D_n/2$.

By the first bound, the first term in \eqref{eq:high-main} is bounded by $(8/3) \, 2^{-(\log n)^{1/4}}$. The second implies that $w_0 \to \infty$. We take $n$ large enough that $w_0 \ge \max(K, 2 e^2 \Lambda)$. Then $\log(w_0 / (2 e \Lambda)) \ge 1$, which gives
\[
\biggl( \frac{2 e \Lambda}{w_0} \biggr)^{w_0} = \exp \biggl( - w_0 \log \frac{w_0}{2 e \Lambda} \biggr) \le e^{-w_0} \le e^{1 - (\log n)^{1/4} / 2}.
\]
Using also $w_0 - K + 2 \ge 2$ and $(n/(n - 1))^2 \le 9/4$, we bound the second term in \eqref{eq:high-main} by $(12 e \Lambda / \log 2) (\log n) e^{-(\log n)^{1/4} / 2}$. As $x^{5/2} e^{-x^{1/4}/2} \to 0$ and $x^{3/2} 2^{-x^{1/4}} \to 0$ as $x \to \infty$, each term is bounded by $(\log n)^{-3/2} / 2$ for all large $n$, with a threshold that depends only on $\Lambda$.

Finally, let $0 \le L' \le U' \le M$. The event in \eqref{eq:high-main-window} is contained in the event in \eqref{eq:high-main}, since $H \le U' \le M$ on this event and the minimum of the two degrees is at most their maximum. This gives \eqref{eq:high-main-window}.
\end{proof}

\begin{remark}\label{rem:high-q2}
For $q = 2$, the conclusion of this section has a short direct proof. Indeed, we now show that $\kappa_j^{\rm piv} = 0$ for all $j > K$ when $q = 2$. Let $G$ be a colorable multigraph and let $u \sim_G v$ with $u \ne v$. A connected bipartite multigraph has exactly two colorings with two colors, which differ by exchanging the colors, and different components can be colored independently. This implies that $u$ and $v$ lie in the same component of $G$, at even distance. Let $e$ be an occurrence at $u$ such that $u$ and $v$ are still connected in $G - e$. Then they are at even distance in $G - e$ as well, since every path in $G - e$ is a path in the bipartite multigraph $G$, and therefore $u \sim_{G - e} v$. A pivotal occurrence at $u$ is thus an occurrence of a pair that lies on every path from $u$ to $v$. We fix one such path without repeated vertices. Its only edge at $u$ is its first edge, so every pivotal occurrence at $u$ is an occurrence of this pair. Since a pair with a pivotal occurrence has multiplicity one, at most one occurrence at $u$ is pivotal for $(u, v)$. If $\deg_G(u) = j > K$, then at least $j - 1 \ge K$ occurrences at $u$ are non-pivotal, and $\kappa_j^{\rm piv}(G) = 0$. If $w_0 \ge K$, all pivotal terms in \cref{prop:high-descent} then vanish, which leaves the bound $(8/3) \, 4^{-(D_n + 1 - w_0)}$ without the potential argument. In particular, the proof of \cref{thm:high-main} does not use $q \ge 3$.
\end{remark}

\chapter{Proof of the main theorems}\label{sec:proof}

In this section, we prove the window bound of \cref{thm:proof-window} and deduce \cref{thm:intro-main,thm:intro-window} from it. The window bound combines the reduction of \cref{sec:red} with the estimates of \cref{sec:low,sec:high}, and the deduction uses the convergence criterion of \cref{thm:red-convergence}, which is proved in \cref{app:conv}. We end with three remarks on the scope of the argument.

\section{Parameters and the window bound}\label{sec:proof-window}

Throughout this section, we retain the fixed integer $q \ge 3$. We now set $\Lambda = \gamma_q + 2$. With this choice, we recall the parameter definitions:
\begin{equation}\label{eq:proof-params}
\begin{gathered}
\gamma_q = \frac{\log q}{\log (q / (q - 1))}, \qquad \Lambda = \gamma_q + 2, \qquad M = \lfloor \Lambda n \rfloor, \qquad \eta_0 = \frac{e^{-2 \Lambda}}{8}, \\
A_q = 2q \log q + 2q + 1, \qquad K = \lceil 8 \Lambda \rceil, \qquad \eps_1 = \eps_2 = 1, \qquad x_n = n \eps_n^2, \\
\eps_n = (\log n)^{-3/2}, \qquad k_n = \biggl\lceil \frac{\log n}{\log \log n} \biggr\rceil, \qquad D_n = \lfloor (\log n)^{1/4} \rfloor, \qquad w_0 = \lfloor D_n / 2 \rfloor,
\end{gathered}
\end{equation}
where the last line is for $n \ge 3$. Each of these parameters was introduced in the section where it is first used, and we collect them here. After choosing $q$, we fix $\eta \in (0, \eta_0]$. We use the quantiles $L_n(\eta)$ and $U_n(\eta)$ and the window $w_n(\eta)$ of \cref{def:pre-quantile} for this value of $\Lambda$. Here $\gamma_q$ is the first-moment bound for the threshold, $A_q$ bounds the expected lifetime in \cref{lem:pre-lifetime}, $\eps_n$ is the target bound for the probabilities in \cref{prop:red-split}, $x_n = n \eps_n^2$, $k_n$ bounds the list lengths retained in the proof in \cref{sec:low}, $D_n$ is the degree cut, and $K$ and $w_0$ are the parameters of \cref{prop:high-descent,thm:high-main}.

Every estimate used below holds for all sufficiently large $n$, and the conditions on $n$ are of three kinds. The conditions of \cref{sec:low} consist of the threshold $n_q$ of \cref{thm:low-main} and, in \cref{cor:low}, a comparison of the negative power $n^{-1/8}$ with a power of $\log n$. Those of \cref{sec:high} use that $w_0 \log(w_0 / (2 e \Lambda))$ grows faster than $\log \log n$. The thresholds in these two sections are determined by $q$, through $\Lambda$. Finally, the bounds on the quantiles in \cref{lem:pre-centre} hold for each fixed $\eta$, and therefore the threshold in \cref{thm:proof-window} depends on $\eta$ as well.

\begin{theorem}\label{thm:proof-window}
Let $q \ge 3$, and let the parameters be given by \eqref{eq:proof-params}. For every $\eta \in (0, \eta_0]$, there exists $n_0$, depending only on $q$ and $\eta$, such that for all $n \ge n_0$,
\begin{equation}\label{eq:proof-window}
w_n(\eta) \le \frac{6 \Lambda}{\eta^2} \, n (\log n)^{-3/2}.
\end{equation}
\end{theorem}

\begin{proof}
Fix $\eta \in (0, \eta_0]$. We abbreviate $L = L_n(\eta)$ and $U = U_n(\eta)$. We will verify the hypotheses \eqref{eq:red-low} and \eqref{eq:red-high} of \cref{prop:red-split} for all large $n$.

For \eqref{eq:red-low}, recall that \cref{cor:low} gives \eqref{eq:red-low-target} for all large $n$, with a threshold depending only on $q$. The event in \eqref{eq:red-low} is contained in the event in \eqref{eq:red-low-target}, so \eqref{eq:red-low} holds for these $n$.

For \eqref{eq:red-high}, note that $0 \le L \le U \le M$ by \cref{lem:pre-quantile}. We apply \eqref{eq:high-main-window} with $L' = L$ and $U' = U$, and we use that $\eps_n = (\log n)^{-3/2}$ for $n \ge 3$. For all large $n$, we obtain
\[
\begin{split}
\Pp \Bigl( L < H \le U, \ \min_{v \in e_H} \deg_\Gamma(v) > D_n, \ r(\Gamma) > x_n \Bigr) &\le \Pp \Bigl( L < H \le U, \ \min_{v \in e_H} \deg_\Gamma(v) > D_n \Bigr)\\
&\le \eps_n,
\end{split}
\]
which is \eqref{eq:red-high}.

The proposition \cref{prop:red-split} now applies for all large $n$. It implies that $w_n(\eta) \le (6 \Lambda / \eta^2) \, n \eps_n$, which is \eqref{eq:proof-window}.
\end{proof}

\section{Proofs of the main theorems}\label{sec:proof-main}

\begin{proof}[Proof of \cref{thm:intro-main}]
We apply \cref{thm:red-convergence} with $\delta_n = \eps_n$, and we first check its hypotheses. Since $\log 3 > 1$, the sequence $(\eps_n)$ takes values in $(0, 1]$ and is nonincreasing. Further, for every integer $j \ge 2$, we have $\lfloor e^j \rfloor \ge 3$ and
\[
\lfloor e^j \rfloor \ge e^j - 1 \ge e^{j - 1}, \qquad \text{and hence} \qquad \eps_{\lfloor e^j \rfloor} \le \bigl( \log e^{j - 1} \bigr)^{-3/2} = (j - 1)^{-3/2}.
\]
Since $\eps_{\lfloor e \rfloor} = \eps_2 = 1$, it follows that
\[
\sum_{j \ge 1} \eps_{\lfloor e^j \rfloor} \le 1 + \sum_{j \ge 2} (j - 1)^{-3/2} < \infty.
\]

Next, \cref{thm:proof-window} shows that for every $\eta \in (0, \eta_0]$, the hypothesis \eqref{eq:red-window-hyp} holds with $K'(\eta) = 6 \Lambda / \eta^2$ and a suitable $n_1(\eta)$.

The hypotheses of \cref{thm:red-convergence} hold, so there exists a constant $\alpha_q \in [1/2, \gamma_q]$ with the properties (1)--(3) stated there. Properties (2) and (3) state that $G(n, \lfloor cn \rfloor)$ is $q$-colorable with probability tending to one for every constant $c \in [0, \alpha_q)$ and tending to zero for every constant $c > \alpha_q$, and that the same holds for $G(n, d/n)$ with the threshold $d_q = 2 \alpha_q$. This is \cref{thm:intro-main}.
\end{proof}

\begin{proof}[Proof of \cref{thm:intro-window}]
Let $\epsilon \in (0, 1)$. We set $\eta_1 = \min(\epsilon / 2, \eta_0)$ and $m_n = L_n(\eta_1)$. As shown in the proof of \cref{thm:intro-main}, \cref{thm:red-convergence} applies. Its property (1) gives $m_n / n \to \alpha_q$.

We now bound the probability that $H$ is far from $m_n$. Recall from \cref{lem:pre-centre}(a) that $U_n(\eta_1) < M$ for all large $n$. For such $n$, we apply \cref{lem:pre-quantile} at the level $\eta_1$, with $m = L_n(\eta_1)$ and with $m = U_n(\eta_1) + 1$, and we obtain
\[
\Pp \bigl( H \le L_n(\eta_1) \bigr) < \eta_1, \qquad \Pp \bigl( H > U_n(\eta_1) + 1 \bigr) < \eta_1.
\]
Hence, with probability at least $1 - 2 \eta_1 \ge 1 - \epsilon$, we have $L_n(\eta_1) < H \le U_n(\eta_1) + 1$, and on this event
\[
0 < H - m_n \le w_n(\eta_1) + 1.
\]

Finally, applying \cref{thm:proof-window} at the level $\eta_1 \in (0, \eta_0]$, we have $w_n(\eta_1) \le (6 \Lambda / \eta_1^2) \, n (\log n)^{-3/2}$ for all large $n$. We set $C = 6 \Lambda / \eta_1^2 + 1$. Using also $n (\log n)^{-3/2} \ge 1$ for $n \ge 3$, we conclude that for all large $n$,
\[
\Pp \bigl( \abs{H - m_n} \le C n (\log n)^{-3/2} \bigr) \ge 1 - \epsilon.
\]
This proves \cref{thm:intro-window}, with a constant $C$ determined by $q$ and $\epsilon$.
\end{proof}

\section{Remarks}\label{sec:proof-remarks}

\begin{remark}\label{rem:proof-q3}
We record where the proof uses the assumption $q \ge 3$. There are two places. The first is the part of \cref{lem:low-frozen} which shows that if a vertex lies in a frozen pair, then its neighbors use at least $q - 1$ colors in every coloring of the multigraph with this vertex removed. Since $q - 1 \ge 2$, the list of its neighbors is then never monochromatic. This is the one use of $q \ge 3$ in the proof of \cref{thm:low-main}, and the implication fails for $q = 2$ (see \cref{rem:low-q2}). The second is \cref{lem:pre-centre}(b), which uses that a multigraph whose components each contain at most one cycle is $q$-colorable when $q \ge 3$, and it gives the lower bound $\alpha_q \ge 1/2$. The remaining steps hold, with the same proofs, for every $q \ge 2$; for the estimates of \cref{sec:high}, see \cref{rem:high-q2}. Some use of $q \ge 3$ is unavoidable, since for $q = 2$ the threshold is not sharp. Indeed, for every constant $d \in (0, 1)$, the probability that $G(n, d/n)$ is bipartite converges to $e^{d/2} ((1 - d)/(1 + d))^{1/4} \in (0, 1)$; see \cite[Section~1, equation~(1.5)]{pittel2010frequently}, and \cite[p.~64]{achlioptas1999sharp} for a discussion of the case $q = 2$.
\end{remark}

\begin{remark}\label{rem:proof-sqrt}
We now explain why the variance bound of \cref{sec:red} cannot certify windows below the order $\sqrt n$. Fix $\eta \in (0, \eta_0]$. In this remark, we write $L = L_n(\eta)$ and $U = U_n(\eta)$. We use the capped variable $Y$ of \eqref{eq:red-Y} and the variables $Y^{(-i)}$ of \eqref{eq:red-Ydel} at this level; in the latter, the occurrence $e_i$ is omitted. The proof of \cref{prop:red-window} rests on the inequalities
\[
\eta^2 w_n(\eta)^2 \le \Var(Y) \le \sum_{i = 1}^U \E \bigl[ \bigl( Y^{(-i)} - Y \bigr)^2 \bigr] = \E \bigl[ H \bigl( 2 r_{U - H + 1}(\Gamma) - 1 \bigr); L < H \le U \bigr].
\]
The first inequality is \cref{lem:red-varlower}, the second is the Efron--Stein inequality used in the proof of \cref{lem:red-deletion}, and the identity is \eqref{eq:red-deletion}. On the event $\{ L < H \le U \}$, the multigraph $\Gamma$ is colorable and $U - H + 1 \ge 1$, so $r_{U - H + 1}(\Gamma) \ge 1$ by \cref{lem:pre-lifetime}(a). Hence the right side is at least $\E [H \ind \{ L < H \le U \}] \ge (L + 1) \Pp(L < H \le U)$, and \cref{lem:pre-centre}(b) with $c_0 = 1/4$ implies that $L > n/4$ for all large $n$. Thus, whenever $\Pp(L < H \le U)$ is bounded below by a positive constant, the right side of \eqref{eq:red-deletion} is at least of order $n$, and the bound on $w_n(\eta)$ that these inequalities can yield is at least of order $\sqrt n$. This is a limitation of the variance bound, not a lower bound on the window. For comparison, Wilson proved that the transition for $q$-colorability has width at least of order $\sqrt n$ in the number of edges \cite[Corollary~6]{wilson2002critical}. The window $n (\log n)^{-3/2}$ of \cref{thm:proof-window} is much larger than this scale.
\end{remark}

\begin{remark}\label{rem:proof-finite}
\Cref{thm:intro-main,thm:intro-window} are asymptotic statements, and the proof gives no useful bounds at practical sizes of $n$. The conditions on $n$ in \cref{sec:low,sec:high} are explicit inequalities, but they hold only for very large $n$. For example, the hypothesis $w_0 \ge K$ of \cref{thm:high-main} is equivalent to $(\log n)^{1/4} \ge 2K$, so it requires $\log n \ge (2K)^4 \ge (16 \Lambda)^4$. Moreover, \cref{thm:red-convergence} identifies $\alpha_q$ only as a limit. It does not provide a rate for the convergence of $L_n(\eta)/n$ to $\alpha_q$, and the proof gives no information on the value of $\alpha_q$ beyond the bounds $1/2 \le \alpha_q \le \gamma_q$.
\end{remark}

\appendix
\chapter{Convergence of the threshold location}\label[appendix]{app:conv}

In this appendix, we prove \cref{thm:red-convergence}, and we first recall its statement. The parameter $\Lambda > \gamma_q + 1$ is fixed as in \cref{sec:pre-process}, and $\eta_0 = e^{-2\Lambda}/8$. Let $(\delta_n)_{n \ge 1}$ be a nonincreasing sequence with values in $(0, 1]$ such that
\begin{equation}\label{eq:conv-summable}
\sum_{k \ge 1} \delta_{\lfloor e^k \rfloor} < \infty,
\end{equation}
and suppose that for every $\eta \in (0, \eta_0]$ there exist $K'(\eta) > 0$ and an integer $n_1(\eta)$ such that
\begin{equation}\label{eq:conv-window}
w_n(\eta) \le K'(\eta) n \delta_n \qquad \text{for all } n \ge n_1(\eta).
\end{equation}
\Cref{thm:red-convergence} asserts that, under these hypotheses, there exists $\alpha_q \in [1/2, \gamma_q]$ with the following properties. For every $\eta \in (0, \eta_0]$, we have $L_n(\eta)/n \to \alpha_q$ and $U_n(\eta)/n \to \alpha_q$ as $n \to \infty$. For every real $c \ge 0$, the probabilities $\Pp(H > \lfloor cn \rfloor)$ and $\Pp(G(n, \lfloor cn \rfloor) \text{ is colorable})$ tend to $1$ if $c < \alpha_q$ and to $0$ if $c > \alpha_q$. For every real $d \ge 0$, the probability $\Pp(G(n, d/n) \text{ is colorable})$ tends to $1$ if $d < 2\alpha_q$ and to $0$ if $d > 2\alpha_q$.

The method adapts the proof of the corresponding statement for random $k$-SAT in \cite[Section~5]{openai2026limiting}. Its main input is a product inequality for colorability probabilities at different system sizes, which is the coloring case of the interpolation of Bayati, Gamarnik, and Tetali \cite[Theorem~5 and Proposition~2]{bayati2013combinatorial}, written for a Poisson number of edges as in \cite[Lemma~5.2]{openai2026limiting}. Translated to quantiles, it shows that the normalized lower quantile $L_n(\eta)/n$ at size $N_1 + N_2$ is at least the smaller of its values at sizes $N_1$ and $N_2$, up to an error. A balanced binary tree of block sizes then shows that such comparisons in minimum form imply convergence, as in \cite[Lemma~5.4]{openai2026limiting}. In \cite{openai2026limiting}, the comparisons concern normalized expectations and have polynomially small errors. Our window bound concerns quantiles and gives only errors of order $\delta_n$. We compare quantiles using errors that are summable on geometric scales, as in the classical theory of nearly subadditive sequences \cite[Theorem~23]{debruijn1952some} and in \cite[arXiv:1006.3786v1, Section~4.3, Lemma~8 and Remark~3]{abbe2014concentration}. This summability cannot be dropped from the tree argument (see \cref{rem:conv-necessity}). In the last step, we use the window bound at every small quantile level. All quantiles then share one limit, and the zero-one laws follow without the sharp-threshold theorem of Achlioptas and Friedgut \cite{achlioptas1999sharp} (see \cref{rem:red-af}).

In \cref{sec:conv-interp}, we prove the product inequality for a Poisson model with loops, and in \cref{sec:conv-quantile} we relate this model to the quantiles of $V$. \Cref{sec:conv-superadd} combines the two to compare the normalized lower quantiles at different sizes, and \cref{sec:conv-tree} shows that these comparisons imply convergence. \Cref{sec:conv-transfer} relates the edge process to $G(n, m)$ and $G(n, p)$, and we complete the proof in \cref{sec:conv-proof}.

\section{The product inequality}\label{sec:conv-interp}

In this appendix, we also use multigraphs with loops. For a finite set $S$, we define a \emph{multigraph with loops on $S$} to be a finite multiset of sets $\{x, y\}$ with $x, y \in S$, allowing $x = y$. We call its elements, counted with multiplicity, \emph{occurrences}, and an occurrence $\{x, x\}$ a \emph{loop}. We define its colorings and colorability by the same inequality as in \cref{def:pre-multigraph}, and let $\Col(G)$ denote its set of colorings. For $x, y \in S$, we define $G + \{x, y\}$ to add one occurrence, including when $x = y$. For $S' \subset S$, we define $G[S']$ by retaining the occurrences with both endpoints in $S'$. In particular, a multigraph with loops that contains a loop is not colorable, and a colorable multigraph with loops is a multigraph in the sense of \cref{def:pre-multigraph}. When $|S| \ge 2$, its frozen classes are defined by \cref{def:pre-frozen}; this includes every graph whose frozen classes are used in the interpolation proof below.

\begin{definition}\label{def:conv-poisson}
Let $n \ge 1$ be an integer and $c \ge 0$ a real number. Let $(u_i, v_i)_{i \ge 1}$ be independent random elements of $[n] \times [n]$ with the uniform law, and let $I$ be an independent random variable with the $\Pois(cn)$ law. We let $A_n(c)$ denote the multigraph with loops on $[n]$ whose occurrences are $\{u_i, v_i\}$ for $1 \le i \le I$, and we set
\begin{equation}\label{eq:conv-rho}
\rho_n(c) = \Pp\big(A_n(c) \text{ is colorable}\big).
\end{equation}
\end{definition}

For $n = 1$, every occurrence of $A_1(c)$ is a loop, so that
\begin{equation}\label{eq:conv-rho1}
\rho_1(c) = \Pp(I = 0) = e^{-c}.
\end{equation}

\begin{lemma}\label{lem:conv-interp}
For all integers $N_1, N_2 \ge 1$ and all real $c \ge 0$, we have
\begin{equation}\label{eq:conv-product}
\rho_{N_1 + N_2}(c) \ge \rho_{N_1}(c) \rho_{N_2}(c).
\end{equation}
\end{lemma}

\begin{proof}
Both sides of \eqref{eq:conv-product} equal $1$ when $c = 0$. We may assume that $c > 0$. Let $N = N_1 + N_2$, $S_1 = [N_1]$, and $S_2 = [N] \setminus [N_1]$. Let $(\omega_i)_{i \ge 1}$ and $(\omega_i^{(j)})_{i \ge 1}$, $j = 1, 2$, be three independent sequences of independent random elements. We give the entries of the first sequence the uniform law on $[N] \times [N]$ and those of sequence $j$ the uniform law on $S_j \times S_j$.

For $\boldsymbol\mu = (\mu_0, \mu_1, \mu_2)$ with $\mu_0, \mu_1, \mu_2 \ge 0$, let $I_0$, $I_1$, and $I_2$ be independent random variables with the $\Pois(\mu_0)$, $\Pois(\mu_1)$, and $\Pois(\mu_2)$ laws, independent of these sequences. We let $A^*(\boldsymbol\mu)$ denote the multigraph with loops on $[N]$ whose occurrences are the sets $\{x, y\}$ for the elements $(x, y)$ of the lists $(\omega_i)_{i \le I_0}$, $(\omega^{(1)}_i)_{i \le I_1}$, and $(\omega^{(2)}_i)_{i \le I_2}$, and we set
\[
\bar\chi(\boldsymbol\mu) = \Pp\big(A^*(\boldsymbol\mu) \text{ is colorable}\big).
\]
For $\upsilon \in [0, c]$, we define the interpolating multigraph and its colorability probability by
\begin{equation}\label{eq:conv-interpolation}
A^{(\upsilon)} = A^*\big(N\upsilon, N_1 (c - \upsilon), N_2 (c - \upsilon)\big), \qquad \chi(\upsilon) = \bar\chi\big(N\upsilon, N_1 (c - \upsilon), N_2 (c - \upsilon)\big).
\end{equation}
The occurrences contributed by $(\omega_i)$ have mean number $N\upsilon$, and those contributed by $(\omega_i^{(j)})$ have mean number $N_j(c - \upsilon)$ for $j = 1, 2$.

We first identify the values of $\chi$ at the endpoints. The multigraph $A^{(0)}$ has no occurrence joining $S_1$ and $S_2$, so it is colorable if and only if $A^{(0)}[S_1]$ and $A^{(0)}[S_2]$ are colorable. These two multigraphs are independent, and after relabeling $S_2$ as $[N_2]$, they have the laws of $A_{N_1}(c)$ and $A_{N_2}(c)$. The multigraph $A^{(c)}$ has the law of $A_N(c)$. Therefore
\[
\chi(0) = \rho_{N_1}(c) \rho_{N_2}(c), \qquad \chi(c) = \rho_N(c),
\]
and it suffices to show that $\chi$ is nondecreasing on $[0, c]$.

Next, we compute the partial derivatives of $\bar\chi$. For $k = (k_0, k_1, k_2)$ with nonnegative integer entries, let $a_k$ be the probability that the multigraph with loops formed by $(\omega_i)_{i \le k_0}$, $(\omega^{(1)}_i)_{i \le k_1}$, and $(\omega^{(2)}_i)_{i \le k_2}$ is colorable. Then
\[
\bar\chi(\boldsymbol\mu) = \sum_{k} \pi_{\mu_0}(k_0) \pi_{\mu_1}(k_1) \pi_{\mu_2}(k_2) a_k = e^{-\mu_0 - \mu_1 - \mu_2} \sum_{k} \frac{\mu_0^{k_0} \mu_1^{k_1} \mu_2^{k_2}}{k_0! \, k_1! \, k_2!} a_k.
\]
Since $|a_k| \le 1$, the last series converges absolutely on $\R^3$. We extend $\bar\chi$ to $\R^3$ by the right side of the preceding display. This extension is smooth, and its partial derivatives may be computed term by term.

Let $(x, y)$ be a uniformly random element of $[N] \times [N]$, independent of $A^*(\boldsymbol\mu)$. The derivative of $\mu_0 \mapsto \pi_{\mu_0}(k_0)$ is $\pi_{\mu_0}(k_0 - 1) - \pi_{\mu_0}(k_0)$. Using this and shifting the index $k_0$ in the terms with $\pi_{\mu_0}(k_0 - 1)$, we obtain
\[
\begin{aligned}
\partial_{\mu_0} \bar\chi(\boldsymbol\mu) &= \sum_{k} \pi_{\mu_0}(k_0) \pi_{\mu_1}(k_1) \pi_{\mu_2}(k_2) \big(a_{(k_0 + 1, k_1, k_2)} - a_k\big)\\
&= -\Pp\big(A^*(\boldsymbol\mu) \text{ is colorable}, \ A^*(\boldsymbol\mu) + \{x, y\} \text{ is not colorable}\big),
\end{aligned}
\]
In the second step, we used that adding an occurrence can only destroy colorability. Repeating this computation for $\mu_1$ and $\mu_2$, we find the analogous formulas, with $(x, y)$ uniform on $S_j \times S_j$. (This is the perturbation formula for Poisson processes \cite[Theorem~19.1]{last2018lectures} in the case that we need.)

We now evaluate these probabilities. Let $A$ be a colorable multigraph with loops on $[N]$. Then $A$ has no loops and is a multigraph on $[N]$. We claim that for $x, y \in [N]$, the multigraph with loops $A + \{x, y\}$ is not colorable if and only if $x \sim _A y$. Indeed, for $x = y$, the multigraph $A + \{x, x\}$ contains a loop, and $x \sim _A x$ holds since $A$ is colorable. For $x \neq y$, the claim is \cref{lem:pre-frozen}(a). Let $C_1, \dots, C_k$ denote the frozen classes of $A$. It follows that the conditional probability that $A + \{x, y\}$ is not colorable is
\[
\begin{aligned}
&\sum_{i = 1}^k \Big(\frac{|C_i|}{N}\Big)^2 &&\text{if } (x, y) \text{ is uniform on } [N] \times [N], \\
&\sum_{i = 1}^k \Big(\frac{|C_i \cap S_j|}{N_j}\Big)^2 &&\text{if } (x, y) \text{ is uniform on } S_j \times S_j,
\end{aligned}
\]
where the sums run over these $k$ classes.

On the event that $A^{(\upsilon)}$ is colorable, let $C_1, \dots, C_k$ now denote its frozen classes. We define the integrand below to be zero otherwise. Applying the chain rule to \eqref{eq:conv-interpolation} and inserting these formulas, we obtain
\begin{equation}\label{eq:conv-derivative}
\begin{aligned}
\chi'(\upsilon) &= N \partial_{\mu_0} \bar\chi - N_1 \partial_{\mu_1} \bar\chi - N_2 \partial_{\mu_2} \bar\chi\\
&= \E\bigg[\ind\{A^{(\upsilon)} \text{ is colorable}\} \sum_{i = 1}^k \bigg(\frac{|C_i \cap S_1|^2}{N_1} + \frac{|C_i \cap S_2|^2}{N_2} - \frac{|C_i|^2}{N}\bigg)\bigg],
\end{aligned}
\end{equation}
where the partial derivatives are evaluated at $(N\upsilon, N_1 (c - \upsilon), N_2 (c - \upsilon))$. For integers $a, b \ge 0$, we have
\begin{equation}\label{eq:conv-square}
\frac{a^2}{N_1} + \frac{b^2}{N_2} - \frac{(a + b)^2}{N_1 + N_2} = \frac{(N_2 a - N_1 b)^2}{N_1 N_2 (N_1 + N_2)} \ge 0,
\end{equation}
as one checks by multiplying both sides by $N_1 N_2 (N_1 + N_2)$. Applying \eqref{eq:conv-square} with $a = |C_i \cap S_1|$ and $b = |C_i \cap S_2|$ in \eqref{eq:conv-derivative}, we find that $\chi'(\upsilon) \ge 0$ for all $\upsilon \in [0, c]$. Hence
\[
\rho_N(c) = \chi(c) \ge \chi(0) = \rho_{N_1}(c) \rho_{N_2}(c). \qedhere
\]
\end{proof}

System-splitting interpolation was used by Guerra and Toninelli to prove thermodynamic limits for mean-field spin glasses \cite[Section~3, Theorems~1--2]{guerra2002thermodynamic}. Franz and Leone \cite[Section~III and appendices]{franz2003replicaa} and Panchenko and Talagrand \cite[Sections~1--2]{panchenko2004bounds} developed free-energy bounds for diluted models under their stated interaction and arity hypotheses. For the coloring problem, Bayati, Gamarnik, and Tetali compare the colorability probability of the global model with a given number of edges with that of the two parts \cite[Theorem~5]{bayati2013combinatorial}, and their proof of \cite[Proposition~2]{bayati2013combinatorial} uses, as above, the classes of vertices that receive the same color in every coloring. Averaging their comparison over a Poisson number of edges gives \eqref{eq:conv-product}. The Poisson form and its proof by differentiation in the means are those of \cite[Lemma~5.2]{openai2026limiting} for $k$-SAT, where the sign of the derivative comes from the convexity of $x \mapsto x^k$. Here it comes from \eqref{eq:conv-square}.

\section{The Poisson model and the quantiles}\label{sec:conv-quantile}

We now relate $\rho_n(c)$ to the hitting time $H$ of the edge process on $[n]$. The proof uses standard concentration bounds for Poisson random variables. For integers $n \ge 1$, we set
\begin{equation}\label{eq:conv-Delta}
\Delta_n = \lceil \sqrt n \log n \rceil.
\end{equation}

\begin{lemma}\label{lem:conv-quantile}
\begin{enumerate}[label=(\alph*)]
\item For all integers $n \ge 2q$ and all real $c \ge 0$, we have
\begin{equation}\label{eq:conv-rho-H}
\rho_n(c) = e^{-c} \sum_{j \ge 0} \pi_{c(n - 1)}(j) \Pp(H > j).
\end{equation}
\item Let $\eta \in (0, 1/2]$. For all large $n$ and all $c \in (0, \Lambda]$, the following holds. If $c(n - 1) \le L_n(\eta) - \Delta_n$, then
\begin{equation}\label{eq:conv-below}
\rho_n(c) \ge e^{-\Lambda} (1 - \eta) \big(1 - e^{-(\log n)^2/(2\Lambda)}\big),
\end{equation}
and if $c(n - 1) \ge U_n(\eta) + \Delta_n$, then
\begin{equation}\label{eq:conv-above}
\rho_n(c) \le \eta + e^{-(\log n)^2/(2\Lambda)}.
\end{equation}
\end{enumerate}
\end{lemma}

\begin{proof}
We begin with (a). Each occurrence of $A_n(c)$ is a loop with probability $1/n$. We apply \cref{lem:pre-split}(a) with $E = [n] \times [n]$, $\nu$ the uniform law, $\mu = cn$, and the partition of $E$ into the diagonal $\{(x, x) : x \in [n]\}$ and its complement. It shows that the number of loops of $A_n(c)$ has the $\Pois(c)$ law, and that, independently, the elements $(u_i, v_i)$ with $u_i \neq v_i$ form a list of $\Pois(c(n - 1))$ independent uniformly random elements of $\{(x, y) \in [n]^2 : x \neq y\}$. For such a uniformly random element $(x, y)$, the set $\{x, y\}$ is a uniformly random pair of $[n]$.

Conditionally on the number $j$ of occurrences that are not loops, the multigraph formed by these occurrences has the law of $G_j$. The multigraph with loops $A_n(c)$ is colorable if and only if it has no loop and this multigraph is colorable. We have
\[
\rho_n(c) = e^{-c} \sum_{j \ge 0} \pi_{c(n - 1)}(j) \Pp(G_j \text{ is colorable}),
\]
which is \eqref{eq:conv-rho-H}.

We now prove (b). We set $\mu = c(n - 1)$. Let $J$ be a random variable with the $\Pois(\mu)$ law. We write $L = L_n(\eta)$ and $U = U_n(\eta)$. Suppose first that $\mu \le L - \Delta_n$. Recall from \cref{lem:pre-quantile} that $\Pp(H > j) > 1 - \eta$ for all $j \le L$. Inserting this into \eqref{eq:conv-rho-H}, we find
\[
\rho_n(c) \ge e^{-c} (1 - \eta) \sum_{j \le L} \pi_\mu(j) = e^{-c} (1 - \eta) \Pp(J \le L) \ge e^{-\Lambda} (1 - \eta) \big(1 - \Pp(J \ge \mu + \Delta_n)\big),
\]
where we used $c \le \Lambda$ and the inclusion $\{J > L\} \subset \{J \ge \mu + \Delta_n\}$. Further, we have $\mu + \Delta_n \le L \le M \le \Lambda n$ and $\Delta_n^2 \ge n (\log n)^2$, so \cref{lem:pre-tails}(c) gives
\[
\Pp(J \ge \mu + \Delta_n) \le \exp\Big(-\frac{\Delta_n^2}{2(\mu + \Delta_n)}\Big) \le \exp\Big(-\frac{n (\log n)^2}{2\Lambda n}\Big) = e^{-(\log n)^2/(2\Lambda)}.
\]
This proves \eqref{eq:conv-below}.

Suppose now that $\mu \ge U + \Delta_n$. By \cref{lem:pre-centre}(a), we have $U + 1 \le M$ for all large $n$. For such $n$, \cref{lem:pre-quantile} and the monotonicity of $\Pp(H > j)$ in $j$ give $\Pp(H > j) < \eta$ for all $j \ge U + 1$. We bound $\Pp(H > j)$ by $1$ for $j \le U$ and by $\eta$ for $j > U$ in \eqref{eq:conv-rho-H}, and we use $e^{-c} \le 1$. Since $0 < \mu \le \Lambda n$, \cref{lem:pre-tails}(c) then gives
\[
\rho_n(c) \le \Pp(J \le U) + \eta \le \Pp(J \le \mu - \Delta_n) + \eta \le \exp\Big(-\frac{\Delta_n^2}{2\mu}\Big) + \eta \le e^{-(\log n)^2/(2\Lambda)} + \eta,
\]
which is \eqref{eq:conv-above}. The only condition on $n$ that we used is $U + 1 \le M$, which does not involve $c$.
\end{proof}

\section{Comparison of the quantiles at different sizes}\label{sec:conv-superadd}

In this section, we fix $\eta \in (0, \eta_0]$ and assume \eqref{eq:conv-window} at this level. We write $K' = K'(\eta)$ and $n_1 = n_1(\eta)$. For each integer $n \ge 2q$, we set $\xi_n = L_n(\eta)/n$. Recall from \cref{lem:pre-quantile} that $L_n(\eta) \le M$, so that $0 \le \xi_n \le \Lambda$. We extend the sequence $(\delta_n)$ to real arguments by $\delta(u) = \delta_{\lfloor u \rfloor}$ for $u \ge 1$, and we define
\begin{equation}\label{eq:conv-g}
g(u) = K' \delta(u) + (\Lambda + 6) \frac{\log u}{\sqrt u}, \qquad u \ge 1.
\end{equation}

\begin{lemma}\label{lem:conv-superadd}
There exists an integer $n_0 \ge 2q$ such that for all integers $N_1, N_2 \ge n_0$, we have
\begin{equation}\label{eq:conv-superadd}
\xi_{N_1 + N_2} \ge \min(\xi_{N_1}, \xi_{N_2}) - g\big(\min(N_1, N_2)\big), \qquad \xi_{N_1 + 1} \ge \xi_{N_1} - g(N_1).
\end{equation}
\end{lemma}

The integer $n_0$ depends on $q$, $\Lambda$, $\eta$, $K'$, $n_1$, and the sequence $(\delta_n)$.

\begin{proof}
We first choose $n_0$. For each integer $v \ge 1$, we have $\Delta_v \le \sqrt v\log v + 1$, and hence
\begin{equation}\label{eq:conv-Delta-bound}
\frac{\Delta_v}{v} \le \frac{\log v}{\sqrt v} + \frac 1v,
\end{equation}
where the right side is nonincreasing in $v \ge e^2$. Let $\vartheta_n = e^{-(\log n)^2/(2\Lambda)}$. Using \cref{lem:pre-centre}(b) with $c_0 = 1/4$, the bound \eqref{eq:conv-Delta-bound}, and the fact that $\vartheta_n \to 0$, we may choose $n_0 \ge \max(2q, n_1, e^2)$ such that the conclusion of \cref{lem:conv-quantile}(b) holds at every size $n \ge n_0$, and such that
\begin{equation}\label{eq:conv-s0}
\xi_n > \frac 14, \qquad \frac{2\Delta_n}{n} < \frac 14, \qquad \vartheta_n \le \min\Big(\eta, \frac 18\Big) \qquad \text{for all } n \ge n_0.
\end{equation}

We begin with the first inequality in \eqref{eq:conv-superadd}. Let $N_1, N_2 \ge n_0$, let $N = N_1 + N_2$ and $N_{\min} = \min(N_1, N_2)$, and set
\[
c = \min\Big(\xi_{N_1} - \frac{2\Delta_{N_1}}{N_1}, \ \xi_{N_2} - \frac{2\Delta_{N_2}}{N_2}\Big).
\]
The bounds \eqref{eq:conv-s0} give $c > 0$, and we have $c \le \xi_{N_1} \le \Lambda$. For $j = 1, 2$, we also have
\[
c (N_j - 1) \le c N_j \le L_{N_j}(\eta) - 2\Delta_{N_j} \le L_{N_j}(\eta) - \Delta_{N_j}.
\]
Therefore \eqref{eq:conv-below} applies at size $N_j$, and together with $\eta \le \eta_0 \le 1/8$ and \eqref{eq:conv-s0} it gives
\[
\rho_{N_j}(c) \ge e^{-\Lambda} (1 - \eta)(1 - \vartheta_{N_j}) \ge e^{-\Lambda} \Big(\frac 78\Big)^2.
\]
By \cref{lem:conv-interp}, we obtain
\begin{equation}\label{eq:conv-product-lower}
\rho_N(c) \ge \rho_{N_1}(c) \rho_{N_2}(c) \ge e^{-2\Lambda} \Big(\frac 78\Big)^4 > \frac{e^{-2\Lambda}}{2}.
\end{equation}

However, recall that $\eta_0 = e^{-2\Lambda}/8$. Together with \eqref{eq:conv-s0}, this gives
\[
\eta + \vartheta_N \le 2\eta \le 2\eta_0 = \frac{e^{-2\Lambda}}{4}.
\]
Suppose that $c(N - 1) \ge U_N(\eta) + \Delta_N$. Then \eqref{eq:conv-above} at size $N \ge n_0$ would give $\rho_N(c) \le e^{-2\Lambda}/4$, which contradicts \eqref{eq:conv-product-lower}. Hence $c(N - 1) < U_N(\eta) + \Delta_N$. Since $N \ge n_1$, the hypothesis \eqref{eq:conv-window} gives
\[
U_N(\eta) = L_N(\eta) + w_N(\eta) \le N \xi_N + K' N \delta_N,
\]
and dividing by $N$, we find
\begin{equation}\label{eq:conv-xiN}
\xi_N > c - \frac cN - K' \delta_N - \frac{\Delta_N}{N} \ge c - \frac{\Lambda}{N} - K' \delta_N - \frac{\Delta_N}{N}.
\end{equation}

We now express the errors through $N_{\min}$. We have $N_1, N_2 \ge N_{\min} \ge e^2$, and the right side of \eqref{eq:conv-Delta-bound} is nonincreasing on $[e^2, \infty)$, so
\[
c \ge \min(\xi_{N_1}, \xi_{N_2}) - 2 \frac{\log N_{\min}}{\sqrt{N_{\min}}} - \frac{2}{N_{\min}}.
\]
Similarly, using $N \ge N_{\min}$, we have $\Delta_N/N \le (\log N_{\min})/\sqrt{N_{\min}} + 1/N_{\min}$ and $\Lambda/N \le \Lambda/N_{\min}$, and $\delta_N \le \delta(N_{\min})$ since $(\delta_n)$ is nonincreasing. Inserting these bounds into \eqref{eq:conv-xiN}, we obtain
\[
\xi_N \ge \min(\xi_{N_1}, \xi_{N_2}) - K' \delta(N_{\min}) - 3 \frac{\log N_{\min}}{\sqrt{N_{\min}}} - \frac{\Lambda + 3}{N_{\min}} \ge \min(\xi_{N_1}, \xi_{N_2}) - g(N_{\min}),
\]
where in the last step we used $1/N_{\min} \le (\log N_{\min})/\sqrt{N_{\min}}$ for $N_{\min} \ge e$. This is the first inequality in \eqref{eq:conv-superadd}.

For the second inequality, let $N_1 \ge n_0$ and $c = \xi_{N_1} - 2\Delta_{N_1}/N_1$, so that $0 < c \le \Lambda$ as before. Using \cref{lem:conv-interp} with $N_2 = 1$, the identity \eqref{eq:conv-rho1}, and \eqref{eq:conv-below} at size $N_1$, we have
\[
\rho_{N_1 + 1}(c) \ge \rho_{N_1}(c) e^{-c} \ge e^{-2\Lambda} \Big(\frac 78\Big)^2 > \frac{e^{-2\Lambda}}{2}.
\]
The argument above, at size $N_1 + 1 \ge n_0$ in place of $N$, shows that $c N_1 < U_{N_1 + 1}(\eta) + \Delta_{N_1 + 1}$, and that
\[
\begin{aligned}
\xi_{N_1 + 1} &\ge c - \frac{\Lambda}{N_1 + 1} - K' \delta_{N_1 + 1} - \frac{\Delta_{N_1 + 1}}{N_1 + 1}\\
&\ge \xi_{N_1} - K' \delta(N_1) - 3 \frac{\log N_1}{\sqrt{N_1}} - \frac{\Lambda + 3}{N_1} \ge \xi_{N_1} - g(N_1).
\end{aligned}
\]
This completes the proof.
\end{proof}

The use of the one-vertex model \eqref{eq:conv-rho1} for the comparison of adjacent sizes follows \cite[Section~5.3]{openai2026limiting}. This comparison cannot be omitted, since \cref{rem:conv-necessity} gives a sequence that satisfies the first inequality in \eqref{eq:conv-superadd} with $g = 0$ but does not converge.

\section{A balanced tree of blocks}\label{sec:conv-tree}

We now prove a general convergence lemma from the two inequalities in \eqref{eq:conv-superadd}. The proof applies those inequalities along a binary tree. The next statement specifies the condition on the error function.

\begin{lemma}\label{lem:conv-tree}
Let $u_0 \ge 1$, and let $g \colon [1, \infty) \to (0, \infty)$ be a function that is nonincreasing on $[u_0, \infty)$ and satisfies $\sum_{k \ge 1} g(e^k) < \infty$. Let $n_0 \ge 1$ be an integer, and let $(\xi_n)_{n \ge n_0}$ be a bounded sequence of real numbers such that the two inequalities in \eqref{eq:conv-superadd} hold for all integers $N_1, N_2 \ge n_0$. Then $(\xi_n)$ converges.
\end{lemma}

\begin{proof}
We first record the consequences of the summability of $g$. Let $u_1 = e^{\lceil \log u_0 \rceil}$. For every integer $k$ with $e^k \ge u_1$, the monotonicity of $g$ gives
\[
\int_{e^k}^{e^{k + 1}} g(v) \frac{\dd v}{v} \le g(e^k),
\]
and summing over these $k$, we find
\begin{equation}\label{eq:conv-integral}
\int_{u_1}^\infty g(v) \frac{\dd v}{v} < \infty.
\end{equation}
For $u \ge 2u_1$, the monotonicity of $g$ also gives
\[
g(u) \log 2 \le \int_{u/2}^{u} g(v) \frac{\dd v}{v},
\]
so \eqref{eq:conv-integral} implies that $g(u) \to 0$ as $u \to \infty$.

For real $\ell \ge 4u_0$, we define the path-error sum
\begin{equation}\label{eq:conv-pathsum}
\mathcal R(\ell) = \sum_{i \ge 1} g\Big(\frac \ell4 \Big(\frac 43\Big)^i\Big).
\end{equation}
For every $i \ge 1$, all points of the interval between $(\ell/4)(4/3)^{i - 1}$ and $(\ell/4)(4/3)^i$ are at least $u_0$, so the monotonicity of $g$ gives
\[
g\Big(\frac \ell4 \Big(\frac 43\Big)^i\Big) \log \frac 43 \le \int_{(\ell/4)(4/3)^{i - 1}}^{(\ell/4)(4/3)^i} g(v) \frac{\dd v}{v}.
\]
Summing over $i$, we obtain
\[
\mathcal R(\ell) \log \frac 43 \le \int_{\ell/4}^\infty g(v) \frac{\dd v}{v},
\]
and by \eqref{eq:conv-integral} the right side tends to zero as $\ell \to \infty$.

Next, we construct the tree. Fix an integer $\ell \ge \max(n_0, 4u_0, 2)$, which we call the block size, and let $n \ge \ell^2$. We write $n = k\ell + j$ with $k = \lfloor n/\ell \rfloor$ and $0 \le j < \ell$. Since $k \ge \ell > j$, we can write $n$ as the sum of $j$ blocks of size $\ell + 1$ and $k - j$ blocks of size $\ell$.

We construct a rooted binary tree whose leaves correspond to the $k$ blocks. For each node $\nu$, we define $k_\nu$ to be the number of blocks below $\nu$, and $N_\nu$ to be the sum of their sizes. At every node with $k_\nu \ge 2$, we split those blocks into two children containing $\lceil k_\nu/2\rceil$ and $\lfloor k_\nu/2\rfloor$ blocks (see \cref{fig:conv-tree}). Thus the root has size $n$, every leaf has size $\ell$ or $\ell + 1$, and the size of an internal node is the sum of the sizes of its two children. We let $\operatorname{ht}(\nu)$ denote the height of $\nu$, the largest number of edges on a path in the tree from $\nu$ down to a leaf.

\begin{figure}[t]
\centering
\begin{tikzpicture}[
  level distance=13mm,
  every node/.style={font=\small},
  inner/.style={draw, circle, minimum size=7mm, inner sep=0pt},
  leaf/.style={draw, rectangle, minimum width=9mm, minimum height=6mm, inner sep=1pt},
  level 1/.style={sibling distance=52mm},
  level 2/.style={sibling distance=26mm},
  level 3/.style={sibling distance=14mm}]
\node[inner] (root) {$5$}
  child {node[inner] {$3$}
    child {node[inner] {$2$}
      child {node[leaf] {$\ell + 1$}}
      child {node[leaf] {$\ell + 1$}}}
    child {node[leaf] {$\ell$}}}
  child {node[inner] {$2$}
    child {node[leaf] {$\ell$}}
    child {node[leaf] {$\ell$}}};
\node[right=3mm of root, align=left, font=\footnotesize] {size $n = 5\ell + 2$};
\end{tikzpicture}

\caption{The tree of the proof of \cref{lem:conv-tree} for $k = 5$ blocks; internal nodes are labeled by their numbers $k_\nu$ of blocks, and leaves by their sizes. Each internal node with $k_\nu$ blocks below it has children with $\lceil k_\nu/2 \rceil$ and $\lfloor k_\nu/2 \rfloor$ blocks, and the leaves are blocks of size $\ell$ or $\ell + 1$. Along every path from a leaf up to the root, the sizes grow by a factor of at least $4/3$ at each step, so the errors in \eqref{eq:conv-superadd} accumulated along the path are bounded by the path-error sum $\mathcal R(\ell)$.}
\label{fig:conv-tree}
\end{figure}

We claim that the children $\nu_1$ and $\nu_2$ of every internal node $\nu$ satisfy
\begin{equation}\label{eq:conv-children}
\frac{N_\nu}{4} \le N_{\nu_i} \le \frac{3 N_\nu}{4}, \qquad i = 1, 2.
\end{equation}
Indeed, each child has at least $\lfloor k_\nu/2 \rfloor \ge 1$ and at most $\lceil k_\nu/2 \rceil \le 2 \lfloor k_\nu/2 \rfloor$ blocks, and each block has size between $\ell$ and $\ell + 1 \le 3\ell/2$, so that
\[
\frac{N_{\nu_1}}{N_{\nu_2}} \le \frac{\lceil k_\nu/2 \rceil (\ell + 1)}{\lfloor k_\nu/2 \rfloor \ell} \le 2 \cdot \frac 32 = 3,
\]
and the same bound holds for $N_{\nu_2}/N_{\nu_1}$. Together with $N_{\nu_1} + N_{\nu_2} = N_\nu$, this gives \eqref{eq:conv-children}.

Now let $\nu = \nu_0, \nu_1, \dots, \nu_{\operatorname{ht}(\nu)}$ be a path of maximum length from $\nu$ down to a leaf. Using \eqref{eq:conv-children}, we have $N_{\nu_{i - 1}} \ge (4/3) N_{\nu_i}$ for all $i$, which implies that
\begin{equation}\label{eq:conv-growth}
N_\nu \ge \Big(\frac 43\Big)^{\operatorname{ht}(\nu)} \ell.
\end{equation}

We now show by induction on the height that every node $\nu$ satisfies
\begin{equation}\label{eq:conv-induction}
\xi_{N_\nu} \ge \xi_\ell - g(\ell) - \sum_{i = 1}^{\operatorname{ht}(\nu)} g\Big(\frac \ell4 \Big(\frac 43\Big)^i\Big).
\end{equation}
A node of height $0$ is a leaf of size $\ell$ or $\ell + 1$. The second inequality in \eqref{eq:conv-superadd}, applied with $N_1 = \ell \ge n_0$, gives $\xi_{\ell + 1} \ge \xi_\ell - g(\ell)$, which proves \eqref{eq:conv-induction} for leaves.

Now let $\nu$ be an internal node with children $\nu_1$ and $\nu_2$, and assume that \eqref{eq:conv-induction} holds for $\nu_1$ and $\nu_2$. Their heights are at most $\operatorname{ht}(\nu) - 1$, and their sizes are at least $\ell \ge n_0$. Using the first inequality in \eqref{eq:conv-superadd} and the induction hypothesis, we obtain
\[
\begin{aligned}
\xi_{N_\nu} &\ge \min(\xi_{N_{\nu_1}}, \xi_{N_{\nu_2}}) - g\big(\min(N_{\nu_1}, N_{\nu_2})\big)\\
&\ge \xi_\ell - g(\ell) - \sum_{i = 1}^{\operatorname{ht}(\nu) - 1} g\Big(\frac \ell4 \Big(\frac 43\Big)^i\Big) - g\big(\min(N_{\nu_1}, N_{\nu_2})\big).
\end{aligned}
\]
The bounds \eqref{eq:conv-children} and \eqref{eq:conv-growth} give
\[
\min(N_{\nu_1}, N_{\nu_2}) \ge \frac{N_\nu}{4} \ge \frac \ell4 \Big(\frac 43\Big)^{\operatorname{ht}(\nu)} \ge u_0,
\]
and since $g$ is nonincreasing on $[u_0, \infty)$, the last term in the previous display is at least $-g((\ell/4)(4/3)^{\operatorname{ht}(\nu)})$. This completes the induction.

Applying \eqref{eq:conv-induction} at the root, we obtain $\xi_n \ge \xi_\ell - g(\ell) - \mathcal R(\ell)$ for all $n \ge \ell^2$. Taking the limit inferior as $n \to \infty$, we find
\[
\liminf_{n \to \infty} \xi_n \ge \xi_\ell - g(\ell) - \mathcal R(\ell)
\]
for every integer $\ell \ge \max(n_0, 4u_0, 2)$. We have shown that $g(\ell) \to 0$ and $\mathcal R(\ell) \to 0$, and taking the limit superior as $\ell \to \infty$ gives $\liminf_n \xi_n \ge \limsup_\ell \xi_\ell$. The sequence is bounded, so it converges.
\end{proof}

The minimum form of the size comparison and the tree with blocks of sizes $\ell$ and $\ell + 1$ are as in \cite[Lemma~5.4]{openai2026limiting}, which assumes polynomially small errors. For the classical additive form of near-superadditivity, sufficient conditions of this summable type go back to de Bruijn and Erd\H{o}s \cite[Theorem~23]{debruijn1952some}; see also \cite[arXiv:1006.3786v1, Section~4.3, Lemma~8 and Remark~3]{abbe2014concentration}. A minimum inequality is weaker than the additive one, and these classical results do not apply directly.

\begin{remark}\label{rem:conv-necessity}
The summability assumption in \cref{lem:conv-tree} and the adjacent-size inequality in \eqref{eq:conv-superadd} cannot be dropped from this convergence argument. First, let $\tilde\delta \colon [1, \infty) \to (0, \infty)$ be nonincreasing with $\int_1^\infty \tilde\delta(v) \dd v/v = \infty$ and, for $t \ge 1$, define $\tilde F(t) = \int_1^t \tilde\delta(v) \dd v/v$. For integers $n \ge 1$, we set $y_n = 1 + \sin(\tilde F(n))/2$. For $3 \le N_1 \le N_2$, we have
\[
|y_{N_1 + N_2} - y_{N_2}| \le \frac 12 \int_{N_2}^{N_1 + N_2} \tilde\delta(v) \frac{\dd v}{v} \le \frac{\log 2}{2} \tilde\delta(N_2) \le \frac{\log 2}{2} \tilde\delta(N_1),
\]
and $|y_{N_1 + 1} - y_{N_1}| \le \tilde\delta(N_1)/(2N_1) \le \tilde\delta(N_1)$. Thus $(y_n)$ satisfies the two inequalities in \eqref{eq:conv-superadd} with $g = \tilde\delta$ and $n_0 = 3$. However, $\tilde F(n) \to \infty$ and $\tilde F(n + 1) - \tilde F(n) \le \tilde\delta(1)/n \to 0$, so the values $\tilde F(n)$ come arbitrarily close to all large real numbers, and therefore $\limsup_n y_n = 3/2$ and $\liminf_n y_n = 1/2$. For nonincreasing $\tilde\delta$, the divergence of $\int_1^\infty \tilde\delta(v) \dd v/v$ is equivalent to $\sum_k \tilde\delta(e^k) = \infty$, and this includes $\tilde\delta(v) = 1/\log(e v)$. Thus, in this argument, a window of order $n/\log n$ would not suffice. The example concerns the tree argument only; it does not show that the quantiles of $V$ fail to converge under a weaker window bound. For the classical additive form, the necessity of the corresponding summability condition is \cite[arXiv:1810.11723v1, Theorem~6]{furedi2020nearly}.

Second, we define $\tilde b_n = \ind\{n\text{ is even}\}$ for integers $n \ge 1$. Then $\tilde b_{N_1 + N_2} \ge \min(\tilde b_{N_1}, \tilde b_{N_2})$ for all $N_1, N_2 \ge 1$, since $N_1 + N_2$ is even whenever $N_1$ and $N_2$ are even. Thus $(\tilde b_n)$ satisfies the first inequality in \eqref{eq:conv-superadd} for every $g \ge 0$, but does not converge. It violates the second inequality for every $g$ that tends to zero.
\end{remark}

\section{Transfer to the uniform and binomial random graphs}\label{sec:conv-transfer}

Recall that, for integers $n \ge 1$ and $0 \le m \le \binom{n}{2}$, we let $G(n, m)$ denote the uniformly random simple graph on $[n]$ with $m$ edges. For $p \in [0, 1]$, we let $G(n, p)$ denote the random simple graph on $[n]$ in which each pair is an edge independently with probability $p$.

\begin{lemma}\label{lem:conv-transfer}
\begin{enumerate}[label=(\alph*)]
\item For all large $n$ and all integers $0 \le m \le \Lambda n$, we have
\begin{equation}\label{eq:conv-transfer-a}
\begin{aligned}
\Pp\big(G(n, m) \text{ is not colorable}\big) &\le e^{2\Lambda^2 + 1} \Pp(H \le m), \\
\Pp\big(G(n, m) \text{ is colorable}\big) &\le e^{2\Lambda^2 + 1} \Pp(H > m).
\end{aligned}
\end{equation}
\item For all integers $n \ge 2$, the probability $\Pp(G(n, m) \text{ is colorable})$ is nonincreasing in $m \in \{0, \dots, \binom{n}{2}\}$. Further, for all integers $0 \le m \le \binom{n}{2}$ and all $p \in [0, 1]$, if $X$ has the $\Bin(\binom{n}{2}, p)$ law, then
\begin{equation}\label{eq:conv-transfer-b}
\begin{aligned}
\Pp\big(G(n, p) \text{ is colorable}\big) &\ge \Pp\big(G(n, m) \text{ is colorable}\big) \Pp(X \le m), \\
\Pp\big(G(n, p) \text{ is colorable}\big) &\le \Pp\big(G(n, m) \text{ is colorable}\big) + \Pp(X < m).
\end{aligned}
\end{equation}
\end{enumerate}
\end{lemma}

\begin{proof}
We begin with (a). Conditioning on a simplicity event of probability bounded away from zero is a standard transfer from nonsimple to simple random graphs; see \cite[Section~2, paragraph on nonsimple graphs]{bayati2013combinatorial}. Let $b = \binom{n}{2}$, and let $\mathcal D_m$ denote the event that the pairs $e_1, \dots, e_m$ are distinct. For all large $n$ and all $m \le \Lambda n$, we have $m/b \le 2\Lambda/(n - 1) \le 1/2$. Using $\log(1 - x) \ge -x - x^2$ for $x \in [0, 1/2]$, we obtain
\[
\begin{aligned}
\Pp(\mathcal D_m) &= \prod_{i = 1}^{m - 1} \Big(1 - \frac{i}{b}\Big) \\
&\ge \exp\Big(-\sum_{i = 1}^{m - 1} \frac{i}{b} - \sum_{i = 1}^{m - 1} \frac{i^2}{b^2}\Big) \\
&\ge \exp\Big(-\frac{m^2}{2b} - \frac{m^3}{b^2}\Big) \\
&\ge e^{-2\Lambda^2 - 1},
\end{aligned}
\]
where in the last step we used $m^2/(2b) = m^2/(n(n - 1)) \le 2\Lambda^2$ and $m^3/b^2 \le 4\Lambda^3 n/(n - 1)^2 \le 1$ for all large $n$.

Conditionally on $\mathcal D_m$, the sequence $(e_1, \dots, e_m)$ is uniformly distributed on the sequences of $m$ distinct pairs of $[n]$. Therefore the set $\{e_1, \dots, e_m\}$ has the law of the edge set of $G(n, m)$. The colorability of $G_m$ is determined by the set of its distinct pairs. If $\mathcal B$ denotes the event that $G_m$ is not colorable, we have
\[
\Pp\big(G(n, m) \text{ is not colorable}\big) = \Pp(\mathcal B \mid \mathcal D_m) \le \frac{\Pp(\mathcal B)}{\Pp(\mathcal D_m)} \le e^{2\Lambda^2 + 1} \Pp(H \le m),
\]
and the second bound in \eqref{eq:conv-transfer-a} follows in the same way from the complement of $\mathcal B$.

We now prove (b). Let $m < \binom{n}{2}$. Adding to $G(n, m)$ a uniformly random pair of $[n]$ that is not an edge of $G(n, m)$ gives a graph with the law of $G(n, m + 1)$. A graph that contains a non-colorable graph is not colorable, which gives the monotonicity.

For this proof, we let $X$ denote the number of edges of $G(n, p)$. Then $X$ has the $\Bin(\binom n2, p)$ law. Conditionally on $X$, the graph $G(n, p)$ has the law of $G(n, X)$ \cite[Lemma~1.1]{frieze2016introduction}, so
\[
\Pp\big(G(n, p) \text{ is colorable}\big) = \sum_{j = 0}^{\binom{n}{2}} \Pp(X = j) \Pp\big(G(n, j) \text{ is colorable}\big).
\]
Bounding the terms with $j \ge m$ by $\Pp(X = j) \Pp(G(n, m) \text{ is colorable})$ and the terms with $j < m$ by $\Pp(X = j)$ gives the second inequality in \eqref{eq:conv-transfer-b}. Dropping the terms with $j > m$ and bounding the others from below by $\Pp(X = j) \Pp(G(n, m) \text{ is colorable})$ gives the first.
\end{proof}

\section{Proof of \texorpdfstring{\cref{thm:red-convergence}}{the convergence theorem}}\label{sec:conv-proof}

We assume the hypotheses recalled at the beginning of this appendix, and we proceed in four steps.

\emph{Step 1 (one level).} Fix $\eta \in (0, \eta_0]$, and let $\xi_n$ and $g$ be as in \cref{sec:conv-superadd}. The function $g$ is positive, since $K' > 0$ and $\delta_n > 0$. It is nonincreasing on $[e^2, \infty)$, because $\delta(u)$ is nonincreasing in $u$ and $(\log u)/\sqrt u$ is nonincreasing for $u \ge e^2$. Further, by \eqref{eq:conv-summable}, we have
\[
\sum_{k \ge 1} g(e^k) = K' \sum_{k \ge 1} \delta_{\lfloor e^k \rfloor} + (\Lambda + 6) \sum_{k \ge 1} k e^{-k/2} < \infty.
\]
The sequence $(\xi_n)_{n \ge 2q}$ takes values in $[0, \Lambda]$. By \cref{lem:conv-superadd}, it satisfies the two inequalities in \eqref{eq:conv-superadd} for all $N_1, N_2 \ge n_0$. Therefore \cref{lem:conv-tree} with $u_0 = e^2$ shows that $(\xi_n)$ converges. We let $\alpha(\eta)$ denote its limit.

\Cref{lem:pre-centre} implies that $c_0 \le \alpha(\eta) \le c_1$ for all $c_0 < 1/2$ and all $c_1 > \gamma_q$, and we conclude that $\alpha(\eta) \in [1/2, \gamma_q]$. The summability \eqref{eq:conv-summable} implies $\delta_{\lfloor e^k \rfloor} \to 0$, and since $(\delta_n)$ is nonincreasing, we have $\delta_n \to 0$. Using \eqref{eq:conv-window}, we conclude that $U_n(\eta)/n = L_n(\eta)/n + w_n(\eta)/n$ also converges to $\alpha(\eta)$.

\emph{Step 2 (all levels share the limit).} Let $0 < \eta_- < \eta_+ \le \eta_0$. Every $m$ with $\Pp(V \le m) \ge \eta_+$ also satisfies $\Pp(V \le m) \ge \eta_-$, which implies that $L_n(\eta_-) \le L_n(\eta_+)$. Further, as in the proof of \cref{lem:pre-quantile}, we have
\[
\Pp\big(V \ge L_n(\eta_+)\big) = 1 - \Pp\big(V \le L_n(\eta_+) - 1\big) > 1 - \eta_+ \ge \eta_-,
\]
and $L_n(\eta_+) \le M$, so $L_n(\eta_+)$ belongs to the set that defines $U_n(\eta_-)$. Hence
\[
\frac{L_n(\eta_-)}{n} \le \frac{L_n(\eta_+)}{n} \le \frac{U_n(\eta_-)}{n} \qquad \text{for all } n \ge 2q.
\]
By Step 1, the outer terms converge to $\alpha(\eta_-)$. It follows that $\alpha(\eta_+) = \alpha(\eta_-)$. We denote the common value by $\alpha_q$. Together with Step 1, this shows that $L_n(\eta)/n \to \alpha_q$ and $U_n(\eta)/n \to \alpha_q$ for every $\eta \in (0, \eta_0]$, and that $\alpha_q \in [1/2, \gamma_q]$.

\emph{Step 3 (the edge process).} Let $0 \le c < \alpha_q$, and let $\eta \in (0, \eta_0]$. Since $L_n(\eta)/n \to \alpha_q > c$, we have $\lfloor cn \rfloor \le L_n(\eta)$ for all large $n$, and for such $n$ \cref{lem:pre-quantile} gives $\Pp(H > \lfloor cn \rfloor) > 1 - \eta$. Therefore $\liminf_n \Pp(H > \lfloor cn \rfloor) \ge 1 - \eta$, and letting $\eta \to 0$, we conclude that $\Pp(H > \lfloor cn \rfloor) \to 1$.

Now let $c > \alpha_q$, and assume first that $c < \Lambda$. The convergence $U_n(\eta)/n \to \alpha_q < c$ implies that $U_n(\eta) < \lfloor cn \rfloor \le M$ for all large $n$, and for such $n$ \cref{lem:pre-quantile} gives $\Pp(H > \lfloor cn \rfloor) < \eta$. Letting first $n \to \infty$ and then $\eta \to 0$, we conclude that $\Pp(H > \lfloor cn \rfloor) \to 0$. For $c \ge \Lambda$, we choose $c' \in (\alpha_q, \Lambda)$ and recall that $\Pp(H > m)$ is nonincreasing in $m$, so that $\Pp(H > \lfloor cn \rfloor) \le \Pp(H > \lfloor c' n \rfloor) \to 0$.

\emph{Step 4 (the uniform and binomial random graphs).} Let $0 \le c < \alpha_q$. Since $\alpha_q \le \gamma_q < \Lambda$, we have $\lfloor cn \rfloor \le \Lambda n$, and \cref{lem:conv-transfer}(a) and Step 3 give
\[
\Pp\big(G(n, \lfloor cn \rfloor) \text{ is not colorable}\big) \le e^{2\Lambda^2 + 1} \Pp\big(H \le \lfloor cn \rfloor\big) \to 0.
\]
For $\alpha_q < c < \Lambda$, we use the second bound in \eqref{eq:conv-transfer-a} and Step 3 to obtain $\Pp(G(n, \lfloor cn \rfloor) \text{ is colorable}) \to 0$. For $c \ge \Lambda$, we choose $c' \in (\alpha_q, \Lambda)$. For all large $n$, we have $\lfloor cn \rfloor \le \binom{n}{2}$, and the monotonicity in \cref{lem:conv-transfer}(b) reduces this case to the previous one.

Now let $d \ge 0$ and let $n$ be so large that $d/n \le 1$. Let $X$ have the $\Bin(\binom{n}{2}, d/n)$ law. Then $\E X = d(n - 1)/2$ and $\Var X \le d(n - 1)/2$.

Suppose first that $d < 2\alpha_q$. We choose $c \in (d/2, \alpha_q)$ and set $m = \lfloor cn \rfloor$. Then $m + 1 - \E X \ge (c - d/2) n > 0$, and Chebyshev's inequality gives
\[
\Pp(X > m) \le \Pp\big(X - \E X \ge (c - d/2) n\big) \le \frac{d n}{2 (c - d/2)^2 n^2} \to 0.
\]
The first inequality in \eqref{eq:conv-transfer-b} and the first part of this step then give
\[
\Pp\big(G(n, d/n) \text{ is colorable}\big) \ge \Pp\big(G(n, m) \text{ is colorable}\big) \Pp(X \le m) \to 1.
\]
Suppose now that $d > 2\alpha_q$. We choose $c \in (\alpha_q, d/2)$ and set $m = \lfloor cn \rfloor$. Then $\E X - m \ge (d/2 - c) n - d/2$, which grows linearly in $n$, and Chebyshev's inequality gives $\Pp(X < m) \to 0$. The second inequality in \eqref{eq:conv-transfer-b} then gives
\[
\Pp\big(G(n, d/n) \text{ is colorable}\big) \le \Pp\big(G(n, m) \text{ is colorable}\big) + \Pp(X < m) \to 0.
\]
This completes the proof of \cref{thm:red-convergence}.

We note that $\alpha_q$ is determined by $q$ alone. Indeed, by Step 3, it is the unique number such that $\Pp(H > \lfloor cn \rfloor)$ tends to $1$ for $c < \alpha_q$ and to $0$ for $c > \alpha_q$, and these probabilities do not involve $\Lambda$, $\eta$, or the sequence $(\delta_n)$. In particular, if the hypotheses of \cref{thm:red-convergence} hold for two values of $\Lambda$, they give the same $\alpha_q$.



\backmatter

\setlength{\bibsep}{3pt plus 1pt minus 1pt}
\bibliographystyle{plainnat}
\bibliography{refs}
\end{document}
